\documentclass[10pt]{amsart}
\usepackage[a4paper,margin=23mm]{geometry}
\usepackage{amsmath,amssymb,amsthm,mathtools,hyperref,graphicx,xspace,letltxmacro,xpatch}
\usepackage[english]{babel}
\newtheorem{theo}{Theorem}[section]
\newtheorem{lemm}[theo]{Lemma}\newtheorem{prop}[theo]{Proposition}\newtheorem{coro}[theo]{Corollary}\newtheorem{conj}[theo]{Conjecture}
\theoremstyle{definition}\newtheorem{defi}[theo]{Definition}\newtheorem{rema}[theo]{Remark}\newtheorem{exam}[theo]{Example}\newtheorem{exem}[theo]{Example}\newtheorem{nota}[theo]{Notation}
\newcommand{\ie}{i.e.}\newcommand{\eg}{e.g.}\newcommand{\resp}{resp.}\newcommand{\skpt}{\par\smallskip\noindent}
\emergencystretch=3em\pagestyle{plain}
%\usepackage{enumerate}
%\usepackage{enumitem}
\usepackage{graphicx}
\usepackage{float}
\usepackage{verbatim}
\usepackage{amssymb}
%\usepackage{color}
%\usepackage{amsthm}
%\usepackage{microtype}
\usepackage{comment}
\usepackage{mathrsfs}







\newcommand{\scC}{\mathscr{C}}
\newcommand{\scL}{\mathscr{L}}
\newcommand{\scS}{\mathscr{S}}
\newcommand{\scT}{\mathscr{T}}
\newcommand{\scD}{\mathscr{D}}
\newcommand*\R{\mathbb{R}}
\newcommand*\N{\mathbb{N}}
\newcommand*\C{\mathbb{C}}
\renewcommand*\S{\mathbb{S}}
\newcommand{\bfR}{\mathbf{R}}
\renewcommand{\i}{\mathbf{i}}

\newcommand{\clV}{\mathcal{V}}
\newcommand{\clO}{\mathcal{O}}
\newcommand{\clM}{\mathcal{M}}
\newcommand{\clN}{\mathcal{N}}
\newcommand{\clT}{\mathcal{T}}
\newcommand{\clF}{\mathcal{F}}
%\newcommand*\Z{{\mathbb Z}}
%\newcommand*\T{{\mathbb{T}}}
\newcommand*\ds{\displaystyle}
\newcommand{\norm}[1]{\left\Vert#1\right\Vert}
\renewcommand{\leq}{\leqslant}
\renewcommand{\geq}{\geqslant}
\renewcommand{\tilde}{\widetilde}
\renewcommand{\hat}{\widehat}
%\renewcommand{\bar}{\overline}
%\renewcommand{\div}{\hbox{div\,}}
%\newcommand{\red}[1]{{\color{red} #1}}
%\newcommand{\RED}[1]{\textsc{\color{red}#1}\\}
%\newcommand{\blue}[1]{{\color{blue} #1}}
%\newcommand{\BLUE}[1]{\textsc{\color{blue}#1}\\}

\DeclareMathOperator{\Hess}{Hess}
\DeclareMathOperator{\supp}{supp}
\renewcommand{\det}{\operatorname{det}}
\renewcommand{\div}{\operatorname{div}}
\newcommand{\divv}{\mathrm{div}}
\DeclareMathOperator{\Span}{Span}
%\AtBeginDocument{
%\newtheorem{myname}[cdrthm]{My beautiful theorem} }


%%%%%%%%%%%%%%%%%%%%%%%%%%%%%%%%%%%%%%%%%%%%%%%%%%%%%%%%%%%%%%%%%%%%%%%%%%%%%%%%%%%%%%%%%%%%%%%%%%%%%%%%%%%%%%%%%%%%%%%%%%%%

\graphicspath{{./figures/}}

\newcommand*{\mk}{\mkern -1mu}
\newcommand*{\Mk}{\mkern -2mu}
\newcommand*{\mK}{\mkern 1mu}
\newcommand*{\MK}{\mkern 2mu}

\hypersetup{urlcolor=purple, linkcolor=blue, citecolor=red}

\newcommand*{\romanenumi}{\renewcommand*{\theenumi}{\roman{enumi}}}
\newcommand*{\Romanenumi}{\renewcommand*{\theenumi}{\Roman{enumi}}}
\newcommand*{\alphenumi}{\renewcommand*{\theenumi}{\alph{enumi}}}
\newcommand*{\Alphenumi}{\renewcommand*{\theenumi}{\Alph{enumi}}}
\let\oldtilde\tilde
\renewcommand*{\tilde}[1]{\mathchoice{\widetilde{#1}}{\widetilde{#1}}{\oldtilde{#1}}{\oldtilde{#1}}}
\let\oldhat\hat
\renewcommand*{\hat}[1]{\mathchoice{\widehat{#1}}{\widehat{#1}}{\oldhat{#1}}{\oldhat{#1}}}
\let\oldexists\exists
\renewcommand*{\exists}{\mathrel{\oldexists}}
\let\oldforall\forall
\renewcommand*{\forall}{\mathrel{\oldforall}}
%%%%%%%%%%%%%%%%%%%%%%%%%%%%%%%%%%%%%%%%%%%%%%%%%%%%%%%%%%%%%%%%%%%%%%%%%%%%%%%%%%%%%%%%%%%%%%%%%%%%%%%%%%%%%%%%%%%%%%%%%%%%

\title{Global Lp Carleman estimates for elliptic boundary problems with general subelliptic weights}
\author{Dehman, Belhassen, Ervedoza, Sylvain, Thabouti, Lotfi}
\date{doi:10.5802/ahl.226; CC BY 4.0}
\begin{document}
\maketitle
\noindent\textit{Editorial note.} Only selected theorem counted; full mathematical body and both appendices retained, without external figure dependencies. Original source macros and original publisher reference keys preserved.\par\medskip
\noindent\textit{Allowed standard prerequisites (editorial).}\begin{itemize}
\item Standard functional analysis, measure theory and the author-cited classical theorems. Standard prior results stay cited; complete source-local proofs and definitions retained.
\end{itemize}\medskip

\section{Introduction}

\subsection*{Main result}
The goal of this article is to prove global $L^p$ Carleman estimates for the flat Laplace operator in a smooth bounded domain of $\R^d$ ($d \geq 3$) for a general weight function satisfying the sub-ellipticity conditions of Hörmander. As an application, we will show how these can be used to obtain quantitative unique continuation results for solutions of elliptic equation with respect to the norms of the potentials.

To be more precise, our main result is the following one:

\begin{theo}\label{Thm-Carleman-General}
Let $d\geqslant 3$. Let $\Omega \subset \R^{d}$ be a bounded domain of class $C^3$, and $\omega$ be a non-empty open subset of $\Omega$ with $\overline \omega \subset \Omega$. Let $\varphi \in C^3(\overline{\Omega})$ be such that
\begin{equation}\label{CarlemanWeight-Cond0}
\forall x \in \partial \Omega, \, \varphi(x) = 0 \text{ and }
\partial_n \varphi(x) < 0,
\end{equation}
and there exists $\alpha, \beta > 0$ for which
\begin{equation}\label{CarlemanWeight-Cond1}
\inf_{\overline{\Omega}\setminus \omega} \vert \nabla \varphi \vert > \alpha,
\end{equation}
and
\begin{multline}\label{CarlemanWeight-Cond2}
\forall x \in \overline{\Omega} \setminus \omega, \, \forall \xi \in \R^d \quad\text{with}\quad \vert \nabla \varphi(x) \vert = \vert \xi \vert \quad\text{and}\quad \nabla \varphi(x) \cdot \xi = 0,
\\
(\Hess \varphi(x)) \nabla \varphi(x) \cdot \nabla \varphi(x) + (\Hess \varphi(x)) \xi \cdot \xi \geq \beta \vert \nabla \varphi(x) \vert^2,
\end{multline}
where $\Hess \varphi$ denotes the Hessian matrix of $\varphi$. Let $\omega_1$ be an open subset of $\Omega$ so that $\overline \omega \subset \omega_1$ and $\overline\omega_1 \subset \Omega$, and $\eta$ be a smooth radial non-negative cut-off function (in $\scC^\infty_c(\R^d)$) vanishing outside the ball of radius $1$ and equal to one in the ball of radius~$\frac{1}{2}$.

Then there exist $C>0$ and $\tau_0 \geq 1$ (depending only on $\alpha$, $\beta$, $\| \varphi \|_{C^3(\overline\Omega)}$, $\eta$ and the geometric configuration $\Omega$, $\omega$, and $\omega_1$) such that for all $u \in H^1(\Omega)$ solution of
\begin{equation} \label{Elliptic-Eq-general}
\begin{cases}
- \Delta u = f_2 + f_{2_*'} + \div F \quad &\text{in } \Omega,\\
u=g \quad &\text{on } \partial \Omega,\\
\end{cases}
\end{equation}
with
\begin{equation*}\label{Reg-Source-Terms}
f_2 \in L^2(\Omega),
\quad f_{2_*'} \in L^{ \frac{2d}{d+2}}(\Omega),
\quad F \in L^2(\Omega; \C^{d}),
\quad\text{ and } \quad g \in H^{\frac{1}{2}}(\partial\Omega),
\end{equation*}
we have, for all $\tau \geq \tau_0$,
\begin{multline}\label{Carleman-General-1}
\tau^{\frac{3}{2}} \left\| e^{\tau \varphi} u \right\|_{L^2(\Omega)} +
\tau^{\frac{1}{2}} \left\| e^{\tau \varphi} \nabla u \right\|_{L^2(\Omega)}
\leq C \left(
\left\| e^{\tau \varphi} f_2 \right\|_{L^2(\Omega)} +
\tau^{\frac{3}{4}-\frac{1}{2d}} \left\| e^{\tau \varphi} f_{2_*'} \right\|_{L^{\frac{2d}{d+2}}_{\eta,\tau} (\Omega)}
\right.
\\
\left. +
\tau \left\| e^{\tau \varphi} F \right\|_{L^2(\Omega)} +
\tau^{\frac{3}{4}} \|g \|_{H^{\frac{1}{2}}(\partial \Omega)} +
\tau^{\frac{3}{2}} \left\| e^{\tau \varphi} u\right\|_{L^2(\omega_1)} +
\tau^{\frac{3}{4}} \left\| e^{\tau \varphi} u \right\|_{L^{\frac{2d}{d-2}}_{\eta,\tau} (\omega_1)}
\right),
\end{multline}
and
\begin{multline}\label{Carleman-General-2}
\tau^{\frac{3}{4}+\frac{1}{2d}} \left\| e^{\tau \varphi} u \right\|_{L^{\frac{2d}{d-2}}_{\eta,\tau} (\Omega)}
\leq C \left(
\left\| e^{\tau \varphi} f_2 \right\|_{L^2(\Omega)} +
\tau^{\frac{3}{4}+\frac{1}{2d}} \left\| e^{\tau \varphi} f_{2_*'} \right\|_{L^{\frac{2d}{d+2}}_{\eta,\tau} (\Omega)}
\right.
\\
\left. +
\tau \left\| e^{\tau \varphi} F \right\|_{L^2(\Omega)} +
\tau^{\frac{3}{4}+\frac{1}{2d}} \|g \|_{H^{\frac{1}{2}}(\partial \Omega)} +
\tau^{\frac{3}{2}} \left\| e^{\tau \varphi} u\right\|_{L^2(\omega_1)} +
\tau^{\frac{3}{4}+\frac{1}{2d}} \left\| e^{\tau \varphi} u \right\|_{L^{\frac{2d}{d-2}}_{\eta,\tau} (\omega_1)}
\right).
\end{multline}
Here, the norms $\| \cdot \|_{L^p_{\eta, \tau}(\Omega)}$ are defined for $p \in [1, \infty]$ for $f \in L^p(\Omega)$ by the formula
\begin{equation}\label{Def-L-p-eta-Omega}
\left\| f\right\|_{L^p_{\eta, \tau}(\Omega)}^2 =
\tau^{\frac{d}{3}} \int_{x_0\,\in\,\Omega} \left\| \eta\left(\tau^{\frac{1}{3}} (\cdot -x_0)\right) f(\cdot) \right\|_{L^p(\Omega)}^2 \, dx_0.
\end{equation}
\end{theo}

\begin{rema}
The notations $2_*$ and $2_*'$ stem from the Sobolev's embedding $H^1(\Omega) \subset L^{2_*}(\Omega)$, with $2_* = \frac{2d}{d-2}$ and $L^{2_*'}(\Omega) \subset H^{-1}(\Omega)$, with $2_*' = \frac{2d}{d+2}$.
\end{rema}


Before going further, let us remark that the existence of a function satisfying the conditions~\eqref{CarlemanWeight-Cond0}--\eqref{CarlemanWeight-Cond1}--\eqref{CarlemanWeight-Cond2} for any arbitrary geometric setting is due to Fursikov and Imanuvilov~\cite[Lemma~1.1]{Fursikov1996-book} (see also~\cite[Proposition~3.31]{LeRousseaubookvolume1}). The conditions~\eqref{CarlemanWeight-Cond1}--\eqref{CarlemanWeight-Cond2} are the \emph{sub-ellipticity conditions} of the weight function $\varphi$ with respect to the Laplace operator, which are known to be necessary and sufficient conditions to get a local $L^2$ Carleman estimate (i.e.~\eqref{Carleman-General-1} for compactly supported functions $u$, and with $f_{2_*'} = 0$) for the Laplace operator with the same powers of the Carleman parameter $\tau$, see~\cite[Chapter~XXVIII]{hormanderanalysis4} and, for instance~\cite[Definition~3.2, Section~3.6 and Section~4.1.2]{LeRousseaubookvolume1} for a more recent perspective.


The Carleman estimate~\eqref{Carleman-General-1} coincides with the one in~\cite{imanuvilov2003global} except for the terms involving the norm
\[
L^{\frac{2d}{d+2}}_{\eta,\tau}(\Omega).
\]
This term and the estimate~\eqref{Carleman-General-2} on
\[
u \text{ in }L^{\frac{2d}{d-2}}_{\eta,\tau}(\Omega)
\]
are the main novelties of our result and allow us to quantify efficiently unique continuation properties for solutions of elliptic equations with respect to the norms of potentials in $L^p(\Omega)$.

Note that estimate~\eqref{Carleman-General-1} implies an estimate on $\tau^{\frac{1}{2}} \| e^{\tau \varphi} u\|_{H^1(\Omega)}$ from the right hand side of~\eqref{Carleman-General-1}, thus on $\tau^{\frac{1}{2}} \|e^{\tau \varphi} u\|_{L^{\frac{2d}{d-2}}(\Omega)}$. Therefore, estimate~\eqref{Carleman-General-1} does not allow to recover estimate~\eqref{Carleman-General-2} directly from classical Sobolev's embeddings.

Local $L^p$ Carleman estimates (i.e Carleman estimates for compactly supported functions) have been derived in many situations, but usually to focus on questions related to unique continuation. We should in particular quote the breakthrough article~\cite{jerison1985unique} obtained for a radial weight $\log(|x|)$, which rather corresponds to a limiting Carleman weight in the sense that the second condition~\eqref{CarlemanWeight-Cond2} is satisfied with $\beta = 0$ (see also the previous results~\cite{Amrein-Berthier-Georgescu-1981, Hormander-Uniqueness-1983}). Later, several works have been devoted to get local Carleman estimates with some specific strictly convex weights, see e.g.~\cite{barcelo1988weighted,Sogge1989uniquecontinuation,Sogge.Strong.Elliptic1990}, which have later been revisited and improved in the works~\cite{koch2001carleman,ferreira2003sharp}. We also point out the more recent works~\cite{davey2020quantitative,davey2019quantitative} for local Carleman estimates with some specific strictly convex weight. Here, we emphasize that we will consider general weight functions satisfying the sub-ellipticity conditions~\eqref{CarlemanWeight-Cond1}--\eqref{CarlemanWeight-Cond2}, similarly as in~\cite{ferreira2003sharp}, which considers the more general case of second order real principal type operators of order 2. In fact, in our context, the article~\cite{ferreira2003sharp} proves that, for all $x_0 \in \R^d$, if $\varphi$ is subelliptic at $x_0$ (that is $|\nabla \varphi(x_0)| \neq 0$ and condition~\eqref{CarlemanWeight-Cond2} at $x = x_0$), there exists a neighborhood $K$ of $x_0$ such that the local Carleman estimate $\| e^{\tau \varphi} u \|_{L^{\frac{2d}{d-2}}} \leq C \| e^{\tau \varphi} \Delta u \|_{L^{\frac{2d}{d+2}}}$ holds for all $u$ compactly supported in $K$. The estimates~\eqref{Carleman-General-1}--\eqref{Carleman-General-2} thus extend the result in~\cite{ferreira2003sharp} by providing a global Carleman estimate, allowing source terms in $H^{-1}(\Omega)$ and boundary conditions in $H^{\frac{1}{2}}(\partial \Omega)$, and estimating $u$ in the $H^1(\Omega)$-norm as well.

Finally, let us also emphasize that Theorem~\ref{Thm-Carleman-General} presents global $L^p$ Carleman estimates, in the sense that the Carleman estimates~\eqref{Carleman-General-1}--\eqref{Carleman-General-2} hold for functions $u$ having possibly non-zero trace on the boundary. To our knowledge, this is new, as all the $L^p$ Carleman estimates with $p \neq 2$ that we have encountered in the literature hold for compactly supported functions.

%\medskip

\subsection*{Properties of the $L^p_{\eta, \tau}(\Omega)$ norms}


To understand the norms $L^{p}_{\eta,\tau}(\Omega)$, let us first remark that for $p = 2$, and $f \in L^2_{\eta, \tau} (\Omega)$, by Fubini's theorem, we have, for $\tau$ large enough,
\begin{align*}
\| f \|_{L^2_{\eta, \tau}(\Omega)}^2 &\leq \|\eta\|_{L^2(\R^d)}^2 \| f\|_{L^2(\Omega)}^2,
\\
 \| f \|_{L^2_{\eta, \tau}(\Omega)}^2
&\geq \| f \|_{L^2(\Omega)}^2
\left(
\inf_{x_0\,\in\,\Omega} \tau^{\frac{d}{3}} \int_{x\,\in \,\Omega} \left|\eta\left(\tau^{\frac{1}{3}} (x -x_0)\right)\right|^2 \, dx
\right)
\geq c_* \| f\|_{L^2(\Omega)}^2.
\end{align*}
Here, $c_*>0$ is independent of $\tau \geq \tau_0$ if $\tau_0$ is chosen so that $\tau_0 \geq \varepsilon_0^{-3}$, where $\varepsilon_0 >0$ is such that for all $\varepsilon \in (0, \varepsilon_0]$ and $x \in \Omega$, there exists a ball of radius $\varepsilon/8$ contained in $B(x, \varepsilon/2) \cap \Omega$ (it is not difficult to check that such an $\varepsilon>0$ exists by compactness and smoothness of the boundary $\partial \Omega$). The norms $L^2_{\eta, \tau}(\Omega)$ are thus equivalent to the usual $L^2(\Omega)$ norm uniformly with respect to the parameter $\tau$.


For other values of $p \in [1, \infty)$, these norms are less easy to describe, as they somehow encode some mean information on the $L^p$-norms localized in balls of radius $\tau^{-\frac{1}{3}}$, as one can see by writing them under the form
\pagebreak
\[
\| f \|_{L^{p}_{\eta,\tau}(\Omega)} = \tau^{\frac{d}{6}} \norm{\norm{ \eta\left(\tau^{\frac{1}{3}} (x - x_0)\right) f(x) }_{L^p_x(\Omega)}}_{L^2_{x_0} (\Omega)}.
\]

In fact, for $p \in (1, \infty) $, again by Fubini's theorem, there exists $C>0$, such that for $\tau\geq \tau_0$ and $f \in L^p(\Omega)$,
\begin{multline}
\frac{\tau^{\frac{d}{3}}}{C} \int_{x\,\in\,\Omega} \int_{x_0\,\in\,\Omega} \left|\eta\left(\tau^{\frac{1}{3}}(x-x_0)\right) f(x)\right|^p \, dx \, dx_0
\leq\|f \|_{L^p(\Omega)}^p \\ \leq C \tau^{\frac{d}{3}} \int_{x\,\in\,\,\Omega} \int_{x_0\,\in\,\Omega} \left|\eta\left(\tau^{\frac{1}{3}}(x-x_0)\right) f(x)\right|^p \, dx \, dx_0,\label{Eq-Lp-Lpetatau}
\end{multline}
i.e. the $L^p$ norm $\| f \|_{L^p(\Omega)}$ is equivalent to the norm
\[
\tau^{\frac{d}{3p}} \norm{\norm{ \eta\left(\tau^{\frac{1}{3}} (x-x_0)\right) f(x) }_{L^p_x(\Omega)}}_{L^p_{x_0} (\Omega)}.
\]
This implies in particular that, for $p>2$,
\begin{multline}\label{Lp-larger-Lpetatau}
\| f \| _{L^{p}_{\eta, \tau}(\Omega)} = \tau^{\frac{d}{6}} \norm{\norm{ \eta\left(\tau^{\frac{1}{3}} (x-x_0)\right) f(x) }_{L^p_x(\Omega)}}_{L^2_{x_0} (\Omega)}
\\
\leq C \tau^{\frac{d}{6}} \norm{\norm{ \eta\left(\tau^{\frac{1}{3}} (x-x_0)\right) f(x) }_{L^p_x(\Omega)}}_{L^p_{x_0} (\Omega)}
\leq C \tau^{\left(\frac{1}{2} - \frac{1}{p}\right) \frac{d}{3}} \| f\|_{L^p(\Omega)}.
\end{multline}
On the other hand, for $p>2$, by Minkowski’s integral inequality (\cite{stein2016singular}, page 271), we have, for $C$ independent of $\tau \geq \tau_0$,
\begin{multline}\label{Lp-smaller-Lpetatau}
\norm{f}_{L^{p} (\Omega) } = \left(\int_{x\,\in \,\Omega} |f(x)|^p \, dx\right)^{\frac{1}{p}}
\\
\begin{aligned}
&\leq C \left(\int_{x\,\in\,\Omega} \left(\tau^{\frac{d}{3}} \int_{x_0\,\in\,\Omega}  \left|\eta\left(\tau^{\frac{1}{3}} (x-x_0)\right) f(x) \right|^{2} \, dx_0 \right)^{\frac{p}{2}} \, dx \right)^{\frac{1}{p}}
\\
&\leq C
\tau^{\frac{d}{6}} \left(\int_{x_0 \,\in\, \Omega} \left(\int_{x\,\in\,\Omega} \left| \eta\left(\tau^{\frac{1}{3}} (x-x_0)\right) f(x) \right|^{p} \, dx \right)^{\frac{2}{p}} \, dx_0 \right)^{\frac{1}{2}} = C \| f \| _{L^{p}_{\eta, \tau}(\Omega)}.
\end{aligned}
\end{multline}
Similarly, for $p<2$, we get, using Minkowski's integral inequality and the norm equivalence~\eqref{Eq-Lp-Lpetatau}, that there exists a constant $C$ independent of $\tau \geq \tau_0$, such that
\begin{equation}\label{Lp-eq-Lpetatau}
\frac{1}{C} \tau^{\left(\frac{1}{2} - \frac{1}{p}\right)\frac{d}{3}} \norm{ f }_{L^{p}(\Omega)}\leq \| f \| _{L^{p}_{\eta, \tau}(\Omega)} \leq C \norm{ f }_{L^{p}(\Omega)}.
\end{equation}

Note that, of course, the estimates~\eqref{Lp-larger-Lpetatau}, \eqref{Lp-smaller-Lpetatau}, and~\eqref{Lp-eq-Lpetatau} can be used to simplify the norms $L^p_{\eta, \tau}$ in the Carleman estimates~\eqref{Carleman-General-1} and~\eqref{Carleman-General-2} and replace them by the classical $L^p$ norms.

Finally, let us point out that, for $p,\, q, \, r$ in $[1, \infty]$ such that $\frac{1}{p} + \frac{1}{q} = \frac{1}{r}$, for $V \in L^q(\Omega)$ and $u \in L^p(\Omega)$, we have the following Hölder type estimate
\begin{equation}\label{Holder-in-Lp-tau-eta}
\| Vu \|_{L^r_{\eta, \tau}(\Omega)}
\leq \| u\|_{L^p_{\eta, \tau}(\Omega)} \sup_{x_0\,\in\,\Omega} \left\{ \| V\|_{L^q\left(B\left(x_0, \tau^{-\frac{1}{3}}\right)\right)} \right\}
\leq \| u\|_{L^p_{\eta, \tau}(\Omega)} \| V \|_{L^q(\Omega)}.
\end{equation}

%\medskip

\subsection*{Application to the quantification of unique continuation with respect to lower order terms}

 Next, as a consequence of the Carleman estimates in Theorem~\ref{Thm-Carleman-General}, we will prove (in Section~\ref{Sec-Proof-Thm-QUCP}) the following result:


\begin{theo}\label{Thm-QUCP}
Let $d\geqslant 3$, $\Omega \subset \R^{d}$ be a bounded domain of class $C^3$, and $\omega$ be a non-empty open subset of $\Omega$ with $\overline \omega \subset \Omega$.
Then there exists a constant $C=C(\Omega,\omega)>0$ depending only on $\Omega$ and $\omega$ such that for any solution $u\in H^{1}_0(\Omega)$ of
\begin{equation*}
\Delta u = V u + W_1 \cdot \nabla u + \div (W_2 u) \quad \text{ in } \Omega,
\end{equation*}
with
\begin{equation*}
V\in L^{q_0}(\Omega),
\quad W_1 \in L^{q_1}\left(\Omega; \C^d\right),
\quad W_2 \in L^{q_2}\left(\Omega; \C^d\right),
\end{equation*}
we have:
\begin{enumerate}
\item\label{theo1.3.1} If $q_0 \in (\frac{d}{2}, \infty]$, $q_1 \in (\frac{3d-2}{2}, \infty]$ and $q_2\in (\frac{3d-2}{2}, \infty]$ and
\[
\frac{1}{q_1} + \frac{1}{q_2} < 4\left(\frac{1-\frac{1}{d}}{3d-2}\right),
\]
the function $u$ satisfies
\begin{multline}\label{Quantitative-0}
\norm{u }_{L^2(\Omega)} \\
\leqslant C e^{ C \left(\norm{V}_{L^{q_0}(\Omega)}^{\gamma(q_0)}+\norm{W_1}_{L^{q_1}(\Omega)}^{\delta(q_1)} + \norm{W_2}_{L^{q_2}(\Omega)}^{\delta(q_2)} + \left(\norm{W_1}_{L^{q_1}(\Omega)}\norm{W_2}_{L^{q_2}(\Omega)}\right )^{\rho(q_1,q_2)}
\right) }\norm{u}_{L^{\frac{2d}{d-2}}(\omega)},
\end{multline}
with
\begin{equation*}
\! \! \! \gamma(q)\!=\! \left\{
\begin{array}{lll} \! \! \!
\ds \frac{1}{\ds \frac{3}{2} \left(1- \frac{d}{2q}\right) + \frac{1}{2q} } & \! \! \text{if } q \geq d,
\quad & \ds \!\! \delta(q)\! = \!\frac{2}{\ds 1 - \frac{3d-2}{2q}},
\medskip
\\
\ds \! \! \! \frac{1}{ \ds \left(\frac{3}{4} + \frac{1}{2d}\right) \!\! \left(2- \frac{d}{q}\right) } & \! \! \text{if } q \in \left(\ds \frac{d}{2},d\right],
\quad & \ds \!\! \rho(q_1,q_2) \!= \! \frac{1}{\ds 1 \!-\! \frac{1}{d}\! - \!\left(\frac{3}{4} \!-\! \frac{1}{2d}\right)\!\! \left(\frac{d}{q_1}\!+\! \frac{d}{q_2} \right) }.
\end{array}
\right.
\end{equation*}

\item\label{theo1.3.2} If $q_0 \in (\frac{d}{2}, \infty]$, $q_1 \in (\frac{3d}{2}, \infty]$ and $q_2\in (\frac{3d}{2}, \infty]$, the function $u$ satisfies
\begin{equation}\label{Quantitative}
\norm{u }_{L^2(\Omega)} \leqslant C e^{ C \left(\norm{V}_{L^{q_0}(\Omega)}^{\gamma(q_0)}+\norm{W_1}_{L^{q_1}(\Omega)}^{\tilde\delta(q_1)} + \norm{W_1}_{L^{q_1}(\Omega)}^{\tilde\delta(q_2)}
\right) }\norm{u}_{L^{\frac{2d}{d-2}}(\omega)},
\end{equation}
with
\begin{equation*}
\tilde \delta(q) = \frac{2}{\ds 1 - \frac{3d}{2q}}.
\end{equation*}
\end{enumerate}
\end{theo}

\begin{rema}
Note that the conditions in item 1 and in item 2 do not overlap, in the sense that there are cases in which the conditions in item 2 are satisfied while conditions in item 1 are not (for instance $q_1 = q_2 = \frac{3d}{2} + \epsilon$ with $\epsilon >0$ small), and reciprocally (for instance $q_1 = \frac{3d-2}{2} + \epsilon$ with $\epsilon>0$ small and $q_2 = \infty$).
\end{rema}

Several remarks are in order.

First, unique continuation is known to hold for general $ V\in L^{q_0}(\Omega)$, $W_1 \in L^{q_1}(\Omega; \C^d),$ and $W_2 \in L^{q_2}(\Omega; \C^d)$ for $q_0 \geq \frac{d}{2}$, $q_1\geq d$ and $q_2 \geq d$, see~\cite{Wolff-1993}, and~\cite{koch2001carleman} where even strong unique continuation is proved in that case when $q_0 > \frac{d}{2}$, $q_1 > d$ and $q_2 > d$. (These classes of integrability for the potentials are sharp, see~\cite{koch2002sharp}.)

These unique continuation results require the use of a Carleman estimate and a delicate argument from harmonic analysis inspired by~\cite{wolff1992property}, see also~\cite{koch2001carleman}. In this argument, the weight function in the Carleman estimate depends on the solution, making the quantification of unique continuation with respect to the norms of the potentials difficult to track. Another related result is the article~\cite{malinnikova2012quantitative}, which quantifies unique continuation properties for the Laplacian operator with lower order terms in the sharp integrability class, but not with respect to the norms of the potentials. In fact, since this work is based on~\cite{koch2001carleman}, as said above, it is not clear how the proof in~\cite{malinnikova2012quantitative} can be made quantitative with respect to the norms of the potentials.

Therefore, when trying to quantify the unique continuation property with respect to the norms of the lower order terms, the known results rely only on the use of a Carleman estimate, which, as pointed out in~\cite{barcelo1988weighted}, does not allow to go beyond $W_1 \in L^{\frac{3d-2}{2}}(\Omega)$. This corresponds to what is done in~\cite{davey2019quantitative,
davey2020quantitative} using $L^p$ Carleman estimate. But the results in~\cite{davey2019quantitative} describing the maximal order of vanishing of solutions of elliptic equations require $V$ and $W_1$ respectively in $L^{q_0}(\Omega)$ with $q_0 > d(3d-2)/(5d-2)$ and in $L^{q_1}(\Omega)$ with $q_1 > \frac{3d-2}{2}$, and $W_2=0$. Also note that~\cite[Theorem~1]{davey2020quantitative}, which applies when $W_1 = W_2 = 0$, exhibits the same dependence in the $L^{q_0}(\Omega)$ norm of $V$ as in Theorem~\ref{Thm-QUCP}.

Let us also mention that taking $L^2$ Carleman estimates, one cannot reach the same integrability class as in our case, see for instance~\cite{Duyckaerts2008}.

Finally, note that using a quantitative Caccioppoli inequality with singular lower order terms, see for instance~\cite[Lemma~5]{davey2019quantitative}, and Sobolev embedding, one can show that the inequalities~\eqref{Quantitative-0} and~\eqref{Quantitative} remain true by replacing $\|u \|_{L^{\frac{2d}{d-2}}(\omega)}$ by $\| u \|_{L^2(\omega_1)}$ for $\omega_1$ an open subset satisfying $\overline\omega \subset \omega_1$ (Since $\omega$ is any arbitrary non-empty open set in Theorem~\ref{Thm-QUCP}, this is of course a harmless condition).

Let us also note that one can be slightly more precise in Theorem~\ref{Thm-QUCP}, in~\eqref{Quantitative-0} and~\eqref{Quantitative}, by using the intermediate bound in~\eqref{Holder-in-Lp-tau-eta} instead of the extremal one in the proof of Theorem~\ref{Thm-QUCP}.


\subsection*{Strategy of the proof of Theorem~\ref{Thm-Carleman-General}}

In order to prove Theorem~\ref{Thm-Carleman-General}, we start with the easy geometric case of a vertical strip, with a linear weight function $x \mapsto x_1$, and a second order operator of the form $\Delta - x_1 \sum_{j = 2}^d \lambda_j \partial_{j}^2$, see Section~\ref{Sec-Carleman-Strip} for the statements.

Although this might seem at first to be a very specific case, we will check in Section~\ref{Sec-Proof-Carl-General} that this is not the case, due to the two following facts. First, if we localize the functions in a ball of radius sufficiently small, one can do a change of variables (in the spirit of the normal geodesic coordinates), such that the conjugated operator $e^{\tau \varphi} \Delta\left(e^{-\tau \varphi}\cdot \right)$ can be recast into the problem in the strip with an operator of the form $\Delta - x_1 \sum_{j = 2}^d \lambda_j \partial_{j}^2$ and the linear weight function $x\mapsto x_1$. Second, one can glue the local and boundary Carleman estimates obtained that way, and the localization terms introduced by the cut-off can be absorbed through that process if the localization is not too strong. Therefore, we have to balance the two processes, and to choose the localization rate appropriately. It turns out that a localization in balls of size $\tau^{-\frac{1}{3}}$ works.

Accordingly, most of the article in fact focuses on the proof of a Carleman estimate in the strip for an operator of the form $\Delta - x_1 \sum_{j = 2}^d \lambda_j \partial_{j}^2$ with linear weight $x \mapsto x_1$. We do that in several steps.

First, due to the specific geometric setting, one can perform a Fourier transform in the tangential variables (which are transverse to the gradient of the weight function, i.e. to the direction $e_1$), and construct explicitly a parametrix, see Section~\ref{Sec-parametrix-Strip}. In fact, this approach is inspired by~\cite{koch2001carleman,
Sogge1989uniquecontinuation} and by recent works on Carleman estimates for Laplace operator with discontinuous conductivities, for instance~\cite{LeRousseau2013interface}.

Once this is done, it is clear that we will have to get estimates on the operators appearing in the parametrix. Dealing with the Hilbertian norms can be done using classical multiplier type arguments and Parseval's identity, see Section~\ref{Sec-Hilbertian-est}.

It thus remains to understand how to get estimates on the operators appearing in the parametrix in $\scL(L^{\frac{2d}{d+2}}(\Omega),L^{\frac{2d}{d-2}}(\Omega)) $ for instance, and other operator norms involving non-Hilbertian spaces. In order to do this, we will rely on the Fourier restriction Stein--Tomas theorem, recalled in Theorem~\ref{Thm-Stein-Tomas}, see e.g.~\cite{tomas1975restriction}, \cite[Theorem~2, p.~352]{stein1993harmonic}, or~\cite[Corollary~2.2.2]{sogge2017fourier}. Similarly, as in~\cite{barcelo1988weighted,kenig1987uniform}, this approach will allow us to give an efficient manner to estimate the norm in $\scL(L^{\frac{2d}{d+2}}(\R^{d-1}),L^{\frac{2d}{d-2}}(\R^{d-1})) $ (among others) of operators given in Fourier, see Section~\ref{Subsec-Fourier-Multipliers}.

Using these results, and the explicit formula obtained for the parametrix, we manage to get $L^p$ Carleman estimates in the strip for an operator of the form $\Delta - x_1 \sum_{j = 2}^d \lambda_j \partial_{j}^2$ with linear weight $x \mapsto x_1$.

Let us finally emphasize that we made the choice of presenting the proofs in a (hopefully) pedagogical manner, and thus of giving all the technical details required to get through the whole proofs. Therefore, some parts, for instance regarding the Hilbertian estimates or the Fourier restriction theorems, might seem merely classical, but we made the choice to present them nevertheless since we did not find them in the literature in the precise version we needed.
%\medskip

\subsection*{Outline} 

The rest of the paper is as follows. Section~\ref{Sec-Carleman-Strip} is devoted to state Carleman estimates (namely Theorems~\ref{Thm-Carleman-Strip} and~\ref{Thm-Carleman-Strip-2}) in the specific case of a vertical strip, with a linear weight function $x \mapsto x_1$, and for an operator of the form $\Delta - x_1 \sum_{j = 2}^d \lambda_j \partial_{j}^2$. Section~\ref{Sec-parametrix-Strip} gives a parametrix of the conjugated operator $e^{\tau x_1}(\Delta - x_1 \sum_{j = 2}^d \lambda_j \partial_{j}^2)(e^{-\tau x_1}\, \cdot \,)$. Section~\ref{Sec-Hilbertian-est} explains how to get Hilbertian estimates on the parametrix. Section~\ref{Sec-Fourier-Restriction-Thm} then recalls Fourier restriction theorems and explains how they can be used in our context to estimate $\scL(L^{\frac{2d}{d+2}}(\R^{d-1}),L^{\frac{2d}{d-2}}(\R^{d-1})) $ norms (among others) of operators given in Fourier. We then derive all the estimates needed on the operators appearing in the parametrix in Section~\ref{Sec-Est-Parametrix} and conclude the proof of Theorems~\ref{Thm-Carleman-Strip} and~\ref{Thm-Carleman-Strip-2}. In Section~\ref{Sec-Proof-Carl-General}, we explain how to derive the proof of Theorem~\ref{Thm-Carleman-General} from Theorem~\ref{Thm-Carleman-Strip-2}. We then provide in Section~\ref{Sec-Proof-Thm-QUCP} the proof of Theorem~\ref{Thm-QUCP}. Finally, in the Appendix, we provide some reminders of classical results, namely the Hardy--Littlewood--Sobolev theorem and the stationary phase lemma (the refined version in~\cite{alazard2017stationary}). We also give the proof of a technical result of interpolation used in Section~\ref{Sec-Proof-Carl-General}.

%\medskip

\subsection*{Notations}

Here is a set of notations we will use throughout the article.

For every $x\in \R^{d}$, $x=(x_1,..,x_d)$, we set $x=(x_1,x')$, where $x'=(x_2,..,x_d) \in \R^{d-1}$.

The notations $\nabla$ and $\Delta$ respectively stand for the gradient and the Laplacian with respect to $x=(x_1,..,x_d)$, and $\nabla'=(\partial_{2},..,\partial_{{d}})$ and $\Delta'= \sum_{j=2}^{d} \partial_{j}^2$ are, respectively, the tangential gradient and Laplacian operators.

In all the document except in Section~\ref{Sec-Fourier-Restriction-Thm}, the Fourier transform is always taken to be the Fourier transform with respect to $x'=(x_2,..,x_d)$, and its dual variable $\xi' \in \R^{d-1}$ is then indexed by $\xi'=(\xi_2,..,\xi_d)$. The Fourier transform of a function $f \in \scS(\R^{d-1})$ will be denoted by $\hat f$:
\[
\hat f(\xi') = \frac{1}{(2\pi)^{\frac{d-1}{2}}} \int_{\R^{d-1}} e^{-\i x' \cdot \xi'} f(x') \, dx', \qquad \qquad\xi' \in \R^{d-1},
\]
and is extended by duality as usual to any $f \in \scS'(\R^{d-1})$.

Note that for a function $f$ defined on $\R^d$ such that $f(x_1,\cdot) \in \scS(\R^{d})$, $\hat f(x_1, \cdot)$ denotes the partial Fourier transform with respect to $x'$, that is:
\begin{equation}\label{Partial-Fourier-Transform}
\hat f(x_1,\xi') = \frac{1}{(2\pi)^{\frac{d-1}{2}}} \int_{\R^{d-1}} e^{-\i x' \cdot \xi'} f(x_1,x') \, dx', \qquad \qquad\xi' \in \R^{d-1}.
\end{equation}



\section{A Carleman estimate in a strip}\label{Sec-Carleman-Strip}

In this section, we focus on the case of a strip
\begin{equation}\label{Omega=Strip}
\Omega = (X_0,X_1) \times \R^{d-1},
\end{equation}
and on the following elliptic problem
\begin{equation}\label{Elliptic-Strip}
\left\{
\begin{array}{ll}
\Delta v-x_{1} \sum\limits_{j=2}^{d} \lambda_{j} \partial_{j}^{2} v = f_2 + f_{2_*'} + \div F & \text { in } \Omega,
\\
v\left(X_{0}, x'\right)=g (x'), &\text{ for } x' \in \R^{d-1},
\\
v\left(X_{1}, x'\right)= 0, &\text{ for } x' \in \R^{d-1},
\end{array}
\right.
\end{equation}
where
\begin{equation}\label{Reg-Source-Terms-Strip}
f_2 \in L^2(\Omega),
\quad f_{2_*'} \in L^{\frac{2d}{d+2}}(\Omega),
\quad F \in L^2\left(\Omega; \C^{d}\right),
\quad\text{ and } \quad g \in H^{\frac{1}{2}}\left(\R^{d-1}\right).
\end{equation}

To be able to solve the elliptic problem~\eqref{Elliptic-Strip}, we assume the coercivity of the operator $-\Delta + x_1 \sum_{j=2}^d \lambda_j \partial_j^2$ in $\Omega$, that is
\begin{equation}\label{Coercivity}
\exists c_0 >0, \quad \forall x_1 \in [X_0, X_1],\, \forall \xi \in \R^{d}, \quad
\frac{1}{c_0^2} |\xi|^2 \leq \sum_{j =1}^d \left(1 - x_1 \lambda_j\right) |\xi_j|^2 \leq c_0^2 |\xi|^2,
\end{equation}
where we have set $\lambda_1 = 0$ for convenience. Under condition~\eqref{Coercivity} and the integrability and regularity assumptions~\eqref{Reg-Source-Terms-Strip}, the problem~\eqref{Elliptic-Strip} has a unique solution $v \in H^1(\Omega)$.

Our goal is to prove the following Carleman estimate:
\begin{theo}\label{Thm-Carleman-Strip}
Let $\Omega$ be as in~\eqref{Omega=Strip} with $X_0 \! <\! 0\! <\! X_1$ and $\max\{|X_0|, |X_1|\} \leq 1$, and assume that the coefficients $(\lambda_j)_{j\,\in\,\{1,\cdots,d\}} \in \R^d$ satisfy $\lambda_1 = 0$, \eqref{Coercivity} and that there exist positive constants $m_*$ and $M_*$ such that
\begin{equation}\label{Sub-ellipticity-Strip}
0< m_*\leq \min_{j\,\in\,\{2,\, \cdots,\, d\}} \lambda_j
\leq \max_{j\,\in\,\{2,\, \cdots,\, d\}} \lambda_j \leq M_*.
\end{equation}
Then there exists a constant $C>0$ depending on $c_0$, $m_*$ and $M_*$ (independent of $X_0, X_1$) such that for all $(f_2, f_{2_*'},F,g)$ as in~\eqref{Reg-Source-Terms-Strip}, if the solution $v$ of~\eqref{Elliptic-Strip} satisfies $(\partial_1 v-F_1) (X_1, x') = 0$ for $x' \in \R^{d-1}$, then we have, for all $\tau \geq 1$,
\begin{multline}\label{Carleman-Strip-1}
\tau^{\frac{3}{2}} \left\| v e^{ \tau x_1}\right\|_{L^2(\Omega)} +
\tau^{\frac{1}{2}} \left\| v e^{\tau x_1} \right\|_{H^1(\Omega)}\\
\leq C\left(
\left\| f_2 e^{\tau x_1} \right\|_{L^2(\Omega)} +
\tau^{\frac{3}{4}-\frac{1}{2d}} \left\| f_{2_*'} e^{\tau x_1} \right\|_{L^{\frac{2d}{d+2}}(\Omega)}
\right.
\\
\left. +
\tau \left\| F e^{\tau x_1} \right\|_{L^2(\Omega)} +
\tau^{\frac{3}{4}} \left\| g e^{\tau X_0} \right\|_{H^{\frac{1}{2}}\left(\{X_0\} \times \R^{d-1}\right)}
\right),
\end{multline}
and
\begin{multline}\label{Carleman-Strip-2}
\tau^{\frac{3}{4}+\frac{1}{2d}} \left\| v e^{\tau x_1}\right\|_{L^{\frac{2d}{d-2}}(\Omega)}
\leq C\left(
\left\| f_2 e^{\tau x_1} \right\|_{L^2(\Omega)} +
\tau^{\frac{3}{4}+\frac{1}{2d}} \left\| f_{2_*'} e^{\tau x_1} \right\|_{L^{\frac{2d}{d+2}}(\Omega)}
\right.
\\
\left. +
\tau \left\| F e^{\tau x_1} \right\|_{L^2(\Omega)} +
\tau^{\frac{3}{4}+\frac{1}{2d}} \left\| g e^{\tau X_0} \right\|_{H^{\frac{1}{2}}\left(\{X_0\} \times \R^{d-1}\right)}
\right).
\end{multline}
\end{theo}

\begin{rema}
A solution $v$ of~\eqref{Elliptic-Strip} with $(f_2, f_{2_*'},F,g)$ as in~\eqref{Reg-Source-Terms-Strip} only belongs a priori to $H^1(\Omega)$. Therefore, trace theorems do not allow defining directly its normal trace. However, $\nabla v - F$ satisfies $\nabla v - F \in L^2(\Omega; \C^d)$ and $\div (\nabla v - F) \in L^2(\Omega) + L^{\frac{2d}{d+2}}(\Omega)$ and it is easy to check that if $R \in L^2(\Omega; \C^d)$ and $\div R \in L^2(\Omega) + L^{\frac{2d}{d+2}}(\Omega)$, then $R \cdot n$ is well-defined as an element of $H^{-\frac{1}{2}}(\partial \Omega)$, see~\cite[Theorem~III.2.43]{Boyer-Fabrie-Book}. Therefore, the trace $(\nabla v - F)\cdot n$ is well-defined as an element of $H^{-\frac{1}{2}}(\partial \Omega)$.
\end{rema}

\begin{rema}
The strict positivity of the coefficients $(\lambda_j)_{j\,\in\,\{2, \,\cdots,\, d\}}$ guaranteed by condition~\eqref{Sub-ellipticity-Strip} is the sub-ellipticity condition for the operator $-\Delta + x_1 \sum_{j =2}^d \lambda_j \partial_j^2$ with respect to the weight function $x\mapsto x_1$, see for instance~\cite[Part~1, Definition~3.30]{LeRousseaubookvolume1}.
\end{rema}

As one easily checks by working on $w$ defined by $w(x) = e^{\tau x_1} v(x)$ in $\Omega$, Theorem~\ref{Thm-Carleman-Strip} is implied by the following result, whose proof is developed from Section~\ref{Sec-parametrix-Strip} to Section~\ref{Sec-Est-Parametrix}:

\begin{theo}\label{Thm-Carleman-Strip-2}
Let $\Omega$ be as in~\eqref{Omega=Strip} with $X_0 \! <\! 0\! <\! X_1$ and $\max\{|X_0|, |X_1|\} \leq 1$, and assume that the coefficients $(\lambda_j)_{j\,\in\,\{1,\,\cdots,\,d\}} \in \R^d$ satisfy $\lambda_1 = 0$, \eqref{Coercivity}, and~\eqref{Sub-ellipticity-Strip}. Then there exist constants $C>0$ and $\tau_0 \geq 1$ depending on $c_0$, $m_*$ and $M_*$ (independent of $X_0, X_1$), such that for all $(f_2, f_{2_*'},F,g)$ as in~\eqref{Reg-Source-Terms-Strip}, if the solution $w$ of
\begin{equation}\label{Elliptic-Strip-w}
\left\{
\begin{array}{ll}
\Delta w-x_{1} \sum\limits_{j=2}^{d} \lambda_{j} \partial_{j}^{2} w - 2 \tau \partial_1 w + \tau^2 w = f_2 + f_{2_*'} + \div F & \text { in } \Omega,
\\
w\left(X_{0}, x'\right)=g (x'), &\text{ for } x' \in \R^{d-1},
\\
w\left(X_{1}, x'\right)= 0, &\text{ for } x' \in \R^{d-1},
\end{array}
\right.
\end{equation}
satisfies $(\partial_1 w - F_1)(X_1, x') = 0$ for $x' \in \R^{d-1}$, then for all $\tau \geq \tau_0$,
\begin{multline}\label{Carleman-Strip-w-1}
\tau^{\frac{3}{2}} \| w\|_{L^2(\Omega)} +
\tau^{\frac{1}{2}} \| \nabla w \|_{L^2(\Omega)}
\\ \leq C\left(
\left\| f_2 \right\|_{L^2(\Omega)} +
\tau^{\frac{3}{4}-\frac{1}{2d}} \| f_{2_*'} \|_{L^{\frac{2d}{d+2}}(\Omega)} +
\tau \| F \|_{L^2(\Omega)} +
\tau^{\frac{3}{4}} \| g \|_{H^{\frac{1}{2}}\left(\{X_0\} \times \R^{d-1}\right)}
\right),
\end{multline}
and
\begin{multline}\label{Carleman-Strip-w-2}
\tau^{\frac{3}{4}+\frac{1}{2d}} \| w\|_{L^{\frac{2d}{d-2}}(\Omega)}
\\
\leq C\left(
\| f_2 \|_{L^2(\Omega)} +
\tau^{\frac{3}{4}+\frac{1}{2d}} \| f_{2_*'} \|_{L^{\frac{2d}{d+2}}(\Omega)}
\right.
\left. +
\tau \| F \|_{L^2(\Omega)} +
\tau^{\frac{3}{4}+\frac{1}{2d}} \| g \|_{H^{\frac{1}{2}}(\{X_0\} \times \R^{d-1})}
\right).
\end{multline}
\end{theo}

In fact, the correspondence between Theorem~\ref{Thm-Carleman-Strip} and~\ref{Thm-Carleman-Strip-2} is given by
\[
(f_2, f_{2_*'},F,g)
\to
\left((f_2 - \tau F_1) e^{\tau x_1}, f_{2_*'} e^{\tau x_1}, F e^{\tau x_1}, g e^{\tau X_0}
\right).
\]
Theorem~\ref{Thm-Carleman-Strip} and Theorem~\ref{Thm-Carleman-Strip-2} are then completely equivalent, and we thus focus only on the latter.

\section{Construction of the parametrix in the case of a strip}\label{Sec-parametrix-Strip}

The goal of this section is to explicitly construct the solution $w$ of~\eqref{Elliptic-Strip-w} for $\tau \geq 1$, $(f_2, f_{2_*'},F,g)$ as in~\eqref{Reg-Source-Terms-Strip}, under the assumptions that the domain $\Omega$ is a strip as in~\eqref{Omega=Strip}, and the coefficients $(\lambda_j)_{j\, \in\, \{1,\,\cdots,\,d\}} \in \R^d$ satisfy $\lambda_1 = 0$, \eqref{Coercivity} and~\eqref{Sub-ellipticity-Strip}.

In order to do that, we take the partial Fourier transform in the variable $x' \in \R^{d-1}$ of~\eqref{Elliptic-Strip-w} with dual variable $\xi' \in \R^{d-1}$:
\begin{equation}\label{Elliptic-Strip-w-Fourier}
\begin{cases}
\ds \left(\partial_1 - \tau\right)^2 \hat w- \sum_{j=2}^{d} (1 - x_{1}\lambda_j) \xi_j^{2} \hat w 
\\
\mkern78mu= \hat f_2 + \hat f_{2_*'} + \partial_1 \hat F_1 + \i \sum\limits_{j = 2}^d \xi_j \hat F_j & \text { for } (x_1, \xi') \in \Omega,\\
\hat w\left(X_{0}, \xi'\right)=\hat g (\xi '), &\text{ for } \xi' \in \R^{d-1},
\\
\hat w\left(X_{1}, \xi'\right)= 0, &\text{ for } \xi' \in \R^{d-1}.
\end{cases}
\end{equation}
We then show the following:

\break
\begin{prop}\label{Prop-parametrix-Strip-1}
Let $\Omega$ be as in~\eqref{Omega=Strip}, and assume that the coefficients $(\lambda_j)_{j \,\in\,\{1,\,\cdots,\,d\}} \in \R^d$ satisfy $\lambda_1 = 0$, \eqref{Coercivity} and~\eqref{Sub-ellipticity-Strip}. We introduce the function $\psi : \overline\Omega \to \R$ defined by
\begin{equation}\label{Def-Psi}
\psi(x_1, \xi')=\sqrt{ \sum_{j=2}^{d}\left(1-x_{1} \lambda_{j}\right) \xi_{j} ^2},
\qquad \qquad x_1 \in [X_0,X_1], \ \xi'\in \R^{d-1}.
\end{equation}

For all $\tau \geq 1$, for all $(f_2, f_{2_*'},F,g)$ as in~\eqref{Reg-Source-Terms-Strip}, the solution $w$ of~\eqref{Elliptic-Strip-w} formally satisfies
\begin{multline}\label{parametrix}
w = K_{\tau,0} \left(f_2 + f_{2_*'}\right) + K_{\tau,1}(F_1) \\
+ \sum_{j = 2}^d K_{\tau,j} (F_j) + R_\tau (w) + G_\tau(g) + H_\tau\big((\partial_1 w -F_1)(X_1, \cdot)\big),
\end{multline}
where, using the partial Fourier transform~\eqref{Partial-Fourier-Transform}, the operators $K_{\tau,j}$, for $j \in \{0,\linebreak\cdots, d\}$, and $R_\tau$ are formally defined for $f$ depending on $(x_1, x') \in \Omega$, by
\begin{align}
\widehat {K_{\tau,j} f}(x_1, \xi') & = \int_{y_1\, \in\,(X_0,X_1)} k_{\tau,j}(x_1,y_1, \xi') \widehat f(y_1, \xi') \, dy_1,
\qquad \qquad(x_1, \xi') \in \Omega,\label{Def-K-tau-j-Fourier}
\\
\widehat {R_{\tau} f}(x_1, \xi') & = \int_{y_1 \,\in\,(X_0,X_1)} r_{\tau}(x_1,y_1, \xi') \widehat f(y_1, \xi') \, dy_1,
\mkern80mu (x_1, \xi') \in \Omega,\label{Def-R-tau-Fourier}
\end{align}
with kernels given, for $(x_1,y_1, \xi') \in [X_0,X_1]^2 \times \R^{d-1}$, by
\begin{multline}\label{Def-k-tau-0-x1-y1-xi'}
k_{\tau,0}(x_1, y_1, \xi')\\
= -1_{\psi(x_1,\xi') \,>\,\tau} \int_{X_0}^{\min \{ x_1, y_1\}} e^{-\tau(y_1-x_1)-\int_{\tilde{x}_1}^{x_1} \psi\left(\tilde{y}_1,\xi'\right) \, d\tilde{y}_1 - \int_{{\tilde{x}_1}}^{y_1} \psi\left(\tilde{y}_1,\xi'\right) \,d \tilde{y}_1 } \, d \tilde{x}_1
\\
+1_{\psi(x_1,\xi')\,\leq\,\tau} 1_{y_1\,>\,x_1} \int_{x_1}^{y_1} e^{-\tau (y_1-x_1)+\int_{x_1}^{\tilde{x}_1} \psi\left(\tilde{y}_1,\xi'\right) \, d \tilde{y}_1-\int_{{\tilde{x}_1}}^{y_1} \psi\left(\tilde{y}_1,\xi'\right) \, d y_1} \, d \tilde{x}_1,
\end{multline}
\begin{multline}\label{Def-k-tau-1-x1-y1-xi'}
k_{\tau,1} (x_1,y_1, \xi') = - 1_{\psi(x_1, \xi')\, \leq\,\tau}\; 1_{y_1\,>\,x_1} e^{-\tau(y_1-x_1)+\int_{x_1}^{y_1} \psi\left(\tilde{y}_1,\xi'\right) \, d \tilde{y}_1}
\\
+ 1_{\psi(x_1,\xi')\,>\,\tau}\; 1_{y_1\,<\,x_1} e^{\tau(x_1-y_1) - \int_{y_1}^{x_1} \psi\left(\tilde{y}_1,\xi'\right) \, d \tilde{y}_1} + k_{\tau,0}(x_1,y_1,\xi') \big(\tau + \psi(y_1,\xi')\big),
\end{multline}
\begin{equation}\label{Def-k-tau-j-x1-y1-xi'}
\begin{split}
k_{\tau,j}(x_1, y_1, \xi') & = \i \xi_j k_{\tau,0}(x_1,y_1, \xi'), \qquad\qquad \qquad j \in \{2, \cdots, d\},
\\
r_\tau(x_1,y_1, \xi) & = k_{\tau, 0}(x_1,y_1, \xi') \partial_1 \psi(y_1,\xi').
\end{split} 
\end{equation}

The operators $G_\tau$ and $H_\tau$ are formally given in Fourier for $g_0 \in \scS(\R^{d-1})$ by
\begin{equation}\label{Def-G-tau-Fourier}
\begin{aligned}
\widehat {G_{\tau} g_0}(x_1, \xi') & = g_{\tau}(x_1, \xi') \widehat g_0(\xi'),
\qquad \qquad(x_1, \xi') \in \Omega,
\\
\widehat {H_{\tau} g_0}(x_1, \xi') & = h_\tau(x_1, \xi') \widehat g_0(\xi'),
\qquad \qquad(x_1, \xi') \in \Omega,
\end{aligned}
\end{equation}
where $g_\tau$ and $r_\tau$ are given, for $(x_1, \xi') \in [X_0,X_1] \times \R^{d-1}$, by
\begin{equation}\label{Def-g-tau-x1-xi'}
\begin{split}
g_{\tau}(x_1, \xi') & = 1_{\psi(x_1,\xi')\, \geq\,\tau} e^{\tau (x_1-X_0)- \int_{X_0}^{x_1} \psi\left(\tilde{y}_1, \xi'\right) \, d \tilde{y}_1 },
\\
h_\tau(x_1,\xi') & = 1_{\psi(x_1,\xi')\,>\,\tau} \int_{X_0}^{x_1} e^{-\tau(X_1-x_1) - \int_{\tilde x_1}^{x_1} \psi\left(\tilde y_1, \xi'\right) \, d\tilde y_1 - \int_{\tilde x_1}^{X_1} \psi\left(\tilde y_1, \xi'\right) \, d\tilde y_1 } d\tilde x_1
\\
& \qquad - 1_{\psi(x_1,\xi')\,\leq\,\tau} \int_{x_1}^{X_1} e^{-\tau(X_1-x_1) +\int_{x_1}^{\tilde x_1} \psi\left(\tilde y_1, \xi'\right) \, d\tilde y_1 - \int_{\tilde x_1}^{X_1} \psi\left(\tilde y_1, \xi'\right) \, d\tilde y_1} d\tilde x_1.
\end{split}
\end{equation}
\end{prop}

\begin{rema}
We emphasize that Proposition~\ref{Prop-parametrix-Strip-1} is formal. We will prove later, in Theorem~\ref{Thm-Hilbertian-est}, Proposition~\ref{Prop-parametrix-Strip-2} and in Proposition~\ref{Prop-OpNorm-Est-K0}, that the operators $K_{\tau,0}$, $(K_{\tau,j})_{j \,\in\,\{1,\,\cdots,\,d\} }$, $R_\tau$, $G_\tau$ and $H_{\tau}$ respectively belong to $\scL(L^{2}(\Omega) + L^{\frac{2d}{d+2}}(\Omega); L^2(\Omega))$, $(\scL(L^{2}(\Omega) ; L^2(\Omega)))^d$, $\scL(L^{2}(\Omega); L^2(\Omega))$, $\scL(H^{1}(\Omega); L^2(\Omega))$, $\scL(H^{\frac{1}{2}}(\partial\Omega); L^2(\Omega))$, and $\scL(H^{-\frac{1}{2}}(\Omega); L^2(\Omega))$. A simple density argument would then allow to justify rigorously formula~\eqref{parametrix}.
\end{rema}



\begin{proof}
The basic strategy of proof of Proposition~\ref{Prop-parametrix-Strip-1} consists in the factorization of the operator $(\partial_1 - \tau)^2 - \sum_{j = 2}^d (1- x_1 \lambda_j) \xi_j^2$:
\[
(\partial_1 - \tau)^2 - \sum_{j = 2}^d (1- x_1 \lambda_j) \xi_j^2 =
\left(
\partial_1 - \tau- \psi(x_1, \xi')
\right)
\left(
\partial_1 - \tau + \psi(x_1, \xi')
\right) - \partial_1 \psi(x_1, \xi'),
\]
where $\psi$ is the function introduced in~\eqref{Def-Psi}, and the last term should be seen as a correction term.

We thus set, for all $ x_1 \in [X_0, X_1], \, \xi' \in \R^{d-1}, $
\[
\widehat H (x_1, \xi') = \hat f_2(x_1, \xi') + \hat f_{2_*'}(x_1, \xi') + \i \sum_{j=2}^d \xi_j \hat F_j(x_1, \xi') + \partial_1 \psi(x_1, \xi') \widehat w(x_1, \xi'),
\]
so that equation~\eqref{Elliptic-Strip-w-Fourier}$_{(1)}$ can be rewritten as
\begin{align*}
\left(
\partial_1 - \tau- \psi(x_1, \xi')
\right)
\left(
\partial_1 - \tau + \psi(x_1, \xi')
\right)
\widehat{w} =
\widehat{H} + \partial_1 \widehat F_1, \quad \text{ in } \Omega.
\end{align*}
Accordingly, introducing the additional unknown $\widehat{z}(x_1, \xi') \!\! =\!\!(\partial_1 - \tau + \psi(x_1, \xi')) \widehat w(x_1,\xi')$, equation~\eqref{Elliptic-Strip-w-Fourier} can be rewritten as a system of two first order ODE indexed by $\xi' \in \R^{d-1}$:
\begin{equation}\label{Eq-System-w-z}
\left\{
\begin{array}{ll}
\left(\partial_1-\tau+\psi(x_1,\xi') \right) \widehat{w}(x_1,\xi')=\widehat{z}(x_1,\xi') &\text{in } \Omega,
\\
\left(\partial_1-\tau-\psi(x_1,\xi') \right)\widehat{z}(x_1,\xi')=\widehat{H}(x_1,\xi') + \partial_1 \hat F_1(x_1, \xi') & \text{in } \Omega,
\\
\widehat{w}(X_0,\xi')= \widehat{g}(\xi'), &\text{on } \R^{d-1},
\\
\widehat{w}(X_1,\xi')= 0 &\text{on } \R^{d-1}.
\end{array}
\right.
\end{equation}

Let $\xi' \in \R^{d-1}$. Solving~\eqref{Eq-System-w-z}$_{2}$ from the right, which can be done easily by working on $(\widehat z - \hat F_1)(\cdot, \xi')$, by Duhamel’s formula we get, for $x_1 \in (X_0,X_1)$,
\begin{multline}\label{Explicit-z}
\widehat{z}(x_1, \xi') = e^{-\tau (X_1 -x_1)- \int_{x_1}^{X_1} \psi\left(\tilde{y}_1,\xi'\right) \, d \tilde{y}_1} \left(\partial_1 \widehat w- \hat F_1\right)(X_1, \xi') +
\hat F_1(x_1, \xi')
\\
-\int_{x_1}^{X_1} e^{-\tau (y_1-x_1)-\int_{x_1}^{y_1} \psi\left(\tilde{y}_1,\xi'\right) \, d \tilde{y}_1} \left(\widehat{H}(y_1, \xi') + (\tau + \psi(y_1, \xi')) \widehat F_1(y_1, \xi')\right) \, d y_1,
\end{multline}
where we did the additional remark that, using~\eqref{Eq-System-w-z}$_{(1,4)}$, $\widehat z(X_1, \xi') = \partial_1 \widehat w (X_1, \xi')$.

We then focus on the equations~\eqref{Eq-System-w-z}$_{(1,3,4)}$ giving $\hat w(\cdot, \xi')$ in terms of $\hat z(\cdot, \xi')$. One should notice here that~\eqref{Eq-System-w-z}$_{(3,4)}$ give two boundary conditions for a first order equation. Therefore, we do a choice when solving~\eqref{Eq-System-w-z}$_{(1)}$ based on the fact that we want formulae involving only exponentials of nonpositive numbers. In order to do such a choice, we analyse the sign of the function $ x_1 \mapsto - \tau + \psi(x_1, \xi') $. Due to conditions~\eqref{Coercivity} and~\eqref{Sub-ellipticity-Strip}, the function $x_1 \mapsto \psi(x_1, \xi')$ is strictly decreasing on $[X_0, X_1]$. Therefore, the function $ x_1 \mapsto - \tau + \psi(x_1, \xi') $ can vanish only once on $[X_0, X_1]$, and if it vanishes at some point $x_{\tau,\xi'} \in [X_0,X_1]$, it is positive in $[X_0, x_{\tau,\xi'})$ and negative for $x_1 \in (x_{\tau,\xi'}, X_1]$.

Accordingly, for $x_1 \in [X_0, X_1]$ such that $\psi(x_1, \xi') \leq \tau$, we use the formula
\begin{equation*}
\widehat{w}(x_1, \xi') = -\int_{x_1}^{X_1} e^{-\tau\left(\tilde{x}_1-x_1\right)+\int_{x_1}^{\tilde{x}_1} \psi\left(\tilde{y}_1,\xi'\right) \, d \tilde{y}_1} \widehat{z}\left(\tilde{x}_1, \xi'\right) \, d \tilde{x}_1,
\end{equation*}
while for $x_1 \in [X_0, X_1]$ such that $\psi(x_1, \xi') >\tau$, we use the formula
\begin{equation*}
\widehat{w}(x_1, \xi') = e^{\tau(x_1-X_0)- \int_{X_0}^{x_1} \psi\left(\tilde{y}_1, \xi'\right) \, d \tilde{y}_1 } \widehat{g}(\xi') +
\int_{X_0}^{x_1} e^{\tau(x_1-\tilde{x}_1) - \int_{\tilde{x}_1}^{x_1} \psi\left(\tilde{y}_1,\xi'\right) \, d \tilde{y}_1} \widehat{z}(\tilde{x}_1, \xi') \, d \tilde{x}_1.
\end{equation*}
These two formulae can be written in one under the form
\begin{multline}\label{Explicit-w-in-terms-of-z}
\hat w(x_1, \xi') = - 1_{\psi(x_1, \xi')\, \leq \,\tau} \int_{x_1}^{X_1} e^{-\tau(\tilde{x}_1-x_1)+\int_{x_1}^{\tilde{x}_1} \psi\left(\tilde{y}_1,\xi'\right) \, d \tilde{y}_1} \widehat{z}(\tilde{x}_1, \xi') \, d \tilde{x}_1
+ 1_{\psi(x_1,\xi')\, >\,\tau}
\\
 \left(e^{\tau(x_1-X_0)- \int_{X_0}^{x_1} \psi\left(\tilde{y}_1, \xi'\right) \, d \tilde{y}_1 } \widehat{g}(\xi') +
\int_{X_0}^{x_1} e^{\tau(x_1-\tilde{x}_1) - \int_{\tilde{x}_1}^{x_1} \psi\left(\tilde{y}_1,\xi'\right) \, d \tilde{y}_1} \widehat{z}(\tilde{x}_1, \xi') \, d \tilde{x}_1 \right)\!.
\end{multline}
The formulae given by Proposition~\ref{Prop-parametrix-Strip-1} are then deduced by putting together formulae~\eqref{Explicit-z} and~\eqref{Explicit-w-in-terms-of-z}. Details are left to the reader.
\end{proof}

\section{Hilbertian estimates}\label{Sec-Hilbertian-est}

The goal of this section is to prove the following result:

\begin{theo}\label{Thm-Hilbertian-est}
Let $\Omega$ be as in~\eqref{Omega=Strip} with $X_0 \! <\! 0\! <\! X_1$ and $\max\{|X_0|, |X_1|\} \leq 1$, and assume that the coefficients $(\lambda_j)_{j\,\in\,\{1,\,\cdots,\,d\}} \in \R^d$ satisfy $\lambda_1 = 0$, \eqref{Coercivity} and~\eqref{Sub-ellipticity-Strip}.

Then there exist constants $C>0$ and $\tau_0 \geq 1$ depending on $c_0$, $m_*$ and $M_*$ (independent of $X_0, X_1$) such that for all $\tau \geq \tau_0$, for all $(f_2,F,g_0, g_1) \in L^2(\Omega) \times L^2(\Omega; \C^d) \times H^{\frac{1}{2}}(\R^{d-1})\times H^{-\frac{1}{2}}(\R^{d-1})$, the function $w$ given by
\begin{equation}\label{parametrix-f2_*=0}
w = K_{\tau,0} (f_2) + K_{\tau,1}(F_1) + \sum_{j = 2}^d K_{\tau,j} (F_j) + G_\tau(g_0) + H_\tau(g_1),
\end{equation}
where $(K_{\tau,j})_{j\,\in\,\{0,\,\cdots,\,d\} }$, $G_\tau$ and $H_\tau$ are given by Proposition~\ref{Prop-parametrix-Strip-1}, satisfies:
\begin{multline}\label{Carleman-L2}
\tau^{\frac{3}{2}} \| w \|_{L^2(\Omega)} + \tau^{\frac{1}{2}} \| \nabla' w \|_{L^2(\Omega)}
\leq C \| f_2 \|_{L^2(\Omega)}
\\
+ C\tau \| F\|_{L^2(\Omega)} + C \tau \| g_0 \|_{H^{\frac{1}{2}}\left(\R^{d-1}\right)} + C\tau \| g_1\|_{H^{-\frac{1}{2}}\left(\R^{d-1}\right)}.
\end{multline}
Besides, for $w \in L^2(\Omega)$ satisfying $\nabla' w \in L^2(\Omega)$, for all $\tau \geq \tau_0$, $R_\tau (w)$ introduced in Proposition~\ref{Prop-parametrix-Strip-1}, satisfies
\begin{equation}\label{Est-R-tau-L2}
\tau^{\frac{3}{2}} \| R_\tau (w) \|_{L^2(\Omega)} + \tau^{\frac{1}{2}} \| \nabla' R_\tau (w) \|_{L^2(\Omega)}
\leq C \| \nabla' w \|_{L^2(\Omega)}.
\end{equation}
\end{theo}

\begin{rema}
The estimate~\eqref{Carleman-L2} above gives estimates on the norm of the operators $K_{\tau,j}$ for $j \in \{0,\, \cdots,\,d\}$ as operators in $\scL(L^2(\Omega), L^2(\Omega))$, and on the norm of the operators $G_{\tau}$ and $H_\tau$ in, respectively $\scL(H^{\frac{1}{2}}(\R^{d-1}), L^2(\Omega))$ and $\scL(H^{-\frac{1}{2}}(\R^{d-1})$, $ L^2(\Omega))$, but with no claim of optimality. In fact, as we will see later in Theorem~\ref{Thm-Improved-Hilbert-Est-G-tau}, the estimate on the operator norm of $G_\tau$ is not sharp. On the contrary, the ones on $K_{\tau,j}$ for $j \in \{0,\,\cdots,\,d\}$ are sharp, see for instance~\cite[Part~1, Theorems~4.4, 4.5 and Theorem~4.10]{LeRousseaubookvolume1} regarding $K_{\tau,0}$.
\end{rema}

In this section and in Section~\ref{Sec-Est-Parametrix}, all the constants $C$ depend only on $c_0$ in~\eqref{Coercivity} and $m_*, M_*$ in~\eqref{Sub-ellipticity-Strip}, and this fact will not be mentioned in the sequel.

\begin{proof}
We first remark that $w$ as in~\eqref{parametrix-f2_*=0} satisfies by construction the following version of~\eqref{Eq-System-w-z}:
\begin{equation}\label{Eq-System-w-z-bis}
\left\{
\begin{array}{ll}
\left(\partial_1-\tau+\psi(x_1,\xi') \right) \widehat{w}(x_1,\xi')=\widehat{z}(x_1,\xi') &\text{for } (x_1, \xi') \in \Omega_{1,\tau},
\\
\left(\partial_1-\tau-\psi(x_1,\xi') \right)(\widehat{z} - \widehat F_1)(x_1,\xi')=\widehat f_2(x_1, \xi')
\\
\qquad\qquad + \i \sum\limits_{j=2}^d \xi_j \hat F_j(x_1, \xi') + (\tau + \psi(x_1, \xi')) \hat F_1(x_1, \xi') &\text{for } (x_1, \xi') \in \Omega,
\\
\widehat{w}(X_0,\xi')= \widehat{g_0}(\xi'), &\text{if } \psi(X_0, \xi') > \tau,
\\
\widehat{w}(X_1,\xi')= 0 &\text{if } \psi(X_1, \xi') \leq \tau,
\\
(\widehat z -\widehat F_1)(X_1, \xi') = \hat g_1 (\xi'), &\text{for } \xi' \in \R^{d-1},
\end{array}
\right.
\end{equation}
where $\Omega_{1,\tau} = (\{ \psi(x_1, \xi') > \tau\} \cap \Omega) \cup (\{ \psi(x_1, \xi') < \tau\} \cap \Omega)$. Note that, given $\xi' \in \R^{d-1}$, due to the conditions~\eqref{Coercivity} and~\eqref{Sub-ellipticity-Strip}, there exists at most one element $x_1^*(\xi') \in [X_0, X_1]$ such that $\psi(x_1^*(\xi'), \xi') = 0$. Consequently, given $\xi' \in \R^{d-1}$, $\{ x_1 \in (X_0, X_1), \, (x_1, \xi') \in \Omega_{1, \tau} \}$ is either the whole interval $(X_0, X_1)$ or the union of two disjoint intervals $(X_0, x_1^*(\xi')) \cup (x_1^*(\xi'), X_1)$.


Since this system is now a family of ODE indexed by the tangential Fourier parameter $\xi' \in \R^{d-1}$, from now on, we see $\xi' \in \R^{d-1}$ as a free parameter.

We then perform estimates on $\hat z(\cdot, \xi') $ using~\eqref{Eq-System-w-z-bis}$_{(2,5)}$ by setting $\tilde z(\cdot, \xi') = \hat z(\cdot, \xi') - \hat F_1(\cdot, \xi')$, which satisfies:
\begin{equation*}\label{Eq-tilde-zs}
\left\{
\begin{array}{ll}
\left(\partial_1-\tau-\psi(x_1,\xi') \right)\tilde {z}(x_1,\xi')=\widehat f_2(x_1, \xi') \\
\qquad\qquad+ \i \sum\limits_{j=2}^d \xi_j \hat F_j(x_1, \xi') + (\tau + \psi(x_1, \xi')) \hat F_1(x_1, \xi') & \text{in } \Omega,
\\
\tilde z(X_1, \xi') = \hat g_1 (\xi'), &\text{on } \R^{d-1}.
\end{array}
\right.
\end{equation*}
We then use a multiplier approach, taking the square of both sides and integrating in $x_1$:
\begingroup
\allowdisplaybreaks
\begin{align*}
\int_{X_0}^{X_1}& \left(\left|\partial_1 \tilde z(x_1, \xi')\right|^2 + \left((\tau + \psi(x_1, \xi'))^2 + \partial_1 \psi(x_1, \xi') \right) \left|\tilde z(x_1, \xi')\right|^2\right) \, dx_1
\\
& +(\tau + \psi(X_0, \xi')) |\tilde z(X_0, \xi')|^2 = (\tau + \psi(X_1, \xi')) |\tilde z(X_1, \xi')|^2
\\
& + \int_{X_0}^{X_1} \left| \widehat f_2(x_1, \xi') + \i \sum_{j=2}^d \xi_j \hat F_j(x_1, \xi') + (\tau + \psi(x_1, \xi')) \hat F_1(x_1, \xi')\right|^2 \, dx_1
\\
\leq&\, (\tau + \psi(X_1, \xi')) |\hat g_1(\xi')|^2 + (d+1) \int_{X_0}^{X_1} \left| \widehat f_2(x_1, \xi') \right|^2 \, dx_1
\\
&+ (d+1) \sum_{j=2}^d \int_{X_0}^{X_1} \left|\xi_j \hat F_j(x_1, \xi')\right|^2 \, dx_1\\
&+ (d+1) \int_{X_0}^{X_1} \left(\tau + \psi(x_1, \xi')\right)^2 \left|\hat F_1(x_1, \xi') \right|^2 \, dx_1.
\end{align*}
\endgroup
Note that, within the setting of Theorem~\ref{Thm-Hilbertian-est}, the function $\psi$ defined as in~\eqref{Def-Psi} is such that there exists $C_1 >0$ for which
\begin{multline}\label{Equivalences-Psi-x-xi-prime}
\forall x_1 \in [X_0,X_1],\, \forall \xi' \in \R^{d-1},
\frac{|\xi'| }{C_1} \leq \psi(x_1, \xi')
\leq C_1 |\xi'|,\\
 - C_1 |\xi'| \leq \partial_1 \psi(x_1, \xi') \leq - \frac{|\xi'| }{C_1}.
\end{multline}
Accordingly, for $\tau \geq \tau_0$ large enough, there exists $C>0$ such that for all $\xi' \in \R^{d-1}$ and $x_1 \in [X_0, X_1]$,
\begin{align*}
&
\frac{1}{C} (\tau + |\xi'|)^2
\leq (\tau + \psi(x_1, \xi'))^2 + \partial_1 \psi(x_1, \xi') \leq {C} (\tau + |\xi'|)^2,
\\
& \frac{1}{C} (\tau + |\xi'|) \leq \tau + \psi(x_1, \xi') \leq C (\tau + |\xi'|).
\end{align*}
Therefore, the above estimate yields:
\begin{multline*}
\int_{X_0}^{X_1} (\tau + |\xi'|)^2 |\tilde z(x_1, \xi')|^2 \, dx_1
\leq C (\tau + |\xi'|) |\hat g_1(\xi')|^2 + C \int_{X_0}^{X_1} | \widehat f_2(x_1, \xi') |^2 \, dx_1
\\
+ C \sum_{j=2}^d \int_{X_0}^{X_1} \left|\xi_j \hat F_j(x_1, \xi')\right|^2 \, dx_1 + C \int_{X_0}^{X_1} (\tau + |\xi'|)^2 \left|\hat F_1(x_1, \xi') \right|^2 \, dx_1.
\end{multline*}
Recalling that $\tilde z(\cdot, \xi') = \hat z(\cdot, \xi') -\hat F_1(\cdot, \xi')$, we obtain
\begin{multline*}
\int_{X_0}^{X_1} (\tau + |\xi'|)^2 |\hat z(x_1, \xi')|^2 \, dx_1
\leq C (\tau + |\xi'|) |\hat g_1(\xi')|^2 + C \int_{X_0}^{X_1} \left| \widehat f_2(x_1, \xi') \right|^2 \, dx_1
\\
+ C \sum_{j=2}^d \int_{X_0}^{X_1} \left|\xi_j \hat F_j(x_1, \xi')\right|^2 \, dx_1 + C \int_{X_0}^{X_1} (\tau + |\xi'|)^2 \left|\hat F_1(x_1, \xi') \right|^2 \, dx_1.
\end{multline*}
We then derive estimates on $\hat w(\cdot, \xi') $ from the equation~\eqref{Eq-System-w-z-bis}$_{(1,3,4)}$, again by taking the square of both sides of~\eqref{Eq-System-w-z-bis}$_{(1)}$ and doing integration by parts. If for all $x_1 \in [X_0, X_0]$, $\psi(x_1, \xi') \neq \tau$, we do the computation at once by doing the integration by parts on $(X_0, X_1)$, and if there exists $x_1^*(\xi') \in [X_0, X_1]$ such that $\psi(x_1^*(\xi'), \xi') = \tau$ (recall that such an $x_1^*(\xi')$ is necessarily unique), we do the computations on $[X_0, x_1^*(\xi'))$ and on $(x_1^*(\xi'), X_1]$, and we sum the estimates. In this latter case, $\partial_1 \hat w$ should be interpreted as $1_{x_1\,<\,x_1^*(\xi')} \partial_1 \hat w(x_1, \xi') + 1_{x_1\,>\,x_1^*(\xi')} \partial_1 \hat w(x_1, \xi')$. There is a priori no reason that this coincides with the derivative of $\hat w(\cdot, \xi')$ in the sense of $\scD'(X_0,X_1)$, which would require some continuity conditions on $\hat w(x_1^*(\xi')^\pm, \xi')$. We get:
\begin{multline*}
\int_{X_0}^{X_1}\left(\left|\partial_1 \hat w(x_1, \xi')\right|^2 + \left((\tau - \psi(x_1, \xi'))^2 - \partial_1 \psi(x_1, \xi') \right) |\hat w(x_1, \xi')|^2\right) \, dx_1
\\
+ (\psi(X_1, \xi') - \tau) \left|\hat w(X_1, \xi')\right|^2 - (\psi(X_0, \xi') - \tau) \left|\hat w(X_0, \xi')\right|^2 =
\int_{X_0}^{X_1}\left| \hat z(x_1, \xi')\right|^2 \, dx_1.
\end{multline*}
Accordingly, using the boundary conditions~\eqref{Eq-System-w-z-bis}$_{(3,4)}$ and~\eqref{Equivalences-Psi-x-xi-prime},
\begin{multline*}
\int_{X_0}^{X_1}\left(\left|\partial_1 \hat w(x_1, \xi')\right|^2 + \left((\tau - \psi(x_1, \xi'))^2 + |\xi'| \right) \left|\hat w(x_1, \xi')\right|^2\right) \, dx_1
\\
\leq C (\psi(X_0, \xi') - \tau) 1_{\psi(X_0, \xi')\,>\,\tau} |\hat g_0(\xi')|^2 + C \int_{X_0}^{X_1}| \hat z(x_1, \xi')|^2 \, dx_1.
\end{multline*}
Therefore, there exists $C>0$ such that for all $\tau \geq \tau_0$ and $\xi' \in \R^{d-1}$,
\begin{multline*}
\int_{X_0}^{X_1}\left(\left|\partial_1 \hat w(x_1, \xi')\right|^2 + \left((\tau - \psi(x_1, \xi'))^2 + |\xi'| \right) \left|\hat w(x_1, \xi')\right|^2\right) \, dx_1
\\
\leq C (\psi(X_0, \xi') - \tau) 1_{\psi(X_0, \xi')\,>\,\tau} |\hat g_0(\xi')|^2 + C \frac{1}{\tau + |\xi'| } \left|\hat g_1(X_1, \xi')\right|^2
\\
+ \frac{C}{(\tau + |\xi'|)^2}\!\! \left(\!\int_{X_0}^{X_1}\left| \widehat f_2(x_1, \xi') \right|^2 dx_1
+ |\xi'|^2\sum_{j=2}^d \int_{X_0}^{X_1} 
\left| \hat F_j(x_1, \xi')\right|^2 dx_1\right)\\ + C  \int_{X_0}^{X_1} \!\left|\hat F_1(x_1, \xi')\right|^2 dx_1.
\end{multline*}
We finally use that there exists a constant such that for all $\xi' \in \R^{d-1}$, $\tau \geq \tau_0$, and $x_1 \in [X_0, X_1]$,
\begin{multline*}
 \left((\tau - \psi(x_1, \xi'))^2 + |\xi'| \right) \geq \frac{1}{C} \left(\tau + \frac{|\xi'|^2}{\tau} \right),
\quad (\psi(X_0, \xi') - \tau) 1_{\psi(X_0, \xi')\,>\,\tau}
\leq C |\xi'|,\\
\frac{1}{\tau + |\xi'| } \leq \frac{1}{1+|\xi'|},
\quad
\frac{1}{(\tau + |\xi'|)^2 } \leq \frac{1}{\tau^2},
\quad
\frac{|\xi'|^2 }{(\tau + |\xi'|)^2 } \leq 1,
\end{multline*}
so that
\begin{multline*}
\int_{X_0}^{X_1}\left(\tau \left|\hat w(x_1, \xi')\right|^2 + \left|\partial_1 \hat w(x_1, \xi')\right|^2 + \frac{1}{\tau} |\xi'|^2 \left|\hat w(x_1, \xi')\right|^2\right) \, dx_1
\\
\leq C |\xi'|  \left|\hat g_0(\xi')\right|^2 +
\frac{C}{1\! + |\xi'| } \left|\hat g_1(\xi')\right|^2 + \frac{C}{\tau^2}\int_{X_0}^{X_1} \left| \widehat f_2(x_1, \xi')\right|^2 dx_1 \\
+ C \sum_{j=1}^d \int_{X_0}^{X_1} \left| \hat F_j(x_1, \xi')\right|^2 dx_1.
\end{multline*}
Integrating in $\xi' \in \R^{d-1}$ and using Parseval's identity, we derive~\eqref{Carleman-L2}.

To prove~\eqref{Est-R-tau-L2}, we simply remark that $\widehat {R_\tau w}\! = \!\widehat {K_{\tau,0} f_2}$ with $\widehat f_2 (x_1, \xi')\! =\! \partial_1 \psi(x_1, \xi')\linebreak \widehat w(x_1, \xi')$, which clearly satisfies $\| f_2 \|_{L^2(\Omega)} \leq C \| \nabla' w\|_{L^2(\Omega)}$. Accordingly the estimates on $K_{\tau,0}$ in~\eqref{Carleman-L2} immediately provide~\eqref{Est-R-tau-L2}.
\end{proof}

In view of the above computations, for $w$ as in~\eqref{parametrix-f2_*=0}, we have good estimates on $\partial_1 \hat w$ in $L^2(\Omega_{1, \tau})$, where
\begin{equation}\label{Def-Omega-1-tau}
\Omega_{1,\tau} = \{ (x_1, \xi') \in \Omega,\text{ with } \psi(x_1, \xi') > \tau \}
\cup \{ (x_1, \xi') \in \Omega,\text{ with } \psi(x_1, \xi') < \tau \}.
\end{equation}
Indeed, from the above computations, for $w$ as in~\eqref{parametrix-f2_*=0}, we have
\begin{multline}\label{Est-d1-Hilbertian-Est}
\tau \| \partial_1 \hat w \|_{L^2(\Omega_{1, \tau})}\\
\leq C \| f_2 \|_{L^2(\Omega)}
+ C \tau \sum_{j = 1}^d \left\| F_j\right\|_{L^2(\Omega)}
+ C \tau \| g_0 \|_{H^{\frac{1}{2}}\left(\R^{d-1}\right)} + C\tau \| g_1\|_{H^{-\frac{1}{2}}\left(\R^{d-1}\right)}.
\end{multline}
In particular, if one knows that $w \in H^1(\Omega)$, we get an estimate on $\partial_1 w$ in $L^2(\Omega)$.

Note that the above proof and the previous remark immediately give the following result, whose proof is left to the reader, since all solutions $w$ of~\eqref{Elliptic-Strip-w} with source terms $(f_2, f_{2_*'},F,g)$ as in~\eqref{Reg-Source-Terms-Strip} belong to $H^1(\Omega)$, and the terms $\| \nabla' w\|_{L^2(\Omega)}^2$ coming from $R_\tau(w)$ and~\eqref{Est-R-tau-L2} can be easily absorbed by taking $\tau$ large enough:

\begin{prop}\label{Prop-parametrix-Strip-2}
Let $\Omega$ be as in~\eqref{Omega=Strip} with $X_0 \! <\! 0\! <\! X_1$ and $\max\{|X_0|, |X_1|\}\linebreak \leq 1$, and assume that the coefficients $(\lambda_j)_{j\, \in\,\{1,\,\cdots,\,d\}} \in \R^d$ satisfy $\lambda_1 = 0$, \eqref{Coercivity} and~\eqref{Sub-ellipticity-Strip}.

Then there exist constants $C>0$ and $\tau_0 \geq 1$ depending on $c_0$, $m_*$ and $M_*$ (independent of $X_0, X_1$) such that for all $(f_2, f_{2_*'},F,g)$ as in~\eqref{Reg-Source-Terms-Strip} with $f_{2_*'} = 0$, if the solution $w$ of~\eqref{Elliptic-Strip-w} satisfies $(\partial_1 w - F_1)(X_1,x') = 0$ in $\R^{d-1}$, then for all $\tau \geq \tau_0$,
\begin{equation*}
w = K_{\tau,0} (f_2) + K_{\tau,1}(F_1) + \sum_{j = 2}^d K_{\tau,j}(F_j) + R_\tau (w) + G_\tau(g),
\end{equation*}
where the operators $(K_{\tau,i})_{i\, \in\,\{0,\,\cdots,\,d\}}$, $R_\tau$ and $G_\tau$ are defined in Proposition~\ref{Prop-parametrix-Strip-1}, and
\begin{multline*}
\tau^{\frac{3}{2}} \| w \|_{L^2(\Omega)} + \tau \| \partial_1 w \|_{L^2(\Omega)} + \tau^{\frac{1}{2}} \left\| \nabla' w \right\|_{L^2(\Omega)}\\
\leq  C \| f_2 \|_{L^2(\Omega)} + C \tau \| F\|_{L^2(\Omega)} + C \tau \| g \|_{H^{\frac{1}{2}}\left(\R^{d-1}\right)}.
\end{multline*}
\end{prop}


We now check that the estimate on $G_{\tau}$ can indeed be improved:
\begin{theo}\label{Thm-Improved-Hilbert-Est-G-tau}
Let $\Omega$ be as in~\eqref{Omega=Strip} with $X_0 \! <\! 0\! <\! X_1$ and $\max\{|X_0|, |X_1|\} \leq 1$, and assume that the coefficients $(\lambda_j)_{j\,\in\,\{1,\,\cdots,\,d\}} \in \R^d$ satisfy $\lambda_1 = 0$, \eqref{Coercivity} and~\eqref{Sub-ellipticity-Strip}, and let $G_\tau$ be the operator in~\eqref{Def-G-tau-Fourier} and~\eqref{Def-g-tau-x1-xi'}.

Then there exists a constant $C>0$ depending on $c_0$, $m_*$ and $M_*$ (independent of $X_0, X_1$) such that for all $\tau \geq 1$, for all $g \in H^{\frac{1}{2}}(\R^{d-1})$,
\begin{equation*}
\tau^{\frac{3}{4}} \| G_{\tau} (g) \|_{L^2(\Omega)} +
\left\| \widehat{\partial_1 G_\tau} (g) \right\|_{L^2(\Omega_{1, \tau})} +
\tau^{-\frac{1}{4}} \left\| \nabla' G_\tau (g) \right\|_{L^2(\Omega)}
\leq C \| g \|_{H^{\frac{1}{2}}\left(\R^{d-1}\right)},
\end{equation*}
where $\Omega_{1, \tau}$ is defined by~\eqref{Def-Omega-1-tau}.
\end{theo}
\begin{rema}
Note that the estimates in Theorem~\ref{Thm-Improved-Hilbert-Est-G-tau} yield better estimates than the ones of Theorem~\ref{Thm-Hilbertian-est} on the $\scL(H^{\frac{1}{2}}(\R^{d-1}), L^2(\R^d))$ norms of the operators $G_\tau$ and $\nabla' G_{\tau}$, and are in agreement with the ones obtained in~\cite{imanuvilov2003global}.
\end{rema}

\begin{proof}
For $\tau \geq 1$ and $\xi' \in \R^{d-1}$ such that $\psi(X_0, \xi') > \tau$, we introduce $x_1^*(\xi') \in (X_0, X_1]$ as the unique solution of $\psi(x_1^*(\xi'), \xi') = \tau$ if it exists, or $x_1^*(\xi') = X_1$ otherwise, and we compute
\pagebreak
\begin{align*}
\int_{X_0}^{x_1^*(\xi')} \left|g_\tau(x_1, \xi')\right|^2 dx_1,
&\text{ and}
\quad
\int_{X_0}^{x_1^*(\xi')} |\xi'|^2 \left|g_\tau(x_1, \xi')\right|^2 dx_1,\\
&\text{ and} 
\quad
\int_{X_0}^{x_1^*(\xi')} \left|\partial_1 g_\tau(x_1, \xi')\right|^2 dx_1.
\end{align*}
In order to do that, we recall that, within the setting of Theorem~\ref{Thm-Improved-Hilbert-Est-G-tau}, we have
\begin{align*}
&\exists C > 0, \forall \xi' \in  \R^{d-1}, \forall (x_1,y_1)  \in  [ X_0, X_1 ]^2\text{ with } y_1  <  x_1,
\\
&\mkern300mu \psi(y_1, \xi') \geq  \psi(x_1, \xi')  + \frac{1}{C} |\xi'| |x_1  -  y_1|,
\\
&
\exists C>0, \ \forall \xi' \in \R^{d-1},\, \forall \tau \geq 1, \ \forall x_1 \in [X_0, x_1^*(\xi')],\\
&\mkern300mu
\frac{1}{C} |\xi'| |x_1^*(\xi') - x_1| \leq \tau - \psi(x_1, \xi'),
\end{align*}
so that
\begin{multline*}
\int_{X_0}^{x_1^*(\xi')} \left|g_\tau(x_1, \xi')\right|^2 \, dx_1\\
\begin{aligned}
&\leq
\int_{X_0}^{x_1^*(\xi')} e^{2 \tau (x_1 -X_0) - 2 \int_{X_0}^{x_1} \psi(y_1, \xi') \, dy_1 } \, dx_1
\\
& \leq
\int_{X_0}^{x_1^*(\xi')} e^{2 \left(\tau - \psi(x_1, \xi')\right) (x_1 -X_0) - |\xi'| (x_1- X_0)^2/C } \, dx_1
\\
& \leq
\int_{X_0}^{x_1^*(\xi')} e^{- 2 |\xi'| \left(x_1^*(\xi') - x_1\right) (x_1 -X_0)/C - |\xi'| (x_1- X_0)^2/C} \, dx_1
\\
& \leq
\int_{X_0}^{\left(X_0+x_1^*(\xi')\right)/2}e^{- 2|\xi'| \left(x_1^*(\xi') - x_1\right) (x_1 -X_0)/C } \, dx_1 \\
&\mkern250mu+
\int_{\left(X_0+x_1^*(\xi')\right)/2}^{x_1^*(\xi')} e^{ - |\xi'| (x_1- X_0)^2/C} \, dx_1
\\
& \leq \!\!
\int_{X_0}^{\left(X_0+x_1^*(\xi')\right)/2}\!\!e^{- |\xi'| \left(x_1^*(\xi') - X_0\right) (x_1 -X_0)/(2C)} dx_1 
\\
&\mkern250mu
 + 
\int_{\left(X_0+x_1^*(\xi')\right)/2}^{x_1^*(\xi')} \!\!e^{- |\xi'| \left(x_1^*(\xi') - X_0\right) (x_1- X_0)/C} dx_1
\\
& \leq 2 \int_{X_0}^{\left(X_0+x_1^*(\xi')\right)/2}e^{- |\xi'| \left(x_1^*(\xi') - X_0\right) (x_1 -X_0)/(2C)} \, dx_1
\\
& \leq C \min\left\{\frac{1}{|\xi'| |x_1^*(\xi')- X_0|}, \left|x_1^*(\xi')- X_0\right| \right\}.
\end{aligned}
\end{multline*}
It is then easy to check that, for $\tau \geq 1$ and $\xi' \in \R^{d-1}$ such that $\psi(X_0, \xi') > \tau$,
\begin{multline*}
\int_{X_0}^{x_1^*(\xi')} \left|g_\tau(x_1, \xi')\right|^2 \, dx_1 |g(\xi')|^2
\\
\begin{aligned}
&\leq C
\left(\!\sup_{\psi(X_0, \xi')\, >\,\tau}
\left\{
\min\left\{\frac{1}{|\xi'|^2 |x_1^*(\xi')- X_0|}, \frac{|x_1^*(\xi')\!-\! X_0|}{|\xi'|} \right\}
\right\} \right) |\xi'| |g(\xi')|^2\\
&\leq
\frac{C}{\tau^{\frac{3}{2}}} \left|\xi'\right|\; |g(\xi')|^2.
\end{aligned}
\end{multline*}

Integrating in $\xi' \in \R^{d-1}$ and using Parseval's identity, we obtain that
\begin{equation*}
\tau^{\frac{3}{4}} \| G_{\tau} (g) \|_{L^2(\Omega)}
\leq C \| g \|_{H^{\frac{1}{2}}\left(\R^{d-1}\right)}.
\end{equation*}
Similarly, we prove that
\begin{multline*}
\int_{X_0}^{x_1^*(\xi')} \left|\xi'\right|^2 \left|g_\tau(x_1, \xi')\right|^2 \, dx_1 \left|g(\xi')\right|^2
\\
\begin{aligned}
&\leq  C 
\left( \sup_{\psi(X_0, \xi')\, >\,\tau}
\left\{
\min\left\{\frac{1}{\left|x_1^*(\xi')
-  X_0\right|}, |\xi'| \left|x_1^*(\xi')- X_0\right| \right\}
\right\} \!\right) \left|\xi'\right|\; |g(\xi')|^2\\
&\leq   C \tau^{\frac{1}{2}} \left|\xi'\right|\; \left|g(\xi')\right|^2,
\end{aligned}
\end{multline*}
so that there exists a constant $C>0$ such that for all $g \in H^{\frac{1}{2}}(\R^{d-1})$,
\begin{equation*}
\tau^{-\frac{1}{4}} \left\|\nabla' G_{\tau} (g) \right\|_{L^2(\Omega)}
\leq C \| g \|_{H^{\frac{1}{2}}\left(\R^{d-1}\right)}.
\end{equation*}
We then check that
\[
\partial_1 g_\tau(x_1, \xi') = \left(\tau - \psi(x_1, \xi')\right) g_{\tau}\left(x_1, \xi'\right).
\]
Then, if $\xi' \in \R^{d-1}$ is such that $\psi(X_0, \xi') > \tau$ and $\psi(X_1, \xi') < \tau$, using that for all $x_1 \in [X_0, x_1^*(\xi'))$, $\tau - \psi(x_1, \xi') \leq C |\xi'| |x_1^*(\xi') - x_1|$, we get
\begin{align*}
& \int_{X_0}^{x_1^*(\xi')} \left|\partial_1 g_\tau(x_1, \xi')\right|^2 \, dx_1
\leq
\int_{X_0}^{x_1^*(\xi')} \left(\tau - \psi(x_1, \xi')\right)^2 e^{2 \tau (x_1 -X_0) - 2 \int_{X_0}^{x_1} \psi(y_1, \xi') \, dy_1 } \, dx_1
\\
& \leq C \int_{X_0}^{x_1^*(\xi')} \left|\xi'\right|^2 \left(x_1^*(\xi') - x_1\right)^2 e^{- 2 \left|\xi'\right| \left(x_1^*(\xi') - x_1\right) (x_1 -X_0)/C - |\xi'| (x_1- X_0)^2/C} \, dx_1
\\
& \leq C \int_{X_0}^{\left(X_0+x_1^*(\xi')\right)/2} |\xi'|^2 \left(x_1^*(\xi') - X_0\right)^2 e^{- 2 |\xi'|\left(x_1^*(\xi') - X_0\right) (x_1 -X_0)/C} \, dx_1
\\
& \qquad + C \int_{\left(X_0+x_1^*(\xi')\right)/2}^{x_1^*(\xi')} |\xi'|^2 \left(x_1^*(\xi') - x_1\right)^2 e^{- |\xi'| \left(x_1^*(\xi') - X_0\right) (x_1- X_0)/(2C)} \, dx_1
\\
& \leq C \int_{X_0}^{\left(X_0+x_1^*(\xi')\right)/2} |\xi'|^2 \left(x_1^*(\xi') - X_0\right)^2 e^{- |\xi'| \left(x_1^*(\xi') - X_0\right) (x_1 -X_0)/(2C)} \, dx_1
\\
& \leq C \min\left\{ |\xi'| (x_1^*(\xi') - X_0), |\xi'|^2 \left(x_1^*(\xi') - X_0\right)^3 \right\}
\\
&
\leq C \min\left\{ \psi(X_0, \xi') - \tau, \frac{(\psi(X_0,\xi') - \tau)^3}{|\xi'|} \right\}.
\end{align*}
In particular,
\begin{multline*}
\int_{X_0}^{x_1^*(\xi')} \left|\partial_1 g_\tau(x_1, \xi')\right|^2 \, dx_1 |g(\xi')|^2
\\
\begin{aligned}
&\leq C
\left(\!\sup_{\psi(X_0, \xi')\,>\,\tau} \left\{ \min\left\{ \frac{\psi(X_0, \xi') - \tau}{|\xi'|}, \frac{(\psi(X_0,\xi') \!-\! \tau)^3}{|\xi'|^2} \right\}\right\} \!\right) \left|\xi'\right|\; \left|g(\xi')\right|^2\\
&\leq \! C \left|\xi'\right|\; \left|g(\xi')\right|^2.
\end{aligned}
\end{multline*}
Similar estimates can be achieved for $\xi' \in \R^{d-1}$ such that $\psi(X_1, \xi') \geq \tau$, and details are left to the reader.

We thus obtained that there exists a constant $C>0$ such that for all $g \in H^{\frac{1}{2}}(\R^{d-1})$ and $\tau \geq 1$,
\begin{equation*}
\left\|\partial_1 \widehat{ G_{\tau} (g)} \right\|_{L^2(\Omega_{1,\tau})}
\leq C \| g \|_{H^{\frac{1}{2}}\left(\R^{d-1}\right)}.
\end{equation*}
This concludes the proof of Theorem~\ref{Thm-Improved-Hilbert-Est-G-tau}.
\end{proof}

\section{Fourier Restriction theorems and applications}\label{Sec-Fourier-Restriction-Thm}

In this section, we first recall the classical Fourier restriction theorem, and present a version adapted to our case. We will then explain how it can be applied to estimate the norms of operators of some specific forms, which will encompass the ones appearing in the parametrix provided in Proposition~\ref{Prop-parametrix-Strip-1}.


In this section, $n \geq 2$ and, for a function $f \in \scS(\R^n)$, the Fourier transform $\widehat f \in \scS(\R^n)$ is given by
\begin{equation*}
\widehat f(\xi) = \frac{1}{(2\pi)^{n/2}} \int_{x\, \in\,\R^n} e^{- \i x \cdot \xi} f(x) \, dx, \qquad \qquad\xi \in \R^n,
\end{equation*}
and is extended by duality to functions in $\scS'(\R^n)$ as usual. We will see later that $n$ in fact corresponds to $d-1$ in the applications we have in mind.

\subsection{Fourier restriction theorems}

We start by recalling the classical Stein--Tomas Fourier restriction theorem:
\begin{theo}
[{\cite{tomas1975restriction}, see also \cite[Theorem~2, p.~352]{stein1993harmonic}}]\label{Thm-Stein-Tomas}
Let $n\geqslant 2$, and $\S^{n-1}$ denote the unit sphere of $\R^n$.

Then the map
\begin{equation*}
\begin{cases}
L^1(\R^n) \to L^2\left(\S^{n-1}\right)
\\
f \mapsto \hat f\, \big|_{\S^{n-1}}
\end{cases}
\end{equation*}
can be extended by continuity on $L^{\frac{2(n+1)}{(n+3)}}(\R^n)$, and there exists a constant $C>0$ such that for all $f \in L^{\frac{2(n+1)}{(n+3)}}(\R^n)$,
\begin{equation*}
\left\| \widehat f \right\|_{L^2\left(\S^{n-1}\right)} \leqslant C \|f\|_{L^{\frac{2(n+1)}{(n+3)}}\left(\R^n\right)}.
\end{equation*}
\end{theo}

It is well-known that this restriction theorem can be extended to any hypersurface with non-vanishing Gaussian curvature, (see, for example, \cite[Corollary~2.2.2]{sogge2017fourier}).

In view of the formulae in Proposition~\ref{Prop-parametrix-Strip-1}, it is interesting for us to analyze Fourier restriction theorems on the family of surfaces
\begin{equation}\label{Def-Sigma-a}
\Sigma_{a} = \left\{ \xi \in \R^n, \, \psi(a, \xi) = 1\right\}, \qquad \qquad a \in [X_0, X_1],
\end{equation}
where by analogy with the function $\psi$ in~\eqref{Def-Psi} and the conditions~\eqref{Coercivity}, we have set
\begin{equation}\label{Def-Psi-R-n}
\psi(a, \xi)=\sqrt{ \sum_{j=1}^{n}(1-a \lambda_{j}) \xi_{j} ^2},
\qquad \qquad a \in [X_0,X_1], \ \xi \in \R^{n}.
\end{equation}
where the family of coefficients $(\lambda_j)_{j\, \in\,\{1,\,\cdots,\,n\}}$ satisfies
\begin{equation}\label{Coercivity-Restriction}
\exists c_0 >0, \quad \forall a \in [X_0, X_1], \, \forall \xi \in \R^{n}, \quad
\frac{1}{c_0} | \xi|^2
\leq
\sum_{j =1}^n \left(1 - a \lambda_j\right) \left|\xi_j\right|^2
\leq c_0 |\xi|^2.
\end{equation}

Note that due to condition~\eqref{Coercivity-Restriction}, for all $a\in [X_0, X_1]$, the surface $\Sigma_{a}$ is an ellipsoid and thus~\cite[Corollary~2.2.2]{sogge2017fourier} applies and yields that for all $a \in [X_0, X_1]$, the map $f \mapsto \hat f\,|_{\Sigma_{a}} $ maps $L^{\frac{2(n+1)}{(n+3)}} (\R^n)$ to $L^2(\Sigma_{a})$.

For our purpose, we need a slightly more refined version of this result, guaranteeing that the norm of this map is independent of $a \in [X_0, X_1]$.
\begin{theo}\label{Thm-Fourier-Rest-Sigma-x-1}
Let $n \geq 2$. Assume that the family of coefficients $(\lambda_j)_{j\, \in\,\{1,\, \cdots,\,n\}}$ satisfies~\eqref{Coercivity-Restriction} for some $c_0 >0$. Then there exists a constant $C>0$ depending only on $c_0$ (and $n$) such that for all $a \in [X_0, X_1]$, for all $f \in L^{\frac{2(n+1)}{(n+3)}}(\R^n)$,
\begin{equation}\label{Uniform-Restriction-Thm}
\left\| \widehat f\right \|_{L^2(\Sigma_{a})} \leqslant C \|f\|_{L^{\frac{2(n+1)}{(n+3)}}\left(\R^n\right)}.
\end{equation}
\end{theo}

Note that the proof below follows the classical one of Theorem~\ref{Thm-Stein-Tomas} and is mainly based on the stationary phase lemma.

\begin{proof}
For $a \in [X_0, X_1]$, we denote by $T_{a}$ the map $T_{a} : f \in L^{1}(\R^n) \mapsto \widehat f \,|_{\Sigma_{a}} \in L^2(\Sigma_{a})$. We then consider its adjoint operator $T_{a}^\star : L^2(\Sigma_{a}) \to L^\infty(\R^n)$: for $g \in L^2(\Sigma_{a})$,
\begin{equation*}
T^\star_{a} g(x)
\\
=
\int_{{\omega}\, \in\,\Sigma_{a} } e^{\i x \cdot {\omega}} g({\omega}) \, d \Sigma_{a}({\omega}),
\qquad \qquad x\in \R^{n}.
\end{equation*}
The operator $T^\star_{a}T_{a}$ then maps $L^1(\R^n)$ to $L^\infty(\R^n)$ and, for $f \in L^1(\R^n)$,
\begin{equation}\label{Def-T-star-T}
T^\star_{a} T_{a} f(x) =
\int_{\R^{n}}
\int_{{\omega}\, \in\,\Sigma_{a} } e^{\i (x-\tilde{x}) \cdot {\omega}}
\, d \Sigma_{a}({\omega}) f(\tilde{x})
\, d \tilde{x},
\qquad \qquad x \in \R^n.
\end{equation}

Next, we will prove that the operator $ T^\star _{a}T_{a}$ can in fact be extended as an operator from $L^{\frac{2(n+1)}{(n+3)}}(\R^n)$ to $L^{\frac{2(n+1)}{(n-1)}}(\R^n)$, uniformly with respect to $a \in [X_0, X_1]$. This will prove~\eqref{Uniform-Restriction-Thm} since
\[
\left\| T^\star _{a}T_{a}\right\|_{\scL\left(L^{\frac{2(n+1)}{(n+3)}}(\R^n),\, L^{\frac{2(n+1)}{(n-1)}}(\R^n)\right)} =
\| T_a \|_{\scL\left(L^{\frac{2(n+1)}{(n+3)}}(\R^n), L^{2}(\Sigma_a)\right)}^2.
\]

To start with, we parametrize the hypersurface $\Sigma_{a}$ through several patches. We first remark that $\Sigma_a$ can be mapped into the sphere $\S^{n-1}$ as follows. For $\omega \in \Sigma_a$, we define $\xi = G_a(\omega)$ by
\[
\forall j \in \{1,\,\cdots,\,n\}, \quad \xi_j = \omega_j \sqrt{1 - a \lambda_j}.
\]
We then choose a spherical cap
\[
\scC_n =
\left\{
\xi \in \S^{n-1}; \quad \xi_n \geq \frac{1}{\sqrt{2n}}
\right\}.
\]
It is easy to check that, if, for $\epsilon \in \{-1,1\}$ and $j \in \{1,\,\cdots,\,n\}$, $R_{\epsilon, j}$ denotes the rotation that maps the basis vector $e_n$ to $\epsilon e_j$, and leaves all the vectors $e_k$ for $k \neq j,n$ invariant, then the family of $R_{\epsilon,j} \scC_n$ for $\epsilon \in \{-1,1\}$ and $j \in \{1,\cdots, n\} $ covers the whole sphere. Therefore, there exists a partition of unity $(\chi_{\epsilon, j})_{\epsilon\, \in\,\{-1,1\},\,j\,\in\,\{1,\,\cdots,\, n\}}$ of the sphere $\S^{n-1}$ such that for each $\epsilon \in \{-1,1\}$ and $j \in \{1,\,\cdots,\,n\} $, the function $\chi_{\epsilon, j}$ is smooth and compactly supported in $R_{\epsilon,j} \scC_n$. Since by construction, $\sum_{\epsilon,j} \chi_{\epsilon,j} (\xi) = 1$ for all $\xi \in \S^{n-1}$, we have
\[
\forall \omega \in \Sigma_a,
\quad
\sum_{\epsilon\,\in \{-1,1\}, \, j\,\in\,\{1,\, \cdots,\,n\} } \chi_{\epsilon,j} (G_a(\omega)) = 1.
\]
Therefore,
\begin{equation*}
T^\star _{a}T_{a} f(x) =\mkern-25mu
\sum_{\epsilon\,\in\,\{-1,1\}, \ j\,\in\,\{1,\, \cdots,\,n\} }
\int_{\R^{n}}
\int_{ {\omega}\,\in\,\Sigma_{a} }
\chi_{\epsilon,j} (G_a(\omega)) e^{\i \left(x-\tilde{x}\right) \cdot {\omega}}
\, d \Sigma_{a}({\omega}) f(\tilde{x})
\, d \tilde{x},
\quad x \in \R^n.
\end{equation*}
Besides, for all $\epsilon \in \{-1,1\}$ and $j\in \{1,\,\cdots,\,n\}$, $\chi_{\epsilon,j}\circ G_a$ is supported in the set of all $\omega \in \Sigma_{a}$ such that $\epsilon \omega_{j}\sqrt{1 - a \lambda_j} \geq 1/\sqrt{2n}$, i.e. the pre-image of the cap $R_{\epsilon,j} (\scC_n)$ by $G_a$, that we denote by $\scC_{\epsilon,j}$. It is clear that this set can be parametrized by $\R^{n-1}$ as follows. Denoting $\check\xi_j = (\xi_1, \cdots, \xi_{j-1}, \xi_{j+1},\,\cdots,\,\xi_n)$, we easily obtain
\begin{equation*}
\scC_{\epsilon,j} =
\left\{
\left(\xi_1,\,..,\,\xi_{j-1}, \epsilon h_{j}\left(a, \check\xi_j\right), \xi_{j+1},\,..,\,\xi_n \right);
\quad
\sum_{\substack{k=1 \\ k \neq j}}^{n} (1-a \lambda_{k}) \xi_k^2 \leq 1- \frac{1}{2n}
\right\},
\end{equation*}
where the function $h_{j}$ is defined by the formula
\begin{equation*}
h_{j}\left(a, \check \xi_j\right) =
\frac{1}{\sqrt{1-a \lambda_j}} \sqrt{ 1- \sum_{\substack{k=1 \\ k \neq j}}^{n} (1-a \lambda_k) \xi_k^ 2},
\qquad a \in [X_0,X_1],\, \check \xi_j \in V_{a,j},
\end{equation*}
with $V_{a,j}$ given by
\[
V_{a,j} =
\left\{
\check\xi_j \in \R^{n-1}; \sum_{\substack{k=1 \\ k \neq j}}^{n} (1-a \lambda_{k}) \xi_k^2 \leq 1 - \frac{1}{2n}
\right\}.
\]
Therefore, the study of $T_a^\star T_a$ is reduced to the study of the family of operators
\begin{multline*}
f
\mapsto
\left(x \mapsto \int_{\R^{n}}
\int_{ \check \xi_j \in V_{a,j}}
\chi_{\epsilon,j} \left(G_a\left(H_{\epsilon,j}\left(a,\check\xi_j\right)\right)\right) e^{\i \left(\check x_j-\check{y}_j\right)\cdot \check \xi_j + \epsilon \i \left(x_j- \tilde x_j\right) h\left(a,\check\xi_j\right)}
\right.
\\
\left.
%\phantom{\int_{\R^{n}}}
\times\sqrt{1 + \left|\nabla h_j\left(a,\check \xi_j\right)\right|^2} d \check \xi_j f(y)
\, d y
\right),
\end{multline*}
where $H_{\epsilon,j}(a, \check \xi_j) = \left(\xi_1,\,..,\,\xi_{j-1}, \epsilon h_{j}(a, \check\xi_j), \xi_{j+1},\,..,\,\xi_n \right)$, for all $j \in \{1,\,\cdots,\,n\}$ and $\epsilon \in \{-1,1\}$.

Thus, up to a renumbering of the coefficients, we can focus without loss of generality on the operator corresponding to $\epsilon = 1$ and $j = n$. Accordingly, we introduce the notation $x' = (x_1,\,\cdots,\,x_{n-1})$ and $\xi' = (\xi_1,\,\cdots,\,\xi_{n-1})$, and we consider the operator, defined for $f \in L^1(\R^{n})$ to $L^\infty(\R^n)$ by
\begin{multline}\label{Def-T-a}
\clT_a f(x) =
\int_{\R^{n}}
\int_{\xi'\,\in\,\R^{n-1}}
\chi\big(G_a(\xi', h_n(a,\xi'))\big) e^{\i \left(x' - \tilde x'\right)\cdot \xi' + \i \left(x_n- \tilde x_n\right) h_n(a,\xi')}
\\
\times\sqrt{1 + |\nabla h_n(a,\xi')|^2} d \xi' f(\tilde x)
\, d \tilde x,
\end{multline}
where $\chi$ is a smooth function on the sphere $\S^{n-1}$ compactly supported in the spherical cap $\scC_n$ and $\chi(G_a(\cdot, h_n(a, \cdot))$ is extended by $0$ for $\xi' \notin V_{a,n}$.

We have reduced the proof of Theorem~\ref{Thm-Fourier-Rest-Sigma-x-1} to the proof of the fact that the maps $\clT_a$ defined in~\eqref{Def-T-a} belong to the space $\scL(L^{\frac{2(n+1)}{(n+3)}}(\R^n), L^{\frac{2(n+1)}{(n-1)}}(\R^n))$ uniformly with respect to $a \in [X_0, X_1]$.

In order to show this property, for $a \in [X_0, X_1]$ and $\delta \in \R$, we introduce the family of operators, defined from $L^{1}(\R^{n-1})$ to $L^\infty(\R^{n-1})$ by
\begin{multline*}\label{Def-T-a-delta}
\scT_{a,\delta} f(x') =
\int_{\R^{n-1}}
\int_{\xi'\,\in\,\R^{n-1}}
\chi\big(G_a(\xi', h_n(a,\xi'))\big) e^{\i \left(x' - \tilde x'\right)\cdot \xi' + \i \delta h_n(a,\xi')}
\\
\times
\sqrt{1 + |\nabla h_n(a,\xi')|^2} d \xi' f(\tilde x')
\, d \tilde x',
\end{multline*}
for which we will show that there exists a constant $C>0$ such that for all $a \in [X_0, X_1]$ and $\delta \in \R$,
\begin{equation}\label{Est-T-a-delta}
\norm{ \scT_{a, \delta}}_{\scL\left(L^{\frac{2(n+1)}{(n+3)}}\left(\R^{n-1}\right), L^{\frac{2(n+1)}{(n-1)}}\left(\R^{n-1}\right)\right)}
\leqslant C | \delta|^{-\frac{n-1}{n+1}}.
\end{equation}

Indeed, if the estimate~\eqref{Est-T-a-delta} holds, then Hardy--Littlewood--Sobolev theorem (recalled in Appendix in Theorem~\ref{Thm-Hardy-Littlewood-Sobolev}) implies that, for $f \in L^1(\R^{n}) \cap L^{\frac{2(n+1)}{(n+3)}}(\R^n)$,
\begin{align*}
\norm{ \clT_a f }_{L^{\frac{2(n+1)}{(n-1)}}(\R^n)} &
\leq
\norm{ \norm{ \int_{\tilde x_n\, \in\,\R} \scT_{a, x_n - \tilde x_n } f(\cdot, \tilde x_n) \, d\tilde x_n }_{L^{\frac{2(n+1)}{(n-1)}}_{x'}\left(\R^{n-1}\right)}}_{L^{\frac{2(n+1)}{(n-1)}}_{x_n}(\R)}
\\
&
\leq
\norm{\int_{\tilde x_n\,\in\,\R} \norm{\scT_{a, x_n -\tilde x_n } f(\cdot, \tilde x_n) }_{L^{\frac{2(n+1)}{(n-1)}}_{x'}(\R^{n-1})}d\tilde x_n}_{L^{\frac{2(n+1)}{(n-1)}}_{x_n}(\R)}
\\
&
\leq C \norm{\int_{\tilde x_n \in \R} \left| \tilde x_n - x_n\right|^{-\frac{n-1}{n+1}} \norm{f(\cdot, \tilde x_n) }_{L^{\frac{2(n+1)}{(n+3)}}_{\tilde x'}\left(\R^{n-1}\right)} d\tilde x_n}_{L^{\frac{2(n+1)}{(n-1)}}_{x_n}(\R)}
\\
&
\leq C \norm{ \norm{f(\cdot, \tilde x_n) }_{L^{\frac{2(n+1)}{(n+3)}}_{\tilde x'}\left(\R^{n-1}\right)}}_{L^{\frac{2(n+1)}{(n+3)}}_{\tilde x_n}(\R)} = C \norm{f}_{L^{\frac{2(n+1)}{(n+3)}}\left(\R^n\right)},
\end{align*}
since
\[
\frac{n-1}{n+1} = 1 - \left(\frac{n+3}{2(n+1)} - \frac{n-1}{2(n+1)}\right).
\]

We thus focus on the proof of estimate~\eqref{Est-T-a-delta}, which, as explained above, would conclude the proof of Theorem~\ref{Thm-Fourier-Rest-Sigma-x-1}. This is done in three steps.

In the first step, we check that $\scT_{a, \delta}$ maps $L^2(\R^{n-1})$ into itself with uniform bounds. Indeed, taking the Fourier transform $x' \to \xi'$ of $\R^{n-1}$, we easily get:
\[
\widehat {\scT_{a, \delta} f}(\xi') =
\chi\big(G_a(\xi', h_n(a,\xi'))\big) e^{\i \delta h_n(a,\xi')}
\,
\sqrt{1 + \left|\nabla h_n(a,\xi')\right|^2}
\widehat {f} (\xi'),
\]
so that by Parseval's identity,
\[
\norm{
\scT_{a,\delta} }_{\scL\left(L^{2}\left(\R^{n-1}\right)\right)}
\leq
\norm{
\chi\big(G_a(\xi', h_n(a,\xi'))\big) e^{\i \delta h_n(a,\xi')}
\sqrt{1 + \left|\nabla h_n(a,\xi')\right|^2} }_{L^\infty_{\xi'}(\R^{n-1})}.
\]
We then immediately get that there exists a constant $C>0$ depending only on $c_0$ in~\eqref{Coercivity-Restriction} such that for all $a \in [X_0, X_1]$ and $\delta \in \R$,
\begin{equation}\label{Est-T-a-delta-L2-L2}
\norm{
\scT_{a,\delta} }_{\scL\left(L^{2}\left(\R^{n-1}\right)\right)}
\leqslant C.
\end{equation}

In a second step, we check that $\scT_{a, \delta}$ maps continuously $L^1(\R^{n-1})$ to $L^\infty(\R^{n-1})$ and get an estimate on its norm. In fact, we clearly have that
\begin{multline}\label{Est-T-a-delta-L1-Linfini}
\norm{
\scT_{a,\delta} }_{\scL\left(L^{1}\left(\R^{n-1}\right), L^\infty\left(\R^{n-1}\right)\right)} \\
\leq
\sup_{\delta'\, \in\,\R^{n-1}}
\left| \int_{\R^{n-1}} e^{\i \delta' \cdot \xi' + \i \delta h_n(a, \xi')} \chi\big(G_a(\xi', h_n(a,\xi'))\big) \sqrt{1 + |\nabla h_n(a,\xi')|^2} \, d \xi' \right|.
\end{multline}
Our goal is then to prove that there exists a constant $C>0$ such that for all $a \in [X_0, X_1]$ and $\delta \in \R$,
\begin{multline}\label{Stationary-Phase-Goal}
\sup_{\delta'\, \in\,\R^{n-1}}
\left|
\int_{\R^{n-1}} e^{\i \delta' \cdot \xi' + \i \delta h_n(a, \xi')} \chi\big(G_a(\xi', h_n(a,\xi'))\big) \sqrt{1 + \left|\nabla h_n(a,\xi')\right|^2} \, d \xi'
\right|\\
\leq
\frac{C}{|\delta|^{\frac{n-1}{2}}}.
\end{multline}
For $\delta' \in \R^{n-1}$ and $\delta \in \R$, we define $\lambda >0$ and $\omega \in \S^{n-1}$ by
\[
\lambda = \sqrt{|\delta'|^2 + \delta^2}
\quad
\text{ and }
\omega = \frac{1}{\lambda}(\delta', \delta).
\]
Accordingly,
\begin{multline*}
\int_{\R^{n-1}} e^{\i \delta' \cdot \xi' + \i \delta h_n(a, \xi')} \chi\big(G_a(\xi', h_n(a,\xi'))\big) \sqrt{1 + |\nabla h_n(a,\xi')|^2} d \xi'\\
= 
\int_{\R^{n-1}}  e^{\i \lambda \Phi(\omega, a, \xi')} \tilde \chi(a,\xi') d \xi',
\end{multline*}
where
\begin{equation*}
\Phi(\omega, a, \xi')\! = \!\omega' \cdot \xi' + \omega_n h_n(a, \xi'),
\text{ and} 
\quad
\tilde \chi(a,\xi')\! =\! \chi\big(G_a\left(\xi', h_n(a,\xi')\right)\big) \sqrt{1 \!+\! |\nabla h_n(a,\xi')|^2}.
\end{equation*}


Note that we immediately have that
\begin{equation}\label{Stationary-Phase-Goal-0}
\left| \int_{\R^{n-1}} e^{\i \lambda \Phi(\omega, a, \xi')} \tilde \chi(a,\xi') \, d \xi'\right|
\leq \sup_{a\, \in\,[X_0,X_1]} \left\|\tilde \chi(a, \cdot) \right\|_{L^1\left(\R^{n-1}\right)},
\end{equation}
so that we are only interested in large values of $\delta$. It is then clear that we have to use the stationary phase lemma to get a suitable estimate on that quantity. Since we need to quantify properly in terms of the parameters $a \in [X_0,X_1]$ and $\delta \in \R$, we will use the refined version of {\cite[Theorem~1]{alazard2017stationary}}, recalled in the appendix in Theorem~\ref{Thm-Stationary-Phase}.



Let $\Omega_n \in (0,\frac{1}{2})$ be such that
\[
\Omega_n \sup_{a\, \in\,[X_0, X_1]}
\left\| \nabla_{\xi'} h_n(a,\xi')\right\|_{L^\infty_{\xi'}(V_{a,n})}
\leq
\frac{1}{2} \sqrt{1 - \Omega_n^2}.
\]
Then
\[
\forall \omega \in \S^{n-1} \text{ with } |\omega_n| \leq \Omega_n,
\quad
\inf_{a\, \in\,[X_0, X_1], \, \xi'\,\in \,V_{a,n}} \left|\nabla_{\xi'} \Phi(\omega,a,\xi')\right| \geq \frac{|\omega'|}{2}.
\]
Therefore, if $\omega \in \S^{n-1}$ with $|\omega_n| \leq \Omega_n$, applying integration by parts based on the formula
\[
e^{\i \lambda \Phi(\omega, a, \xi')} =
\frac{1}{\i \lambda |\nabla_{\xi'} \Phi(\omega,a,\xi')|^2} \nabla_{\xi'} \Phi\left(\omega,a,\xi'\right) \cdot \nabla_{\xi'} e^{\i \lambda \Phi(\omega, a, \xi')},
\]
we get that, for all $k \in \N$, there exists $C_k>0$ and a decreasing function $\clF_k: \R \to \R$ such that
\begin{multline*}
\!\!\left|
\int_{\R^{n-1}} e^{\i \lambda \Phi(\omega, a, \xi')} \tilde \chi(a,\xi') \, d \xi'
\right|
\\
\leq 
\frac{C_k}{|\lambda|^k } \clF_k\left(\| \Phi (\omega, \cdot, \cdot) \|_{W^{k+1,\infty}\left([X_0,X_1] \times V_{a,n}\right)}\!\right) \|\tilde \chi \|_{W^{k+1,\infty}\left([X_0,X_1] \times V_{a,n}\right)}.
\end{multline*}
Therefore, for all $k \in \N$, there exists a constant $C_k>0$ such that for $\omega \in \S^{n-1} \text{ with } |\omega_n| \leq \Omega_n$ and $a \in [X_0, X_1]$,
\begin{equation}\label{Stationary-Phase-Goal-1}
\left|
\int_{\R^{n-1}} e^{\i \lambda \Phi\left(\omega, a, \xi'\right)} \tilde \chi(a,\xi') \, d \xi'
\right|
\leq
\frac{C_k}{|\delta|^k}.
\end{equation}
It remains to analyze what happens when $\omega \in \S^{n-1}$ satisfies $|\omega_n| \geq \Omega_n$. There, we use that
\[
\left(\Hess_{\xi'} \Phi\right)(\omega, a, \xi') = \omega_n \left(\Hess_{\xi'} h_n\right)(a,\xi').
\]
Since there exists $a_0 >0$ such that
\[
\inf_{a\, \in\,[X_0,X_1]}\; \inf_{\xi'\,\in\,V_{a,n}} \left|\text{det} \left(\Hess_{\xi'} h_n\right)(a, \xi') \right| \geq a_0.
\]
a direct application of Theorem~\ref{Thm-Stationary-Phase} yields the existence of a constant $C>0$ such that for all $a \in [X_0, X_1]$, for all $\omega \in \S^{n-1}$ with $|\omega_n | \geq \Omega_n$,
\begin{equation}\label{Stationary-Phase-Goal-2}
\left|
\int_{\R^{n-1}} e^{\i \lambda \Phi\left(\omega, a, \xi'\right)} \tilde \chi(a,\xi') \, d \xi'
\right|
\leq
\frac{C}{|\lambda|^{\frac{n-1}{2}}}
\leq
\frac{C}{|\delta|^{\frac{n-1}{2}}}.
\end{equation}
Combining~\eqref{Stationary-Phase-Goal-0}--\eqref{Stationary-Phase-Goal-1}--\eqref{Stationary-Phase-Goal-2}, we get~\eqref{Stationary-Phase-Goal}, and thus from~\eqref{Est-T-a-delta-L1-Linfini}, the existence of a constant $C$ such that for all $a \in [X_0, X_1]$ and $\delta \in \R$,
\begin{equation}\label{Est-T-a-delta-L1-Linfini-Final}
\norm{
\scT_{a,\delta} }_{\scL\left(L^{1}\left(\R^{n-1}\right), L^\infty\left(\R^{n-1}\right)\right)}
\leq
\frac{C}{|\delta|^{\frac{n-1}{2}}}.
\end{equation}

In a third and last step, we conclude the estimate~\eqref{Est-T-a-delta} by M.~Riesz interpolation theorem (\cite[Theorem~0.1.13]{sogge2017fourier}) combining~\eqref{Est-T-a-delta-L2-L2} and~\eqref{Est-T-a-delta-L1-Linfini-Final}. This concludes the proof of Theorem~\ref{Thm-Fourier-Rest-Sigma-x-1}.
\end{proof}


\subsection{Fourier multiplier operators}\label{Subsec-Fourier-Multipliers}

The goal of this section is to show how Theorem~\ref{Thm-Fourier-Rest-Sigma-x-1} can be applied to get estimates on some families of Fourier multipliers operators.

To be more precise, for $X_0 <X_1$ and coefficients $(\lambda_j)_{j\, \in\,\{1,\, \cdots,\,n\}}$ satisfying~\eqref{Coercivity-Restriction}, we define $\psi$ as in~\eqref{Def-Psi-R-n} and $\Sigma_a$ the ellipsoid defined for $a \in [X_0,X_1]$ by~\eqref{Def-Sigma-a}.

For $a \in [X_0, X_1]$ and $k \in L^\infty(\R_+, L^\infty(\Sigma_a))$, we consider operators given as follows:
\begin{multline}\label{Class-Fourier-Multipliers}
K_{a,k} : L^2(\R^n) \to L^2(\R^n),\\
\text{ given by }
\widehat{K_{a,k}(f)} (\xi) = k\left(\psi(a,\xi),\frac{\xi}{\psi(a,\xi)} \right)
\hat f(\xi),
\quad \xi \in \R^n.
\end{multline}


We prove the following result:
\begin{prop}\label{Prop-Est-Fourier-Op}
Let $n\in \N$, $n\geq 2$. Let $X_0 <X_1$, and the coefficients $(\lambda_j)_{j\, \in\,\{1, \,\cdots,\,n\}}$ satisfy~\eqref{Coercivity-Restriction}. For $a \in [X_0,X_1]$, let $\psi$ and $\Sigma_a$ be as in~\eqref{Def-Sigma-a}--\eqref{Def-Psi-R-n}. Then there exists a constant $C>0$ such that, for all $a \in [X_0,X_1]$, for all $k \in L^\infty(\R_+, L^\infty(\Sigma_a))$,
\begin{itemize}
\item the Fourier multiplier operator $K_{a,k}$ in~\eqref{Class-Fourier-Multipliers} maps $L^2(\R^{n})$ to itself and
\begin{equation}\label{Est-K-a-k-L2-L2}
\norm{K_{a,k}}_{\scL\left(L^2(\R^n)\right)}
\leq
\| k \|_{L^\infty\left(\R_+, L^\infty(\Sigma_a)\right)}.
\end{equation}
\item if moreover, $k$ satisfies
\[
\int_0^\infty \| k(\lambda, \cdot)\|_{L^\infty(\Sigma_a)} \, \lambda^{\frac{n-1}{n+1}} \, d\lambda < \infty,
\]
the operator $K_{a,k}$ in~\eqref{Class-Fourier-Multipliers} belongs to $\scL(L^{\frac{2(n+1)}{(n+3)}}(\R^n),L^{\frac{2(n+1)}{(n-1)}}(\R^n))$ and
\begin{equation}\label{Est-K-a-k-Lp-Lp'}
\norm{K_{a,k}}_{\scL\left(L^{\frac{2(n+1)}{(n+3)}}\left(\R^n\right),L^{\frac{2(n+1)}{(n-1)}}\left(\R^n\right)\right)}
\leq C \int_0^\infty \| k(\lambda, \cdot)\|_{L^\infty(\Sigma_a)} \, \lambda^{\frac{n-1}{n+1}} \, d\lambda.
\end{equation}
\item if moreover, $k$ satisfies
\[
\int_0^\infty \| k(\lambda, \cdot)\|_{L^\infty(\Sigma_a)}^2 \, \lambda^{\frac{n-1}{n+1}} \, d\lambda < \infty,
\]
the operator $K_{a,k}$ in~\eqref{Class-Fourier-Multipliers} belongs to
\[
\scL\left(L^{\frac{2(n+1)}{(n+3)}}\left(\R^n\right),L^{2}\left(\R^n\right)\right) \cap \scL\left(L^{2}\left(\R^n\right), L^{\frac{2(n+1)}{(n-1)}}\left(\R^n\right)\right),
\]
and
\begin{align}
\norm{K_{a,k}}_{\scL\left(L^{\frac{2(n+1)}{(n+3)}}\left(\R^n\right),L^{2}\left(\R^n\right)\right)}
\leq C \sqrt{\int_0^\infty \| k(\lambda, \cdot)\|_{L^\infty(\Sigma_a)}^2 \, \lambda^{\frac{n-1}{n+1}} \, d\lambda},
\label{Est-K-a-k-Lp-L2}
\\
\norm{K_{a,k}}_{\scL\left(L^{2}\left(\R^n\right), L^{\frac{2(n+1)}{(n-1)}}\left(\R^n\right)\right)}
\leq C \sqrt{\int_0^\infty \| k(\lambda, \cdot)\|_{L^\infty(\Sigma_a)}^2 \, \lambda^{\frac{n-1}{n+1}} \, d\lambda}.\label{Est-K-a-k-L2-Lp'}
\end{align}
\end{itemize}
\end{prop}

\begin{rema}
The estimates of Proposition~\ref{Prop-Est-Fourier-Op} will play a similar role as Sogge's spectral projection bounds~\cite{sogge1988concerning} for the spherical Laplacian to prove $L^p$ Carleman estimates in the elliptic case as in~\cite{jerison1986carleman, Sogge.Strong.Elliptic1990,koch2001carleman}.
\end{rema}

\begin{proof}
The proof of~\eqref{Est-K-a-k-L2-L2} follows immediately from Parseval's identity.

The proof of~\eqref{Est-K-a-k-Lp-Lp'} is more subtle and is done in several steps. First, based on Theorem~\ref{Thm-Fourier-Rest-Sigma-x-1}, we analyze, for $\lambda>0$, the map $T_{a,\lambda}: f \in L^1(\R^n) \mapsto \hat f|_{\Sigma_{a,\lambda}}$, where $\Sigma_{a,\lambda} = \{ \xi \in \R^n, \, \psi(a, \xi) = \lambda\}$. We then explain how this yields estimate~\eqref{Est-K-a-k-Lp-Lp'}.

The first step is based on the fact that for $a \in [X_0, X_1]$ and $\lambda >0$, for $f \in L^1(\R^n)$,
\begin{equation*}
T_{a,\lambda}^\star T_{a, \lambda} f(x) =
\int_{\tilde x\, \in\,\R^{n}}\;
\int_{\xi\,\in\,\Sigma_{a,\lambda}} e^{\i (x-\tilde x). \xi } \, d \Sigma_{a,\lambda }(\xi) f(\tilde x) \, d \tilde x,
\qquad \qquad x \in \R^n.
\end{equation*}
Since the function $\xi \mapsto \psi(a, \xi)$ is homogeneous of degree~$1$, by using a scaling argument, we get
\begin{align*}
T_{a,\lambda}^\star T_{a, \lambda} f(x) & =
\lambda^{n-1} \int_{\tilde x\, \in\,\R^{n}} \int_{\omega\,\in\,\Sigma_{a,1}} e^{\i \lambda \left(x- \tilde x\right) \cdot \omega} \, d \Sigma_{a,1}(\omega) f(\tilde x) \, d \tilde x
\\
& =
\lambda^{-1} \int_{\tilde x\,\in\,\R^{n}} \int_{\omega \,\in\,\Sigma_{a,1}} e^{\i \left(\lambda x- \tilde x\right) \cdot \omega} \, d \Sigma_{a,1}(\omega) f\left(\frac{\tilde x}\lambda\right) \, d \tilde x, \qquad \qquad x \in \R^n.
\end{align*}
We thus obtain
\begin{equation*}
\left(T_{a,\lambda}^\star T_{a, \lambda}f \right)\left(x\right) =
\lambda^{-1} \left(T_{a,1}^\star T_{a, 1} \left(f\left(\frac{\cdot}{\lambda}\right)\right)\right)(\lambda x),
\qquad \qquad x \in \R^n.
\end{equation*}
From this identity, a simple scaling argument shows that for all $a \in [X_0, X_1] $ and all $\lambda>0 $,
\begin{equation*}
\norm{T_{a,\lambda}^\star T_{a, \lambda} }_{\scL\left(L^{\frac{2(n+1)}{(n+3)}}(\R^n),L^{\frac{2(n+1)}{(n-1)}}(\R^n)\right)}
\leq
\lambda^{\frac{n-1}{n+1}} \norm{T_{a,1}^\star T_{a, 1} }_{\scL\left(L^{\frac{2(n+1)}{(n+3)}}(\R^n),L^{\frac{2(n+1)}{(n-1)}}(\R^n)\right)}.
\end{equation*}
Since $T_{a,1}^\star T_{a, 1} =T_{a}^\star T_{a} $ is the operator defined in~\eqref{Def-T-star-T}, and since $T_a$ belongs to $\scL(L^{\frac{2(n+1)}{(n+3)}}(\R^n),L^{2}(\R^n))$ from Theorem~\ref{Thm-Fourier-Rest-Sigma-x-1}, we deduce that there exists $C>0$ such that for all $a \in [X_0, X_1]$, for all $\lambda >0$,
\begin{align*}
& \norm{T_{a,\lambda}^\star T_{a, \lambda} }_{\scL\left(L^{\frac{2(n+1)}{(n+3)}}(\R^n),L^{\frac{2(n+1)}{(n-1)}}(\R^n)\right)}
\leq C \lambda^{\frac{n-1}{n+1}},
\\
& \norm{T_{a, \lambda} }_{\scL\left(L^{\frac{2(n+1)}{(n+3)}}(\R^n), L^2(\Sigma_{a,\lambda})\right)} = \norm{T_{a,\lambda}^\star}_{\scL\left(L^{2}(\Sigma_{a,\lambda}),L^{\frac{2(n+1)}{(n-1)}}(\R^n)\right)}
\leq C \lambda^{\frac{(n-1)}{2(n+1)}}.
\end{align*}

The second step then consists on rewriting the operator $K_{a,k}$ as follows:
\[
K_{a,k} = \int_{\lambda >0} T_{a,\lambda}^\star M_{a,k(\lambda, \cdot)} T_{a, \lambda} \, d\lambda,
\]
where $M_{a, k(\lambda, \cdot)}$ is the operator defined from $ L^2(\Sigma_{a, \lambda})$ to itself as follows: for $g \in L^2(\Sigma_{a,\lambda})$,
\[
M_{a, k(\lambda, \cdot)} g (\xi) = k\left(\lambda, \frac{\xi}{\psi(a, \xi)} \right) g(\xi), \qquad \qquad \xi \in \Sigma_{a,\lambda}.
\]
Accordingly, we have
\begingroup
\allowdisplaybreaks
\begin{align*}
\lefteqn{\norm{K_{a,k} }_{\scL\left(L^{\frac{2(n+1)}{(n+3)}}(\R^n),L^{\frac{2(n+1)}{(n-1)}}(\R^n)\right)} }
\\
& \leq
\int_{\lambda\, >\,0}
\norm{T_{a,\lambda}^\star M_{a,k(\lambda, \cdot)} T_{a, \lambda} }_{\scL\left(L^{\frac{2(n+1)}{(n+3)}}(\R^n),L^{\frac{2(n+1)}{(n-1)}}(\R^n)\right)}
\, d\lambda
\\
&\leq \int_{\lambda\, >\,0}
\norm{T_{a,\lambda}^\star}_{\scL\left(L^{2}(\Sigma_{a,\lambda}),L^{\frac{2(n+1)}{(n-1)}}(\R^n)\right)}
\norm{M_{a,k(\lambda, \cdot)} }_{\scL\left(L^{2}(\Sigma_{a,\lambda})\right)}\\
&\mkern60mu
\norm{T_{a, \lambda} }_{\scL\left(L^{\frac{2(n+1)}{(n+3)}}(\R^n),L^{2}(\Sigma_{a,\lambda})\right)}
\! d\lambda
\\
& \leq \int_{\lambda\,>\,0} \| k (\lambda, \cdot)\|_{L^\infty(\Sigma_a)} \lambda^{\frac{n-1}{n+1}} \, d\lambda,
\end{align*}
\endgroup
where we used the straightforward estimates:
\[
\norm{M_{a,k(\lambda, \cdot)} }_{\scL\left(L^{2}(\Sigma_{a,\lambda})\right)}
\leq
\| k(\lambda, \cdot) \|_{L^\infty(\Sigma_a)}.
\]
This concludes the proof of the estimate~\eqref{Est-K-a-k-Lp-Lp'}.

Estimates~\eqref{Est-K-a-k-Lp-L2}--\eqref{Est-K-a-k-L2-Lp'} are based on the estimate~\eqref{Est-K-a-k-Lp-Lp'} and the facts that the map $K_{a,k}$ satisfies $K_{a,k}^\star = K_{a,\overline{k}}$ when computing the adjoint with respect to the $L^2(\R^n)$ scalar product, and $K_{a,k}^\star K_{a,\overline k} = K_{a,|k|^2}$. Therefore,
\begin{multline*}
\norm{K_{a,k} }_{\scL\left(L^{\frac{2(n+1)}{(n+3)}}(\R^n),L^{2}(\R^n)\right)}^2\\
\begin{aligned}
&=
\norm{K_{a,\overline k} }_{\scL\left(L^{2}(\R^n),L^{\frac{2(n+1)}{(n-1)}}(\R^n)\right)}^2
\\
&=
\norm{K_{a,|k|^2} }_{\scL\left(L^{\frac{2(n+1)}{(n+3)}}(\R^n),L^{\frac{2(n+1)}{(n-1)}}(\R^n)\right)} =
\int_0^\infty \| k(\lambda, \cdot)\|_{L^\infty(\Sigma_a)}^2 \, \lambda^{\frac{n-1}{n+1}} \, d\lambda,
\end{aligned}
\end{multline*}
which concludes the proof of Proposition~\ref{Prop-Est-Fourier-Op}, up to exchanging $k$ and $\overline k$ in the above formulae.
\end{proof}


\section{\texorpdfstring{$L^p$}{Lp}-Estimates on the parametrix, and proofs of Theorem~\ref{Thm-Carleman-Strip} and Theorem~\ref{Thm-Carleman-Strip-2}}\label{Sec-Est-Parametrix}

This section is devoted to give estimates on the norms of the various operators appearing in Proposition~\ref{Prop-parametrix-Strip-1}, especially in the spaces $\scL(L^p(\Omega), L^q(\Omega))$ for suitable values of $p$ and $q$. This will be done in particular by using the results in Proposition~\ref{Prop-Est-Fourier-Op} with $n =d -1$ and the Hardy--Littlewood--Sobolev theorem (Theorem~\ref{Thm-Hardy-Littlewood-Sobolev}).

We will also repeatedly use the straightforward lemma below, whose proof is left to the reader.

\begin{lemm}\label{lem-general-inequalities}\leavevmode
\begin{enumerate}
\item \label{item-1-lemma-6.1}
For all $\alpha \! \in \R$, and $a \! \in \! [\frac{1}{2}, 2]$, there exists $C>0$ such that for all $\mu >1$,\linebreak $
%\ds
\int_1^\mu e^{a\lambda} \lambda^\alpha \, d\lambda \leq C e^{a \mu} (1+ \mu)^\alpha. $

\item\label{item-2-lemma-6.1}
For all $\alpha >-1$, there exists $C>0$ such that for all $\mu >0$, $
%\ds
\int_0^\mu e^{\lambda} \lambda^\alpha d\lambda \leq C e^{ \mu} \frac{\mu^{\alpha+1}}{1+ \mu}. $

\item \label{item-3-lemma-6.1}
For all $\alpha > -1$, there exists $C >0$ such that for all $\gamma >0$, $
%\ds
\int_{\lambda\,>\,\gamma} e^{-\lambda} \lambda^{\alpha} d \lambda \leqslant C e^{-\gamma} (1+ \gamma)^\alpha. $
\item \label{item-4-lemma-6.1}
For all $\alpha < -1$, there exists $C >0$ such that for all $\gamma >0$, $
%\ds
\int_{\lambda\,>\,\gamma} e^{-\lambda} \lambda^{\alpha} d \lambda \leqslant C e^{-\gamma} \frac{\gamma^{\alpha+1}}{1+\gamma}. $
\item \label{item-5-lemma-6.1}
For all $\alpha \in (-3,0)$ and $a \in [\frac{1}{2}, 2]$, there exists $C >0$ such that for all $\gamma >0$,
$
%\ds
\int_{\lambda\,>\,\gamma} (\lambda - \gamma)^2 e^{-a \lambda} \lambda^{\alpha} \, d \lambda
\leq C \frac{e^{- a \gamma}}{1+ \gamma^{- \alpha }}. $
\end{enumerate}
\end{lemm}



In the whole section, we assume the setting of Theorem~\ref{Thm-Carleman-Strip}. Within this setting, with $\psi$ defined as in~\eqref{Def-Psi}, there exists $c_1 >0$ depending only on $c_0$ in~\eqref{Coercivity} such that
\begin{equation}\label{Equivalences-Psi-x-xi}
\forall (x_1,y_1) \in [X_0,X_1]^2,\, \forall \xi' \in \R^{d-1},
\quad
\frac{1}{c_1} \psi(y_1,\xi') \leq \psi(x_1, \xi') \leq c_1 \psi(y_1,\xi').
\end{equation}

We also recall that in this section, all the constants $C$ depend only on $c_0$ in~\eqref{Coercivity} and $m_*, M_*$ in~\eqref{Sub-ellipticity-Strip}.



\subsection{Estimates on the operator \texorpdfstring{$K_{\tau, 0}$}{K tau, 0} in~\eqref{Def-K-tau-j-Fourier}--\eqref{Def-k-tau-0-x1-y1-xi'}}

The goal of this section is to estimate the norm of the operator $K_{\tau, 0}$ in~\eqref{Def-K-tau-j-Fourier}--\eqref{Def-k-tau-0-x1-y1-xi'}, more precisely:

\begin{prop}\label{Prop-OpNorm-Est-K0}
Let $\Omega$ be as in~\eqref{Omega=Strip} with $X_0 \! <\! 0\! <\! X_1$ and $\max\{|X_0|, |X_1|\}\linebreak \leq 1$, and assume that the coefficients $(\lambda_j)_{j \in \{1,\cdots,d\}} \in \R^d$ satisfy $\lambda_1 = 0$, \eqref{Coercivity} and~\eqref{Sub-ellipticity-Strip}. Then there exist $C>0$ and $\tau_0 \geq 1$ independent of $X_0, X_1$ (and depending only on $c_0$, $m_*$ and $M_*$ in~\eqref{Coercivity} and~\eqref{Sub-ellipticity-Strip}), such that for all $\tau \geq \tau_0$, for all $f \in L^{\frac{2d}{d+2}}(\Omega)$,
\begin{multline}\label{Est-K-tau-0-f-L-2_*}
\left\| K_{\tau,0} f\right\|_{L^{\frac{2d}{d-2}}(\Omega)} +
\tau^{\frac{3}{4}+\frac{1}{2d}}\left \| K_{\tau,0} f\right\|_{L^{2}(\Omega)} +
\left\| \partial_1 \widehat{ K_{\tau,0} f}\right\|_{L^2(\Omega_{1, \tau})}\\
+
\tau^{-\frac{1}{4} +\frac{1}{2d}} \left\| \nabla' K_{\tau,0} f\right\|_{L^2(\Omega)}
\leq C \| f \|_{L^{\frac{2d}{d+2}}(\Omega)},
\end{multline}
and, for all $f \in L^2(\Omega)$,
\begin{multline}\label{Est-K-tau-0-f-L-2}
\tau^{\frac{3}{4}+\frac{1}{2d}} \left\| K_{\tau,0} f\right\|_{L^{\frac{2d}{d-2}}(\Omega)} +
\tau^{\frac{3}{2}} \left\| K_{\tau,0} f\right\|_{L^{2}(\Omega)} +
\tau
\left\| \partial_1 \widehat{ K_{\tau,0} f}\right\|_{L^2(\Omega_{1, \tau})} \\
+
\tau^{\frac{1}{2}} \left\| \nabla' K_{\tau,0} f\right\|_{L^2(\Omega)}
\leq C \| f \|_{L^2(\Omega)},
\end{multline}
with $\Omega_{1, \tau}$ as in~\eqref{Def-Omega-1-tau}.
\end{prop}
\begin{rema}
In the above estimates, we point out that the $\scL(L^2(\Omega),L^{\frac{2d}{d-2}}(\Omega))$ and $\scL(L^{\frac{2d}{d+2}}(\Omega),L^2(\Omega))$ bounds of the operator $K_{\tau,0}$ are estimated by a power of the Carleman parameter that depends on $d$. This fact, which does not occur for the Hilbertian estimates, has been already observed in several cases, and we refer for instance to~\cite{barcelo1988weighted, koch2001carleman, koch2005dispersive, Sogge1989uniquecontinuation}.
\end{rema}


\begin{proof}
In view of the results in Proposition~\ref{Prop-Est-Fourier-Op}, we first estimate weighted norms of $k_{\tau,0}(x_1, y_1, \cdot)$ for $x_1$ and $y_1$ in $[X_0,X_1]$ (recall the definition of $k_{\tau,0}$ in~\eqref{Def-k-tau-0-x1-y1-xi'}). We also identify $\xi' \in \R^{d-1}$ with pairs $(\lambda, \omega') \in \R_+ \times \Sigma_{x_1}$, where $\Sigma_{x_1} = \{ \omega' \in \R^{d-1}, \psi(x_1, \omega')\linebreak = 1\}$, through the formula $\xi' = \lambda \omega'$, or equivalently $\lambda = \psi(x_1,\xi')$ and $\omega' = \xi' /\psi(x_1,\xi')$. With a slight abuse of notations, we denote $k_{\tau,0}$ similarly whether it is written in terms of $\xi' \in \R^{d-1}$ or in terms of $(\lambda, \omega') \in \R_+ \times \Sigma_{x_1}$, that is
\[
k_{\tau,0}(x_1, y_1, \lambda, \omega') = k_{\tau,0}(x_1, y_1, \lambda \omega').
\]

We begin with the following lemma:

\begin{lemm}\label{Lem-Est-k-0-x1-y1}
There exist constants $C>0$ and $C_1 >0$ independent of $X_0, X_1$ (and depending only on $c_0$, $m_*$ and $M_*$ in~\eqref{Coercivity} and~\eqref{Sub-ellipticity-Strip}), such that for all $x_1$ and $y_1$ in $[X_0, X_1]$, for all $\tau \geq 1$, and $\lambda >0$,
\begin{itemize}
\item If $\lambda \leq \tau$, then the kernel $k_{\tau,0}$ defined in~\eqref{Def-k-tau-0-x1-y1-xi'} satisfies
\begin{multline}\label{Est-k-tau-0-x1-y1-Low-Freq}
\| k_{\tau, 0}(x_1,y_1,\lambda, \cdot) \|_{L^\infty\left(\Sigma_{x_1}\right)}\\
\leq
\left\{
\begin{array}{ll}
\ds C | y_1- x_1 | e^{-\tau | y_1- x_1 | }, \quad & \text{ if } \ \lambda |y_1 -x_1| \leqslant 1,
\smallskip
\\
\ds \frac{C}{\lambda }e^{- (\tau - \lambda) |y_1- x_1| - \lambda (y_1- x_1) ^ 2/C_1}, \quad & \text{ if } \lambda |y_1 -x_1| \geqslant 1.
\end{array}
\right.
\end{multline}
\item If $\lambda \geq \tau$, then $k_{\tau,0}$ satisfies
\begin{equation}\label{Est-k-tau-0-x1-y1-High-Freq}
\| k_{\tau, 0}(x_1,y_1,\lambda, \cdot) \|_{L^\infty\left(\Sigma_{x_1}\right)}
\leq
\left\{
\begin{array}{ll}
\ds \frac{C}{\lambda} e^{-(\lambda-\tau) |x_1-y_1|- \lambda (x_1-y_1)^2/C_1}, & \text{ if } y_1<x_1,
\smallskip
\\
\ds \frac{C}{\lambda} e^{-(\lambda/C + \tau)|y_1 - x_1|}, & \text{ if } y_1>x_1.
\end{array}
\right.
\end{equation}
\end{itemize}
Setting, for $x_1$ and $y_1$ in $[X_0, X_1]$, $\xi' \in \R^{d-1}$,
\begin{multline}\label{Def-k-tau-0-d1}
k_{\tau,0,\partial_1} (x_1, y_1, \xi') = -1_{x_1 < y_1} e^{-\tau(y_1-x_1) - \int_{x_1}^{y_1} \psi\left(\tilde y_1, \xi'\right) \, d\tilde y_1}
\\
\begin{aligned}
& - 1_{\psi(x_1, \xi')\,>\,\tau} (\tau - \psi(x_1, \xi'))s
\int_{X_0}^{\min\{x_1,y_1\}} e^{-\tau(y_1-x_1) - \int_{\tilde x_1}^{x_1} \psi\left(\tilde y_1, \xi'\right) \, d\tilde y_1 - \int_{\tilde x_1}^{y_1} \psi\left(\tilde y_1, \xi'\right) \, d\tilde y_1} d\tilde x_1
\\
& + 1_{\psi(x_1, \xi')\,\leq\,\tau}\; 1_{x_1\,<\,y_1} \left(\tau - \psi(x_1, \xi')\right) \int_{x_1}^{y_1} e^{-\tau(y_1-x_1) + \int_{x_1}^{\tilde x_1} \psi\left(\tilde y_1, \xi'\right) \, d\tilde y_1 - \int_{\tilde x_1}^{y_1} \psi\left(\tilde y_1, \xi'\right) \, d\tilde y_1} d\tilde x_1,
\end{aligned}
\end{multline}
$k_{\tau,0,\partial_1}$ satisfies the following bounds:
\begin{itemize}
\item If $\lambda \leq \tau$,
\begin{multline}\label{Est-k-tau-0-d1-x1-y1-Low-Freq}
\|k_{\tau, 0, \partial_1}(x_1,y_1,\lambda, \cdot) \|_{L^\infty\left(\Sigma_{x_1}\right)}\\
\leq C e^{-(\lambda/C+\tau)|y_1 -x_1|} + (\tau -\lambda)
\| k_{\tau,0}(x_1, y_1,\lambda, \cdot)\|_{L^\infty\left(\Sigma_{x_1}\right)}.
\end{multline}
\item If $\lambda > \tau$,
\begin{multline}\label{Est-k-tau-0-d1-x1-y1-High-Freq}
\|k_{\tau, 0, \partial_1}(x_1,y_1,\lambda, \cdot) \|_{L^\infty\left(\Sigma_{x_1}\right)}\\
\leq
\left\{
\begin{array}{ll}
\ds
\frac{C}{\lambda} (\lambda - \tau) e^{-(\lambda-\tau) |x_1-y_1|- \lambda (x_1-y_1)^2/C_1}, & \text{ if } y_1<x_1,
\smallskip
\\
\ds C e^{-(\lambda/C + \tau)(y_1 - x_1)}, & \text{ if } y_1>x_1.
\end{array}
\right.
\end{multline}
\end{itemize}
\end{lemm}

\begin{rema}
The kernel $k_{\tau,0, \partial_1}$ corresponds to the kernel of $\partial_1 K_{\tau,0}$ in the following sense: for all $f \in L^2((X_0,X_1); L^2(\R^{d-1}))$, and all $(x_1, \xi') \in \Omega_{1, \tau}$,
\begin{equation}\label{d1-Hat-K}
\partial_1 \widehat{K_{\tau, 0} f}(x_1,\xi') = \int_{y_1 \in (X_0,X_1)} k_{\tau,0, \partial_1}(x_1, y_1, \xi') \widehat f(y_1, \xi') \, dy_1.
\end{equation}
\end{rema}

\begin{proof}[Proof of Lemma~\ref{Lem-Est-k-0-x1-y1}]
Let us first prove~\eqref{Est-k-tau-0-x1-y1-Low-Freq} corresponding to $\lambda \leq \tau$. Let $\xi' \in \R^{d-1}$ and $x_1 \in [X_0, X_1]$ be such that $\lambda = \psi(x_1,\xi') \leq \tau$. We then have to estimate, for $y_1 > x_1$,
\[
\int_{x_1}^{y_1} e^{-\tau (y_1-x_1)+\int_{x_1}^{ \tilde{x}_1} \psi\left(\tilde{y}_1,\xi'\right) \, d \tilde{y}_1- \int_{{\tilde{x}_1}}^{y_1} \psi\left(\tilde{y}_1,\xi'\right) \, d \tilde{y}_1} \, d \tilde{x}_1.
\]
For $ y_1 \in [x_1, X_1]$, we introduce the map
\[
\rho(x_1,\tilde x_1, y_1,\xi') =
\int_{x_1}^{\tilde{x}_1} \psi(\tilde{y}_1,\xi') \, d \tilde{y}_1 -
\int_{{\tilde{x}_1}}^{y_1} \psi(\tilde{y}_1,\xi') \,d \tilde{y}_1,
\]
which clearly satisfies $\rho(x_1,y_1, y_1, \xi') = - \rho(x_1,x_1,y_1, \xi') \geq |y_1 - x_1| \lambda /c_1$, and \\ $\partial_{\tilde x_1} \rho(x_1, \tilde x_1, y_1, \xi') = 2 \psi(\tilde x_1, \xi')$ and $\partial_{\tilde x_1}^2 \rho(x_1, \tilde x_1, y_1, \xi') = 2 \partial_1 \psi(\tilde x_1, \xi') <0$ by \eqref{Sub-ellipticity-Strip}. Therefore, by concavity in $\tilde x_1$, for $x_1 < y_1$ and $\tilde x_1 \in [x_1, y_1]$,
\begin{align*}
\rho(x_1, \tilde x_1, y_1, \xi')
&\leq \rho(x_1,y_1,y_1, \xi') - 2 \psi(y_1,\xi') |y_1 - \tilde x_1|\\
&\leq \int_{x_1}^{y_1} \psi(\tilde{y}_1,\xi') \, d \tilde{y}_1- \frac{2}{c_1} \lambda |y_1 - \tilde x_1|,
\end{align*}
where the last estimate follows from~\eqref{Equivalences-Psi-x-xi}. Hence, we obtain
\begin{multline*}
\left|
\int_{x_1}^{y_1} e^{-\tau (y_1-x_1)+\int_{x_1}^{ \tilde{x}_1} \psi\left(\tilde{y}_1,\xi'\right) \, d \tilde{y}_1- \int_{{\tilde{x}_1}}^{y_1} \psi\left(\tilde{y}_1,\xi'\right) \, d \tilde{y}_1} \, d \tilde{x}_1
\right|
\\
\begin{aligned}
&\leq e^{-\tau (y_1-x_1)} \int_{x_1}^{y_1} e^{\int_{x_1}^{y_1} \psi\left(\tilde{y}_1,\xi'\right) \, d \tilde{y}_1- {2} \lambda \left|y_1 - \tilde x_1\right|/{c_1}}
\, d \tilde{x}_1
\\
&\leq e^{-\tau(y_1 - x_1) + \int_{x_1}^{y_1} \psi\left(\tilde{y}_1,\xi'\right) \, d \tilde{y}_1} \min\left\{ |y_1 - x_1|, \frac{c_1}{2 \lambda}\right\}.
\end{aligned}
\end{multline*}
We then use that the function $\tilde y_1 \mapsto \psi(\tilde y_1, \xi') $ is concave, so that for $x_1 < \tilde y_1$, we have $\psi(\tilde y_1, \xi') \leq \psi(x_1, \xi') + \partial_1 \psi(x_1, \xi') (\tilde y_1 -x_1)$. From~\eqref{Coercivity-Restriction} and~\eqref{Sub-ellipticity-Strip}, there exists a constant $c_2>0$ depending only on $c_0$, $m_*$ and $M_*$ such that $\partial_1 \psi(x_1, \xi') \leq - c_2 \psi(x_1, \xi') = - c_2 \lambda$. Therefore, for $y_1 > x_1$, $
\int_{x_1}^{y_1} \psi(\tilde{y}_1,\xi') \, d \tilde{y}_1
\leq (y_1 -x_1) \lambda - c_2 \lambda (y_1 -x_1)^2, $ and~\eqref{Est-k-tau-0-x1-y1-Low-Freq} follows immediately.


We then prove~\eqref{Est-k-tau-0-x1-y1-High-Freq} corresponding to $\lambda > \tau$. Let $\xi' \in \R^{d-1}$ and $x_1 \in [X_0,X_1]$ be such that $\lambda = \psi(x_1, \xi') > \tau$. We then have to estimate, for $y_1 \in [X_0,X_1]$, the quantity
\[
\int_{X_0}^{\min \{ x_1, y_1\}} e^{-\tau(y_1-x_1)-\int_{\tilde{x}_1}^{x_1} \psi\left(\tilde{y}_1,\xi'\right) \, d\tilde{y}_1 - \int_{{\tilde{x}_1}}^{y_1} \psi\left(\tilde{y}_1,\xi'\right) \,d \tilde{y}_1 } \, d \tilde{x}_1.
\]
As before, one easily checks that the map $\tilde x_1 \mapsto -\int_{\tilde{x}_1}^{x_1} \psi(\tilde{y}_1,\xi') \, d\tilde{y}_1 - \int_{{\tilde{x}_1}}^{y_1} \psi(\tilde{y}_1,\xi') \,d \tilde{y}_1$ has derivative $2 \psi(\tilde x_1, \xi')$ and is thus strictly increasing and concave. Therefore, for all $\tilde x_1 \in [X_0, \min\{x_1,y_1\}]$,
\begin{multline*}
-\int_{\tilde{x}_1}^{x_1} \psi(\tilde{y}_1,\xi') \, d\tilde{y}_1 - \int_{{\tilde{x}_1}}^{y_1} \psi(\tilde{y}_1,\xi') \,d \tilde{y}_1
\\
\begin{aligned}
& \leq - \int_{\min\{x_1,y_1\}}^{\max\{x_1,y_1\}} \psi(\tilde{y}_1,\xi') \, d\tilde{y}_1 + 2 \psi(\min\{x_1, y_1\}, \xi') (\tilde x_1 - \min\{x_1, y_1\})
\\
& \leq - \int_{\min\{x_1,y_1\}}^{\max\{x_1,y_1\}} \psi(\tilde{y}_1,\xi') \, d\tilde{y}_1 - \frac{2}{c_1} \lambda (\min\{x_1, y_1\} - \tilde x_1).
\end{aligned}
\end{multline*}
Accordingly,
\begin{multline*}
\int_{X_0}^{\min \{ x_1, y_1\}} e^{-\tau(y_1-x_1)-\int_{\tilde{x}_1}^{x_1} \psi\left(\tilde{y}_1,\xi'\right) \, d\tilde{y}_1 - \int_{{\tilde{x}_1}}^{y_1} \psi\left(\tilde{y}_1,\xi'\right) \,d \tilde{y}_1 } \, d \tilde{x}_1
\\
\leq e^{-\tau (y_1 - x_1) - \int_{\min\{x_1,y_1\}}^{\max\{x_1,y_1\}} \psi\left(\tilde{y}_1,\xi'\right) \, d\tilde{y}_1 } \min\left\{ \frac{c_1}{2 \lambda}, |\min\{x_1, y_1\} - X_0| \right\}.
\end{multline*}
If $y_1 > x_1$, we simply use
\[
- \int_{\min\{x_1,y_1\}}^{\max\{x_1,y_1\}} \psi\left(\tilde{y}_1,\xi'\right) \, d\tilde{y}_1
\leq - \frac{1}{c_1} (y_1- x_1) \lambda.
\]
If $y_1 < x_1$, we use that $\psi$ is concave decreasing and thus for all $\tilde y_1 \in [y_1, x_1]$, $\psi(\tilde y_1, \xi') \geq \psi(x_1, \xi') + (\tilde y_1 - x_1) \partial_1 \psi(y_1, \xi')$. But there exists a constant $c_3>0$ depending on $c_0$, $m_*$ and $M_*$ such that $\partial_1 \psi(y_1, \xi') \leq - \psi(x_1, \xi')/c_3$, so that we easily get
\[
- \int_{\min\{x_1,y_1\}}^{\max\{x_1,y_1\}} \psi(\tilde{y}_1,\xi') \, d\tilde{y}_1
\leq - (y_1- x_1) \lambda - \frac{1}{c_3} (y_1- x_1)^2 \lambda.
\]
in this case. Combining the last three estimates immediately yields~\eqref{Est-k-tau-0-x1-y1-High-Freq}.

The proof of estimates~\eqref{Est-k-tau-0-d1-x1-y1-Low-Freq}--\eqref{Est-k-tau-0-d1-x1-y1-High-Freq} follows from the fact that, for $y_1 > x_1$,
\[
\left\| e^{-\tau(y_1-x_1) - \int_{x_1}^{y_1} \psi\left(\tilde y_1, \xi'\right) \, d\tilde y_1}\right\|_{L^\infty_{\omega'}\left(\Sigma_{x_1}\right)}
\leq e^{-(\tau +\lambda/c_1) (y_1 - x_1)},
\]
and from the estimates already proved above. Details are left to the reader.
\end{proof}

Using the bounds in Lemma~\ref{Lem-Est-k-0-x1-y1}, we prove the following lemma:

\begin{lemm}\label{Lem-Est-k0-x1-y1-integratednorms}
There exists a constant $C >0$ independent of $X_0, X_1$ (and depending only on $c_0$, $m_*$ and $M_*$ in~\eqref{Coercivity} and~\eqref{Sub-ellipticity-Strip}), such that for all $x_1$ and $y_1$ in $[X_0, X_1]$, for all $\tau \geq 1$,
\begin{align}
\int_{\lambda\,>\,0} \norm{k_{\tau,0}(x_1, y_1, \lambda, \cdot) }_{L^\infty\left(\Sigma_{x_1}\right)} \, \lambda^{1 -\frac{2}{d}}\, d\lambda
&\leq
\frac{C}{|x_1 - y_1|^{1-\frac{2}{d}}},\label{est-k0-x1-y1-Lp-Lp'}
\\
\left(\int_{\lambda\,>\,0} \norm{k_{\tau,0}(x_1, y_1, \lambda, \cdot) }^2_{L^\infty\left(\Sigma_{x_1}\right)} \, \lambda^{1 -\frac{2}{d}}\, d\lambda\right)^{\frac{1}{2}}
&\leq C \tilde k_{\tau,0}(x_1-y_1),\label{est-k0-x1-y1-Lp-L2}
\\
\text{ with }\tilde k_{\tau,0} \in L^{\frac{d}{d-1}} (\R)\text{ and }
\left\|\tilde k_{\tau,0}\right\|_{L^{\frac{d}{d-1}} (\R)} &\leq C \tau^{-\frac{3}{4}-\frac{1}{2d}},
\notag
\\
\left(
\int_{\lambda\,>\,0} \norm{ k_{\tau,0, \partial_1 }(x_1, y_1, \lambda, \cdot) }^2_{L^\infty\left(\Sigma_{x_1}\right)} \, \lambda^{1 -\frac{2}{d}}\, d\lambda
\right)^{\frac{1}{2}}
&\leq
\frac{C}{|x_1 - y_1|^{1-\frac{1}{d}} },\label{est-k0-d1-x1-y1-Lp-L2}
\\
\bigg(\int_{\lambda\,>\,0} \norm{\lambda \omega' k_{\tau,0}(x_1, y_1, \lambda, \omega') }^2_{L^\infty_{\omega'}\left(\Sigma_{x_1}\right)} \, \lambda^{1 -\frac{2}{d}}\, d\lambda&\bigg)^{\frac{1}{2}}\label{est-k0-x1-y1-Lp-nabla'-L2}\\
&\leq
\frac{C}{\left|x_1 - y_1\right|^{1-\frac{1}{d}}} + \check k_{\tau,0}(x_1 -y_1),\nonumber
\\
\text{ with }\check k_{\tau,0} \in L^{\frac{d}{d-1}} (\R)\text{ and }
\left\|\check k_{\tau,0}\right\|_{L^{\frac{d}{d-1}} (\R)} &\leq C \tau^{\frac{1}{4}-\frac{1}{2d}}.
\notag
\end{align}
\end{lemm}

\begin{proof}[Proof of Lemma~\ref{Lem-Est-k0-x1-y1-integratednorms}]
We start by simply noticing that we can always impose that $C_1$ in Lemma~\ref{Lem-Est-k-0-x1-y1} is large enough to get for all $(x_1, y_1) \in [X_0, X_1]^2$, that $|x_1 - y_1|^2/C_1 \leq |x_1 - y_1|/4$. This can be done by assuming for instance $C_1 \geq 8$ since $|X_0|, |X_1| \leq 1$. This will make some of the estimates below easier to prove properly. For convenience, this constant $C_1$ will next be denoted by $C$, similarly as generic constants which depend only on the dimension and the parameters $c_0$, $m_*$ and $M_*$ in~\eqref{Coercivity}, and~\eqref{Sub-ellipticity-Strip}.

\begin{proof}[{Proof of~\eqref{est-k0-x1-y1-Lp-Lp'}.}]
We decompose the integral in the left-hand side of~\eqref{est-k0-x1-y1-Lp-Lp'} in three terms more suitable to use the results in Lemma~\ref{Lem-Est-k-0-x1-y1}. The first is easily estimated as follows:
\begin{multline*}
\int_{0}^{ \min\{ \tau, 1/|y_1-x_1|\} } \norm{k_{\tau,0}(x_1, y_1, \lambda, \cdot) }_{L^\infty\left(\Sigma_{x_1}\right)} \, \lambda^{1 -\frac{2}{d}}\, d\lambda\\
\begin{aligned}
& \leq C \int_{0}^{ \min\{ \tau, 1/|y_1-x_1|\} } |y_1-x_1| e^{-\tau |y_1 -x_1|} \, \lambda^{1 -\frac{2}{d}}\, d\lambda
\\
&
\leq C |y_1-x_1| e^{-\tau |y_1 -x_1|} \min\left\{ \tau, \frac{1}{|y_1-x_1|}\right\}^{2 - \frac{2}{d}}
\\
&
\leq C |y_1-x_1|^{-1+\frac{2}{d}} e^{-\tau |y_1 -x_1|} \min\left\{ \tau |y_1 -x_1|, 1\right\}^{2 - \frac{2}{d}}
\leq C |y_1-x_1|^{-1+\frac{2}{d}}.
\end{aligned}
\end{multline*}
Next if $1/|y_1 -x _1| < \tau$, that is $\tau |y_1 - x_1| > 1$, we get
\begin{multline*}
\int_{ 1/|y_1-x_1|}^{ \tau } \norm{k_{\tau,0}(x_1, y_1, \lambda, \cdot) }_{L^\infty\left(\Sigma_{x_1}\right)} \, \lambda^{1 -\frac{2}{d}}\, d\lambda \\
\begin{aligned}
&\leq C \int_{ 1/|y_1-x_1|}^{ \tau } e^{-(\tau - \lambda)|y_1 -x_1|- \lambda |y_1-x_1|^2/C_1} \lambda^{ -\frac{2}{d}}\, d\lambda
\\
& \leq C e^{-\tau |y_1-x_1|}\int_{1/|y_1 -x_1|}^{\tau} e^{\lambda |y_1-x_1|} \lambda^{-\frac{2}{d}} \, d \lambda
\\
& \leq C |y_1 - x_1|^{-1+\frac{2}{d}} e^{-\tau |y_1-x_1|}\int_1^{\tau |y_1 -x_1|} e^\lambda \lambda^{-\frac{2}{d}} \, d\lambda
\\
& \leq C |y_1 - x_1|^{-1+\frac{2}{d}} \frac{1}{(1 + \tau |y_1 -x_1|)^{\frac{2}{d}}}
\leq C |y_1 -x_1|^{-1+\frac{2}{d}},
\end{aligned}
\end{multline*}
where, from the fourth to the fifth lines, we have used Lemma~\ref{lem-general-inequalities}$\MK$\eqref{item-1-lemma-6.1}. Finally, we also have, for $y_1 >x_1$,
\begin{multline*}
\int_{\lambda\,>\,\tau} \norm{k_{\tau,0}(x_1, y_1, \lambda, \cdot) }_{L^\infty\left(\Sigma_{x_1}\right)} \, \lambda^{1 -\frac{2}{d}}\, d\lambda\\
\begin{aligned}
&\leq C \int_{\lambda\,>\,\tau} e^{-(\lambda/C + \tau)(y_1 - x_1)} \lambda^{-\frac{2}{d}} \, d \lambda
\\
& \leq C e^{-\tau (y_1-x_1)} \int_{\lambda > 0} e^{-\lambda(y_1 - x_1)/C} \lambda^{-\frac{2}{d}} \, d\lambda
\\
&\leq Ce^{-\tau (y_1-x_1)} |y_1 -x_1|^{-1 + \frac{2}{d}}
\leq C |y_1 -x_1|^{-1+\frac{2}{d}}.
\end{aligned}
\end{multline*}
Accordingly, estimate~\eqref{est-k0-x1-y1-Lp-Lp'} holds for $y_1 > x_1$. Then for $y_1 < x_1$, it only remains to prove the following estimate, in which we use Lemma~\ref{lem-general-inequalities}$\MK$\eqref{item-3-lemma-6.1}:
\begin{multline*}
\int_{\lambda > \tau} \norm{k_{\tau,0}(x_1, y_1, \lambda, \cdot)}_{L^\infty\left(\Sigma_{x_1}\right)} \, \lambda^{1 -\frac{2}{d}}\, d\lambda\\
\begin{aligned}
&\leq C \int_{\lambda\, >\,\tau} e^{-(\lambda- \tau)(x_1 - y_1)- \lambda (x_1 - y_1)^2/C_1} \lambda^{-\frac{2}{d}} \, d \lambda
\\
& \leq C e^{\tau(x_1-y_1)}
\int_{\lambda>\tau} e^{-\lambda (x_1 - y_1)} \lambda^{-\frac{2}{d}} \, d \lambda
\\
& \leq C \left|x_1 - y_1\right|^{-1+\frac{2}{d}} e^{\tau(x_1-y_1)} \int_{\lambda\,>\,\tau (x_1 - y_1)} e^{-\lambda} \lambda^{-\frac{2}{d}} \, d \lambda\\
&\leq C \left|x_1 - y_1\right|^{-1+\frac{2}{d}}.
\end{aligned}
\end{multline*}

This concludes the proof of~\eqref{est-k0-x1-y1-Lp-Lp'} for $y_1 < x_1$ as well.
\end{proof}

\begin{proof}[{Proof of~\eqref{est-k0-x1-y1-Lp-L2}.}]
Of course, the proof of~\eqref{est-k0-x1-y1-Lp-L2} is very similar to the one of~\eqref{est-k0-x1-y1-Lp-Lp'}. We first have
\begin{multline*}
\int_{0}^{ \min\{ \tau, 1/|y_1-x_1|\} } \norm{k_{\tau,0}(x_1, y_1, \lambda, \cdot) }_{L^\infty\left(\Sigma_{x_1}\right)}^2 \, \lambda^{1 -\frac{2}{d}}\, d\lambda
\\ 
\begin{aligned}
&
\leq C \int_{0}^{ \min\{ \tau, 1/|y_1-x_1|\} } |y_1-x_1|^2 e^{-2 \tau |y_1 -x_1|} \, \lambda^{1 -\frac{2}{d}}\, d\lambda
\\
&
\leq C |y_1-x_1|^2 e^{-2\tau |y_1 -x_1|} \min\left\{ \tau, \frac{1}{|y_1-x_1|}\right\}^{2 - \frac{2}{d}}
\\
&\leq C \tau^{-\frac{2}{d}} \left((\tau |y_1-x_1|)^{\frac{2}{d}} e^{-2\tau |y_1 -x_1|} \min\left\{ \tau|y_1-x_1|, 1\right\}^{2 - \frac{2}{d}}\right).
\end{aligned}
\end{multline*}
Next if $1/|y_1 -x _1| < \tau$, that is $\tau |y_1 - x_1| > 1$, using Lemma~\ref{lem-general-inequalities}$\MK$\eqref{item-1-lemma-6.1}, we have
\begin{multline*}
\int_{1/|y_1 -x_1|}^{\tau}\norm{k_{\tau,0}(x_1, y_1, \lambda, \cdot) }_{L^\infty\left(\Sigma_{x_1}\right)}^2 \, \lambda^{1 -\frac{2}{d}}\, d\lambda \\
\begin{aligned}
&\leq C \int_{1/|y_1 -x_1|}^{\tau}e^{-2 (\tau - \lambda)|y_1 -x_1|- 2\lambda(y_1-x_1)^2/C_1} \lambda^{-1 -\frac{2}{d}}\, d\lambda
\\
&\leq C e^{-2 \tau |y_1-x_1|}\int_{1/|y_1 -x_1|}^{\tau} e^{2 \lambda\left(|y_1-x_1|- (y_1-x_1)^2/C_1\right)} \lambda^{-1-\frac{2}{d}} \, d \lambda
\\
&\leq C |y_1 - x_1|^{\frac{2}{d}} e^{-2\tau |y_1-x_1|}\int_1^{\tau |y_1 -x_1|} e^{2\lambda \left(1-|y_1-x_1|/C_1\right) } \lambda^{-1-\frac{2}{d}} \, d\lambda
\\
&\leq C |y_1 -x_1|^{\frac{2}{d}} \frac{e^{-\tau (y_1-x_1)^2/C}}{(1 + \tau |y_1 -x_1|)^{1+\frac{2}{d}}}\\
&\leq C \tau^{-\frac{2}{d}} \frac{e^{-\tau (y_1-x_1)^2/C}}{(1 + \tau |y_1 -x_1|)}.
\end{aligned}
\end{multline*}
Finally, we also have, for $y_1 > x_1$,
\begin{multline*}
\int_{\lambda\,>\,\tau} \norm{k_{\tau,0}(x_1, y_1, \lambda, \cdot) }^2_{L^\infty\left(\Sigma_{x_1}\right)} \, \lambda^{1 -\frac{2}{d}}\, d\lambda\\
\begin{aligned}
& \leq C \int_{\lambda\,>\,\tau} e^{-2(\lambda/C + \tau)(y_1 - x_1)} \lambda^{-1-\frac{2}{d}} \, d \lambda\\
& \leq C e^{- 2 \tau (y_1 - x_1)}
\int_{\lambda\,>\,\tau} e^{-2\lambda(y_1 - x_1)/C} \lambda^{-1-\frac{2}{d}} \, d \lambda\\
&\leq Ce^{- 2 \tau (y_1 - x_1)}
\int_{\lambda\,>\,\tau} \lambda^{-1-\frac{2}{d}} \, d \lambda\\
&
\leq C e^{-2\tau(y_1 -x_1)} \tau^{-\frac{2}{d}}.
\end{aligned}
\end{multline*}
On the other hand, for $x_1 > y_1$, we have that
\begin{multline*}
\int_{\lambda\,>\,\tau} \norm{k_{\tau,0}(x_1, y_1, \lambda, \cdot) }^2_{L^\infty\left(\Sigma_{x_1}\right)} \, \lambda^{1 -\frac{2}{d}}\, d\lambda\\
\begin{aligned}
&\leq C \int_{\lambda\,>\,\tau} e^{-2(\lambda- \tau)(x_1 - y_1)- 2\lambda (x_1 - y_1)^2/C_1} \lambda^{-1-\frac{2}{d}} \, d \lambda
\\
& \leq C e^{2\tau(x_1-y_1)}
\int_{\lambda\,>\,\tau} e^{-2 \lambda \left((x_1 - y_1)+(x_1- y_1)^2/C_1\right) } \lambda^{-1-\frac{2}{d}} \, d \lambda
\\
& \leq C e^{2\tau(x_1-y_1)} | x_1- y_1|^{\frac{2}{d}}
\int_{\lambda\,>\,\tau | x_1- y_1| } e^{-2 \lambda \left(1+(x_1- y_1)^2/C_1\right) } \lambda^{-1-\frac{2}{d}} \, d \lambda,
\end{aligned}
\end{multline*}
so that using
Lemma~\ref{lem-general-inequalities}$\MK$\eqref{item-4-lemma-6.1}, for $x_1>y_1$, we obtain
\begin{equation*}
\int_{\lambda\,>\,\tau} \norm{k_{\tau,0}(x_1, y_1, \lambda, \cdot) }^2_{L^\infty(\Sigma_{x_1})} \, \lambda^{1 -\frac{2}{d}}\, d\lambda
\leq C \tau^{-\frac{2}{d}} \frac{e^{-2 \tau (x_1-y_1)^2/C_1}}{1 + \tau (x_1 - y_1)}.
\end{equation*}
Therefore, combining the above estimates, we have
\begin{multline}\label{def-tilde-k-tau-0}
\left(
\int_{\lambda\,>\,0} \norm{k_{\tau,0}(x_1, y_1, \lambda, \cdot) }^2_{L^\infty(\Sigma_{x_1})} \, \lambda^{1 -\frac{2}{d}}\, d\lambda
\right)^{\frac{1}{2}}
\leq C \tilde k_{\tau,0}(x_1-y_1),
\\
\text{ where } \tilde k_{\tau,0}(z_1) = 1_{|z_1|\,<\,1/\tau} \tau^{-\frac{1}{d}} + 1_{|z_1| \,>\,1/\tau} \tau^{-\frac{1}{d}} \left((\tau |z_1|)^{\frac{1}{d}}e^{-\tau |z_1|} + \frac{e^{- \tau |z_1|^2/C} }{(\tau |z_1|)^{\frac{1}{2}}} \right).
\end{multline}
Easy computations then yield $
\| \tilde k_{\tau,0} \|_{L^{\frac{d}{d-1}}(\R)} \leq C \tau^{-\frac{3}{4}-\frac{1}{2d}} $ as announced.
\end{proof}

\begin{proof}[{Proof of~\eqref{est-k0-d1-x1-y1-Lp-L2}.}]
We have
\[
\int_{\lambda\,>\,0} e^{-2 (\lambda/C + \tau)|y_1-x_1|} \lambda^{1-\frac{2}{d}} \, d\lambda
\leq e^{-2 \tau |y_1 - x_1|} |y_1-x_1|^{-2+\frac{2}{d}}
\leq |y_1-x_1|^{-2+\frac{2}{d}}.
\]
We also have, as before,
\begin{multline*}
\int_{0}^{ \min\{ \tau, 1/|y_1-x_1|\} } (\tau-\lambda)^2 |y_1-x_1|^2 e^{-2 \tau |y_1 -x_1|} \, \lambda^{1 -\frac{2}{d}}\, d\lambda \\
\begin{aligned}
&
\leq
\tau^2 \int_{0}^{ \min\{ \tau, 1/|y_1-x_1|\} } |y_1-x_1|^2 e^{-2 \tau |y_1 -x_1|} \, \lambda^{1 -\frac{2}{d}}\, d\lambda
\\
& \leq C \tau^2 |y_1-x_1|^2 e^{-2\tau |y_1 -x_1|} \min\left\{ \tau, \frac{1}{|y_1-x_1|}\right\}^{2 - \frac{2}{d}}
\\
& \leq C |y_1 - x_1|^{\frac{2}{d}} \tau^2 e^{-2\tau |y_1 -x_1|} \min\left\{ \tau |y_1 - x_1|,1\right\}^{2 - \frac{2}{d}}
\\
& \leq C |y_1 - x_1|^{-2+\frac{2}{d}} (\tau |y_1 - x_1|) ^2 e^{-2\tau |y_1 -x_1|} \min\left\{ \tau |y_1 - x_1|,1\right\}^{2 - \frac{2}{d}}
\\
&
\leq C |y_1 - x_1|^{-2+\frac{2}{d}}.
\end{aligned}
\end{multline*}
If $1/|y_1 -x _1| < \tau$, that is $\tau |y_1 - x_1| > 1$, we distinguish the cases $1/|y_1-x_1 | \leq \tau/2$ and $1/|y_1-x_1| \geq \tau/2$.

%\break
\noindent If $1/|y_1 - x_1| \in [\tau/2, \tau]$, we have
\begin{multline*}
\int_{1/|y_1-x_1|}^{ \tau } (\tau - \lambda)^2 e^{-2 (\tau - \lambda)|y_1 -x_1|- 2\lambda(y_1-x_1)^2/C_1} \lambda^{-1 -\frac{2}{d}}\, d\lambda
\\
\begin{aligned}
& \leq C e^{-2 \tau |y_1-x_1|}
\int_{\tau/2}^\tau (\tau - \lambda) ^2 e^{2 \lambda\left(|y_1-x_1|- (y_1-x_1)^2/C_1\right)} \lambda^{-1-\frac{2}{d}} \, d \lambda
\\
&\leq C e^{-2 \tau |y_1-x_1|}
\tau^{-1-\frac{2}{d}} \int_{\tau/2}^\tau (\tau - \lambda) ^2 e^{2 \lambda\left(|y_1-x_1|- (y_1-x_1)^2/C_1\right)} \, d \lambda
\\
&\leq C e^{-2 \tau |y_1-x_1|}
\frac{\tau^{-1-\frac{2}{d}} e^{2 \tau \left(|y_1-x_1|- (y_1-x_1)^2/C_1\right)}}{|y_1 -x_1|^3}
\\
&
\leq C\frac{\tau^{-1-\frac{2}{d}} e^{-2 \tau (y_1-x_1)^2/C_1}}{|y_1 -x_1|^3}\\
&
\leq C |y_1 - x_1|^{-2+\frac{2}{d}} (\tau |y_1-x_1|)^{-1-\frac{2}{d}} e^{-2 \tau (y_1-x_1)^2/C}
\leq C |y_1 - x_1|^{-2+\frac{2}{d}}.
\end{aligned}
\end{multline*}
If ${1}/{| y_1 - x_1 | } \leqslant \tau /2 $, we split the integral into two parts, $\int_{1/| y_1 - x_1 | } ^{ \tau / 2 }$ and $\int_{\tau/2} ^{ \tau }$. The second integral has been estimated above, and, using
Lemma~\ref{lem-general-inequalities}$\MK$\eqref{item-1-lemma-6.1}, the first one is estimated as follows:
\begin{multline*}
\int_{\lambda\, \in\,(1/|y_1-x_1|, \tau/2) } (\tau - \lambda)^2 e^{-2 (\tau - \lambda)|y_1 -x_1|- 2\lambda(y_1-x_1)^2/C} \lambda^{-1 -\frac{2}{d}}\, d\lambda
\\
\begin{aligned}
&\leq e^{-2\tau |y_1 - x_1| } \tau^{2} \int_{1/| y_1 - x_1 | }^{\tau /2} e^{2 \lambda \left(|y_1 - x_1 | - (y_1 - x_1)^2/C\right) } \lambda^{-1 - 2 /d} d \lambda
\\
&\leq C e^{-2\tau |y_1-x_1|} \tau^{2} |y_1 - x_1|^{\frac{2}{d}} \int_1^{\tau |y_1 -x_1|/2} e^{2\lambda \left(1-(y_1-x_1)/C\right) } \lambda^{-1-\frac{2}{d}} \, d\lambda
\\
& \leq C |y_1 -x_1|^{\frac{2}{d}}\tau^{2} e^{ -\tau|y_1-x_1|}(\tau |y_1 -x_1|)^{-1-\frac{2}{d}}
\\
& \leq C |y_1 - x_1|^{-2+\frac{2}{d}} (\tau |y_1 - x_1|)^{1 - \frac{2}{d}} e^{ - \tau |y_1 - x_1|} \leq C |y_1 - x_1|^{-2+\frac{2}{d}}.
\end{aligned}
\end{multline*}

It remains to estimate, for $y_1 < x_1$,
\begin{multline*}
\int_{\lambda\,>\,\tau} (\lambda - \tau)^2 e^{-2(\lambda- \tau)(x_1 - y_1)- 2\lambda (x_1 - y_1)^2/C} \lambda^{-1-\frac{2}{d}} \, d \lambda
\\
\begin{aligned}
& \leq C e^{2\tau(x_1-y_1)}
\int_{\lambda\,>\,\tau} (\lambda- \tau)^2 e^{-2 \lambda ((x_1 - y_1)+(x_1- y_1)^2/C) } \lambda^{-1-\frac{2}{d}} \, d \lambda
\\
& \leq C e^{2\tau(x_1-y_1)}| x_1- y_1|^{-2+\frac{2}{d}}
\int_{\lambda\,>\,\tau |x_1- y_1 |} (\lambda- \tau| x_1- y_1|)^2 e^{-2 \lambda \left(1+(x_1- y_1)/C\right) } \lambda^{-1-\frac{2}{d}} \, d \lambda.
\end{aligned}
\end{multline*}

We now use Lemma~\ref{lem-general-inequalities}$\MK$\eqref{item-5-lemma-6.1} and obtain for $x_1 > y_1$,
\begin{multline*}
\int_{\lambda\,>\,\tau} (\lambda - \tau)^2 e^{-2(\lambda- \tau)(x_1 - y_1)- 2\lambda (x_1 - y_1)^2/C} \lambda^{-1-\frac{2}{d}} \, d \lambda \\
\leq
\frac{C |x_1 - y_1|^{-2+\frac{2}{d}} e^{-2 \tau (x_1-y_1)^2/C}}{ 1+ (\tau |x_1 -y_1|)^{1+\frac{2}{d}} }
\leq C |x_1 - y_1|^{-2+\frac{2}{d}}.
\end{multline*}
\end{proof}

\begin{proof}[{Proof of~\eqref{est-k0-x1-y1-Lp-nabla'-L2}.}]
We have
\begin{multline*}
\int_{0}^{\min\{ \tau, 1/|y_1-x_1|\}} 
\norm{\lambda \omega' k_{\tau,0}(x_1, y_1, \lambda, \omega') }^2_{L^\infty_{\omega'}\left(\Sigma_{x_1}\right)}
\, \lambda^{1 -\frac{2}{d}}\, d\lambda\\
\begin{aligned}
&
\leq C \int_{0}^{\min\{ \tau, 1/|y_1-x_1|\}} |y_1-x_1|^2 e^{-2 \tau (y_1 -x_1)} \, \lambda^{3 -\frac{2}{d}}\, d\lambda
\\
&
\leq C |y_1-x_1|^2 e^{-2\tau (y_1 -x_1)} \min\left\{ \tau, \frac{1}{|y_1-x_1|}\right\}^{4 - \frac{2}{d}}
\\
& \leq C |y_1-x_1|^{-2+\frac{2}{d}} e^{-2\tau (y_1 -x_1)} \!\min\left\{ \tau|y_1-x_1|, 1\right\}^{4 - \frac{2}{d}}
\leq \! C |y_1-x_1|^{-2+\frac{2}{d}}.
\end{aligned}
\end{multline*}
If $1/|y_1 -x _1| < \tau$, that is $\tau |y_1 - x_1| > 1$, using Lemma~\ref{lem-general-inequalities}$\MK$\eqref{item-1-lemma-6.1}, we obtain
\begin{multline*}
\int_{1/|y_1 -x_1|}^{\tau}  \norm{\lambda \omega' k_{\tau,0}(x_1, y_1, \lambda, \omega') }^2_{L^\infty_{\omega'}\left(\Sigma_{x_1}\right)} \, \lambda^{1 -\frac{2}{d}}\, d\lambda\\
\begin{aligned}
&
\leq  C \int_{1/|y_1 -x_1|}^{\tau} e^{-2 (\tau - \lambda)|y_1 -x_1|- 2\lambda(y_1-x_1)^2/C} \lambda^{1 -\frac{2}{d}}\, d\lambda
\\
&\leq  C e^{-2 \tau|y_1-x_1|}\int_{1/|y_1 -x_1|}^{\tau} e^{2 \lambda\left(|y_1-x_1|- (y_1-x_1)^2/C\right)} \lambda^{1-\frac{2}{d}} \, d \lambda
\\
& \leq  C |y_1 - x_1|^{-2+\frac{2}{d}} e^{-2\tau |y_1-x_1|} \int_1^{\tau |y_1 -x_1|}   e^{2\lambda \left(1-|y_1-x_1|/C\right)} \lambda^{1-\frac{2}{d}} \, d\lambda
\\
& \leq  C |y_1 -x_1|^{-2+\frac{2}{d}} e^{-\tau (y_1-x_1)^2/C} (\tau |y_1 -x_1|)^{1-\frac{2}{d}}
\\
& \leq  C \tau^{1-\frac{2}{d}}e^{-\tau (y_1-x_1)^2/C} |y_1 -x_1|^{-1}.
\end{aligned}
\end{multline*}
Finally, using Lemma~\ref{lem-general-inequalities}$\MK$\eqref{item-3-lemma-6.1}, we also have, for $y_1 > x_1$,
\begin{multline*}
\int_{\lambda\, >\,\tau} \norm{\lambda \omega' k_{\tau,0}(x_1, y_1, \lambda, \omega')}^2_{L^\infty_{\omega'}\left(\Sigma_{x_1}\right)} \, \lambda^{1 -\frac{2}{d}}\, d\lambda\\
\begin{aligned}
&\leq C \int_{\lambda\,>\,\tau} e^{-2(\lambda/C + \tau)(y_1 - x_1)} \lambda^{1-\frac{2}{d}} \, d \lambda
\\
&\leq C e^{- 2 \tau (y_1 - x_1)}
\int_{\lambda\,>\,\tau} e^{-2\lambda(y_1 - x_1)/C} \lambda^{1-\frac{2}{d}} \, d \lambda
\\
&\leq Ce^{- 2 \tau (y_1 - x_1)} |y_1 -x_1|^{-2+\frac{2}{d}} \int_{\lambda\,>\,\tau (y_1 - x_1)} e^{-2 \lambda/C} \lambda^{1-\frac{2}{d}} \, d\lambda
\\
&\leq C e^{-2\tau(y_1 -x_1)}|y_1 -x_1|^{-2+ \frac{2}{d}} (1+ (\tau |y_1 -x_1|))^{1-\frac{2}{d}}
\\
&
\leq C |y_1-x_1|^{-2+\frac{2}{d}}.
\end{aligned}
\end{multline*}

Similarly, for $x_1 > y_1$,
\begin{multline*}
\int_{\lambda\,>\,\tau}  \norm{\lambda \omega' k_{\tau,0}(x_1, y_1, \lambda, \omega') }^2_{L^\infty_{\omega'}\left(\Sigma_{x_1}\right)} \, \lambda^{1 -\frac{2}{d}}\, d\lambda\\
\begin{aligned}
&\leq C \int_{\lambda\,>\,\tau}  e^{-2(\lambda- \tau)(x_1 - y_1)- 2\lambda (x_1 - y_1)^2/C} \lambda^{1-\frac{2}{d}} \, d \lambda
\\
&\leq C e^{2\tau(x_1-y_1)}
\int_{\lambda\,>\,\tau} e^{-2 \lambda \left((x_1 - y_1)+(x_1- y_1)^2/C\right)} \lambda^{1-\frac{2}{d}} \, d \lambda
\\
&\leq C e^{2\tau(x_1-y_1)} | x_1- y_1 |^{-2 + \frac{2}{d}}
\int_{\lambda\,>\,\tau | x_1 - y_1 | } e^{-2 \lambda \left(1+(x_1- y_1)/C\right)} \lambda^{1-\frac{2}{d}} \, d \lambda,
\end{aligned}
\end{multline*}
so that using Lemma~\ref{lem-general-inequalities}$\MK$\eqref{item-3-lemma-6.1}, we get, for $x_1 > y_1$,
\begin{multline*}
\int_{\lambda\,>\,\tau} \norm{\lambda \omega' k_{\tau,0}(x_1, y_1, \lambda, \cdot) }^2_{L^\infty\left(\Sigma_{x_1}\right)} \, \lambda^{1 -\frac{2}{d}}\, d\lambda\\
\begin{aligned}
&\leq C |y_1 - x_1|^{-2 + \frac{2}{d}} e^{-2\tau(y_1-x_1)^2/C} (1 + \tau |y_1-x_1|)^{1 - \frac{2}{d}}
\\
&\leq C \left(1_{| y_1 - x_1 |\,\leq\,1/\tau } |y_1-x_1|^{-2+\frac{2}{d}} + 1_{| y_1 - x_1 | \,\geq\,1/\tau} \tau ^{1-\frac{2}{d}} | y_1 - x_1 |^{-1}e^{-2\tau(y_1-x_1)^2/C} \right).
\end{aligned}
\end{multline*}
Therefore, by combining the above estimates, we have
\begin{multline*}
\left(\int_{\lambda\,>\,0} \norm{\lambda \omega' k_{\tau,0}(x_1, y_1, \lambda, \cdot) }^2_{L^\infty(\Sigma_{x_1})} \, \lambda^{1 -\frac{2}{d}}\, d\lambda\right)^{\frac{1}{2}}
\\
\leq C\left(\frac{1}{|x_1-y_1|^{1-\frac{1}{d}}} + \check k_{\tau,0}(x_1-y_1)\right),
\\
\text{where } \check k_{\tau,0}(z_1) = 1_{|z_1| \,>\,1/\tau}
\tau^{\frac{1}{2} - \frac{1}{d}} e^{-\tau z_1^2/C} |z_1|^{-\frac{1}{2}}.
\end{multline*}
Easy computations then give $
\| \check k_{\tau,0}(z_1)\|_{L^{\frac{d}{d-1}}(\R)} \leq C \tau^{\frac{1}{4}-\frac{1}{2d}}, $ as announced. This concludes the proof of Lemma~\ref{Lem-Est-k0-x1-y1-integratednorms}.
\end{proof}\let\qed\relax
\end{proof}

Now we are in position to conclude the proof of Proposition~\ref{Prop-OpNorm-Est-K0}. First, from Proposition~\ref{Prop-Est-Fourier-Op} with $n = d-1$, estimate~\eqref{est-k0-x1-y1-Lp-Lp'} and the one-dimensional Hardy--Littlewood--Sobolev inequality (recall Theorem~\ref{Thm-Hardy-Littlewood-Sobolev}), we have, for $f \in L^{\frac{2d}{d+2}}(\Omega)$,
\begin{multline*}
\norm{ K_{\tau, 0} f }_{L^{\frac{2d}{d-2}}(\Omega)}\\
\begin{aligned}
&\leq
\norm{ \norm{ K_{\tau, 0} f(x_1, \cdot) }_{L_{x'}^{\frac{2d}{d-2}}(\R^{d-1})}}_{L_{x_1}^{\frac{2d}{d-2}}(X_0,X_1)}
\\
&\leq
\norm{ \int_{X_0}^{X_1} \left(\int_{\lambda} \| k_{\tau, 0}(x_1, y_1, \lambda,\cdot)\|_{L^\infty(\Sigma_{x_1}) } \lambda^{1 - \frac{2}{d}} \, d\lambda\right) \norm{ f(y_1, \cdot)}_{L_{y'}^{\frac{2d}{d+2}}(\R^{d-1}) } d y_1 }_{L_{x_1}^{\frac{2d}{d-2}}(X_0,X_1)}
\\
&\leq C
\norm{ \int_{X_0}^{X_1} \frac{1}{|x_1 - y_1|^{1-\frac{2}{d}}} \norm{ f(y_1, \cdot)}_{L_{y'}^{\frac{2d}{d+2}}(\R^{d-1}) } d y_1 }_{L_{x_1}^{\frac{2d}{d-2}}(X_0,X_1)}
\\
&\leq C \norm{ \norm{ f(y_1, \cdot)}_{L_{y'}^{\frac{2d}{d+2}}(\R^{d-1}) }}_{L_{y_1}^{\frac{2d}{d+2}}(X_0,X_1)} = C \norm{f}_{L^{\frac{2d}{d+2}}(\Omega)}.
\end{aligned}
\end{multline*}

Using the estimate~\eqref{est-k0-x1-y1-Lp-L2} in Proposition~\ref{Prop-Est-Fourier-Op}, and Young's inequality, we have, for $f \in L^{\frac{2d}{d+2}}(\Omega)$,
\begingroup
\allowdisplaybreaks
\begin{align*}
&\norm{K_{\tau, 0} f}_{L^{2}(\Omega)}\\
&\leq
\norm{ \norm{ K_{\tau, 0} f(x_1, \cdot) }_{L_{x'}^{2}\left(\R^{d-1}\right)}}_{L_{x_1}^{2}(X_0, X_1)}
\\
&
\leq
\norm{ \int_{X_0}^{X_1} \left(\int_{\lambda} \left\| k_{\tau, 0}(x_1, y_1, \lambda,\cdot)\right\|^2_{L^\infty(\Sigma_{x_1}) } \lambda^{1 - \frac{2}{d}} \, d\lambda\right)^{\frac{1}{2}} \norm{ f(y_1, \cdot)}_{L_{y'}^{\frac{2d}{d+2}}(\R^{d-1}) }d y_1}_{L_{x_1}^{2}(X_0,X_1)}
\\
&\leq C
\norm{ \int_{X_0}^{X_1} \tilde k_{\tau,0}(x_1 -y_1) \norm{ f(y_1, \cdot)}_{L_{y'}^{\frac{2d}{d+2}}\left(\R^{d-1}\right)}d y_1 }_{L_{x_1}^{2}(X_0,X_1)}
\\
& \leq C \norm{\tilde k_{\tau,0}}_{L^{\frac{d}{d-1}}(\R)} \norm{ \norm{ f(y_1, \cdot)}_{L_{y'}^{\frac{2d}{d+2}}\left(\R^{d-1}\right)}}_{L_{y_1}^{\frac{2d}{d+2}}(X_0,X_1)} \leq C\tau^{-\frac{3}{4}-\frac{1}{2d}} \norm{f}_{L^{\frac{2d}{d+2}}(\Omega)}.
\end{align*}
\endgroup
Similarly, for $f \in L^{2}(\Omega)$, we get
\begin{equation}\label{Est-K-tau-0-L2-L2_*'}
\norm{ K_{\tau, 0} f}_{L^{\frac{2d}{d-2}}(\Omega)}
\leq
\tau^{-\frac{3}{4}-\frac{1}{2d}} \norm{f}_{L^{2}(\Omega)}.
\end{equation}

Using the relation~\eqref{d1-Hat-K}, the estimate~\eqref{est-k0-d1-x1-y1-Lp-L2} in Proposition~\ref{Prop-Est-Fourier-Op}, and Hardy--Little\-wood--Sobolev theorem (Theorem~\ref{Thm-Hardy-Littlewood-Sobolev}), we get, for all $f \in L^{\frac{2d}{d+2}}(\Omega)$,
\begin{equation*}
\norm{\partial_1 \widehat{ K_{\tau,0} f}}_{L^2(\Omega_{1, \tau})}
\leq C \norm{f}_{L^{\frac{2d}{d+2}}(\Omega)}.
\end{equation*}

Using the estimate~\eqref{est-k0-x1-y1-Lp-nabla'-L2} in Proposition~\ref{Prop-Est-Fourier-Op}, Hardy--Littlewood--Sobolev theorem (Theorem~\ref{Thm-Hardy-Littlewood-Sobolev}) and Young's inequality, we get, for all $f \in L^{\frac{2d}{d+2}}(\Omega)$,
\begin{equation*}
\norm{\nabla' K_{\tau, 0} f}_{L^{2}(\Omega)}
\leq C{\tau^{\frac{1}{4}-\frac{1}{2d}}} \norm{f}_{L^{\frac{2d}{d+2}}(\Omega)}.
\end{equation*}
The above estimates allow to conclude the estimate~\eqref{Est-K-tau-0-f-L-2_*}.

Estimate~\eqref{Est-K-tau-0-f-L-2} simply consists in the combination of~\eqref{Est-K-tau-0-L2-L2_*'}, Theorem~\ref{Thm-Hilbertian-est}, and estimate~\eqref{Est-d1-Hilbertian-Est}.
\end{proof}



\subsection{Estimates on the operator \texorpdfstring{$K_{\tau, 1}$}{K tau, 1}} The goal of this section is to estimate the norm of the operator $K_{\tau, 1}$ in~\eqref{Def-K-tau-j-Fourier}--\eqref{Def-k-tau-1-x1-y1-xi'}, more precisely:
\begin{prop}\label{Prop-OpNorm-Est-K1}
Let $\Omega$ be as in~\eqref{Omega=Strip} with $X_0 \! <\! 0\! <\! X_1$ and $\max\{|X_0|, |X_1|\} \leq 1$, and assume that the coefficients $(\lambda_j)_{j\, \in\,\{1,\,\cdots,\,d\}} \in \R^d$ satisfy $\lambda_1 = 0$, \eqref{Coercivity} and~\eqref{Sub-ellipticity-Strip}. Then there exist $C>0$ and $\tau_0 \geq 1$ independent of $X_0, X_1$ (and depending only on $c_0$, $m_*$ and $M_*$ in~\eqref{Coercivity} and~\eqref{Sub-ellipticity-Strip}), such that for all $\tau \geq \tau_0$ and for all $f \in L^2(\Omega)$,
\begin{multline*}
\tau^{-\frac{1}{4}+\frac{1}{2d}}\norm{ K_{\tau,1} f}_{L^{\frac{2d}{d-2}}(\Omega)} +
\tau^{\frac{1}{2}} \norm{K_{\tau,1} f}_{L^{2}(\Omega)} +
\norm{\partial_1 \widehat{ K_{\tau,1} f}}_{L^2(\Omega_{1, \tau})} +
\tau^{-\frac{1}{2}} \norm{\nabla' K_{\tau,1} f}_{L^2(\Omega)} \\
\leq C \| f \|_{L^2(\Omega)}.
\end{multline*}
\end{prop}

\begin{proof}
We use the same notations as in the proof of the previous proposition. From the definition of $k_{\tau,1}$ in~\eqref{Def-k-tau-1-x1-y1-xi'}, adapting the proof of Lemma~\ref{Lem-Est-k0-x1-y1-integratednorms} to $k_{\tau,1}$, we can easily deduce the following result, whose proof is left to the reader:

\begin{lemm}\label{Lem-Est-k-1-x1-y1}
There exists a constant $C >0$ independent of $X_0, X_1$ (and depending only on $c_0$, $m_*$ and $M_*$ in~\eqref{Coercivity} and~\eqref{Sub-ellipticity-Strip}), such that for all $x_1$ and $y_1$ in $[X_0, X_1]$, for all $\tau \geq 1$, and $\lambda >0$,
\begin{multline}\label{Est-k-tau-1-x1-y1}
\| k_{\tau, 1}(x_1,y_1,\lambda, \cdot) \|_{L^\infty\left(\Sigma_{x_1}\right)}\\ 
\leq C e^{-|\tau - \lambda| | y_1 - x_1 | - \lambda (y_1 - x_1)^2/C }
+ C (\tau + \lambda) \| k_{\tau, 0}(x_1,y_1,\lambda, \cdot) \|_{L^\infty\left(\Sigma_{x_1}\right)}.
\end{multline}
\end{lemm}

We shall then prove the following lemma:
\begin{lemm}\label{Lem-Est-k1-x1-y1-integratednorms}
There exists a constant $C >0$ independent of $X_0, X_1$ (and depending only on $c_0$, $m_*$ and $M_*$ in~\eqref{Coercivity} and~\eqref{Sub-ellipticity-Strip}), such that for all $x_1$ and $y_1$ in $[X_0, X_1]$, for all $\tau \geq 1$,
\begin{multline}\label{est-k1-x1-y1-Lp-L2}
\left(\int_{\lambda\,>\,0} \norm{k_{\tau,1}(x_1, y_1,\lambda, \cdot) }^2_{L^\infty(\Sigma_{x_1})} \, \lambda^{1 -\frac{2}{d}}\, d\lambda\right)^{\frac{1}{2}}
\leq C\frac{1}{|x_1-y_1|^{1-\frac{1}{d}}} + \tilde k_{\tau,1}(x_1-y_1),
\\
 \text{ with }\tilde k_{\tau,1} \in L^{\frac{d}{d-1}} (\R)\text{ and }
\left\|\tilde k_{\tau,1}\right\|_{L^{\frac{d}{d-1}} (\R)} \leq C \tau^{\frac{1}{4}-\frac{1}{2d}}.
\end{multline}
\end{lemm}

\begin{proof}
Using Lemma~\ref{lem-general-inequalities}$\MK$\eqref{item-2-lemma-6.1}, we get
\begin{multline*}
\int_{0}^{\tau} e^{-2 (\tau - \lambda)|y_1 -x_1|- 2\lambda(y_1-x_1)^2/C}\lambda^{1 -\frac{2}{d}} \, d\lambda\\
\begin{aligned}
&\leq C e^{-2 \tau |y_1-x_1|}\int_{ 0}^{\tau} e^{2 \lambda\left(|y_1-x_1|- (y_1-x_1)^2/C\right)} \lambda^{1-\frac{2}{d}} \, d \lambda
\\
&\leq C |y_1 - x_1|^{-2 + \frac{2}{d}} e^{-2\tau |y_1-x_1|} \! \! \int_0^{\tau |y_1 -x_1|}  e^{2\lambda \left(1-|y_1-x_1|/C\right)} \lambda^{1-\frac{2}{d}} \, d\lambda
\\
& \leq C |y_1 - x_1|^{-2+\frac{2}{d}} e^{-2\tau(y_1-x_1)^2/C} \frac{(\tau |y_1-x_1|)^{2-\frac{2}{d}}}{1 + \tau |y_1 - x_1|}
\\
& \leq C \tau^{1 -\frac{2}{d}} |y_1-x_1|^{-1} e^{- 2 \tau (y_1 - x_1)^2/C}.
\end{aligned}
\end{multline*}
Similarly, using Lemma~\ref{lem-general-inequalities}$\MK$\eqref{item-3-lemma-6.1}, we get
\begin{multline*}
\int_{\tau}^\infty e^{-2 (\lambda -\tau)|y_1 -x_1|- 2\lambda(y_1-x_1)^2/C}\lambda^{1 -\frac{2}{d}}\, d\lambda\\
\begin{aligned}
&\leq \! C e^{2 \tau|y    _1-x_1|}\int_{\tau}^\infty \!\!\! e^{-2 \lambda\left(|y_1-x_1|+ (y_1-x_1)^2/C\right)} \lambda^{1-\frac{2}{d}} \, d \lambda
\\
&\leq C |y_1 - x_1|^{-2 + \frac{2}{d}} e^{-2\tau |y_1-x_1|}\int_{\tau |y_1 -x_1|}^\infty e^{-2\lambda \left(1+|y_1-x_1|/C\right)} \lambda^{1-\frac{2}{d}} \, d\lambda
\\
& \leq C |y_1 - x_1|^{-2 + \frac{2}{d}} e^{-2\tau(y_1-x_1)^2/C} (1 + \tau |y_1-x_1|)^{1 - \frac{2}{d}}
\\
& \leq C \left(1_{| y_1 - x_1 |\,\leq\,1/\tau } |y_1-x_1|^{-2+\frac{2}{d}} + 1_{| y_1 - x_1 | \,\geq\,1/\tau } \tau ^{1-\frac{2}{d}} | y_1 - x_1 |^{-1}e^{-2\tau(y_1-x_1)^2/C} \right).
\end{aligned}
\end{multline*}
Now, from~\eqref{est-k0-x1-y1-Lp-L2}, \eqref{est-k0-x1-y1-Lp-nabla'-L2}, \eqref{Est-k-tau-1-x1-y1} and the explicit formula of $\tilde k_{\tau,0}$ in~\eqref{def-tilde-k-tau-0}, we get
\begin{multline*}
\left(\int_{\lambda\,>\,0} (\tau + \lambda)^2 \left\| k_{\tau, 0} (x_1, y_1, \lambda) \right\|^2 \, \lambda^{1-\frac{2}{d}} d\lambda\right)^{\frac{1}{2}}
\\
\leq C \left(\frac{1}{|x_1 - y_1|^{1-\frac{1}{d}}} + \tau \tilde k_{\tau,0}(x_1 - y_1) + \check k_{\tau, 0} (x_1 - y_1) \right).
\end{multline*}
We then easily obtain~\eqref{est-k1-x1-y1-Lp-L2}.
\end{proof}

We now conclude the proof of Proposition~\ref{Prop-OpNorm-Est-K1}. As in the proof of Proposition~\ref{Prop-OpNorm-Est-K0}, the estimate~\eqref{est-k1-x1-y1-Lp-L2} easily implies that there exists $C>0$ such that for all $\tau \geq 1$ and for all $f \in L^2(\Omega)$,
\begin{equation*}
\tau^{-\frac{1}{4}+\frac{1}{2d}} \left\| K_{\tau,1} f\right\|_{L^{\frac{2d}{d-2}}(\Omega)}
\leq C \| f \|_{L^2(\Omega)}.
\end{equation*}
The estimates on $K_{\tau,1}$ and $\nabla' K_{\tau,1}$ as operators from $L^2(\Omega)$ to $L^2(\Omega)$, and on $\partial_1\widehat { K_{\tau,1}}$ from $L^2(\Omega)$ to $L^2(\Omega_{1, \tau})$ can then be deduced immediately from Theorem~\ref{Thm-Hilbertian-est} and~\eqref{Est-d1-Hilbertian-Est}.
\end{proof}

\subsection{Estimates on the operator \texorpdfstring{$K_{\tau,j}$}{K tau,j} for \texorpdfstring{$j \geq 2$}{j >= 2}}


The goal of this section is to estimate the norm of the operator $K_{\tau, j}$ in~\eqref{Def-K-tau-j-Fourier}--\eqref{Def-k-tau-j-x1-y1-xi'} for $j \geq 2$, more precisely:

\begin{prop}\label{Prop-OpNorm-Est-Kj}
Let $\Omega$ be as in~\eqref{Omega=Strip} with $X_0 \! <\! 0\! <\! X_1$ and $\max\{|X_0|, |X_1| \}\linebreak \leq 1$, and assume that the coefficients $(\lambda_j)_{j\,\in\,\{1,\,\cdots,\,d\}} \in \R^d$ satisfy $\lambda_1 = 0$, \eqref{Coercivity} and~\eqref{Sub-ellipticity-Strip}. Then there exist $C>0$ and $\tau_0 \geq 1$ independent of $X_0, X_1$ (and depending only on $c_0$, $m_*$ and $M_*$ in~\eqref{Coercivity} and~\eqref{Sub-ellipticity-Strip}), such that for all $j \in \{2,\,\cdots,\,d\}$, for all $\tau \geq \tau_0$ and for all $f \in L^2(\Omega)$,
\begin{multline*}
\tau^{-\frac{1}{4}+\frac{1}{2d}} \left\| K_{\tau,j} f\right\|_{L^{\frac{2d}{d-2}}(\Omega)} +
\tau^{\frac{1}{2}}\left\| K_{\tau,j} f\right\|_{L^{2}(\Omega)} +
\left\| \partial_1 \widehat{ K_{\tau,j} f}\right\|_{L^2(\Omega_{1, \tau})} \\ +
\tau^{-\frac{1}{2}} \left\| \nabla' K_{\tau,j} f\right\|_{L^2(\Omega)}
\leq C \| f \|_{L^2(\Omega)}.
\end{multline*}
\end{prop}

\begin{proof}
We start by noticing that there exists a constant $C >0$ such that for all $j \in \{2, \,\cdots,\,d\}$, for all $x_1$ and $y_1$ in $[X_0, X_1]$, for all $\tau \geq 1$, and $\lambda >0$,
\begin{equation*}\label{Est-k-tau-j-x1-y1}
\left\| k_{\tau, j}(x_1,y_1,\lambda, \cdot) \right\|_{L^\infty\left(\Sigma_{x_1}\right)}
\leq C \lambda \left\| k_{\tau, 0}(x_1,y_1,\lambda, \cdot) \right\|_{L^\infty\left(\Sigma_{x_1}\right)}.
\end{equation*}
From~\eqref{est-k0-x1-y1-Lp-nabla'-L2}, there exists a constant $C >0$ such that for all $x_1$ and $y_1$ in $[X_0, X_1]$, for all $\tau \geq 1$, and for all $j \in \{2, \cdots, d\}$,
\begin{multline}\label{est-kj-x1-y1-Lp-L2}
\left(\int_{\lambda\,>\,0} \norm{k_{\tau,j}(x_1, y_1,\lambda, \cdot) }^2_{L^\infty(\Sigma_{x_1})} \, \lambda^{1 -\frac{2}{d}}\, d\lambda\right)^{\frac{1}{2}}
\leq C\frac{1}{|x_1-y_1|^{1-\frac{1}{d}}} + \check k_{\tau,0}(x_1-y_1),
\\
\text{ with }\check k_{\tau,0} \in L^{\frac{d}{d-1}} (\R)\text{ and }
\left\|\check k_{\tau,1}\right\|_{L^{\frac{d}{d-1}} (\R)} \leq C \tau^{\frac{1}{4}-\frac{1}{2d}}.
\end{multline}
Accordingly, there exists $C>0$, such that for all $j \in \{2, \cdots, d\}$, for all $\tau \geq 1$ and for all $f \in L^2(\Omega)$,
\begin{equation*}
\tau^{-\frac{1}{4}+\frac{1}{2d}} \left\| K_{\tau,j} f\right\|_{L^{\frac{2d}{d-2}}(\Omega)}
\leq C \| f \|_{L^2(\Omega)}.
\end{equation*}
We then conclude Proposition~\ref{Prop-OpNorm-Est-Kj} by combining this estimate with the ones in Theorem~\ref{Thm-Hilbertian-est} and~\eqref{Est-d1-Hilbertian-Est}.
\end{proof}


\subsection{Estimates on the operator \texorpdfstring{$G_\tau$}{G tau} in~\eqref{Def-G-tau-Fourier}--\eqref{Def-g-tau-x1-xi'}}

\begin{prop}\label{Prop-OpNorm-Est-G-tau}
Let $\Omega$ be as in~\eqref{Omega=Strip} with $X_0 \! <\! 0\! <\! X_1$ and $\max\{ |X_0|, |X_1|\}\linebreak \leq 1$, and assume that the coefficients $(\lambda_j)_{j\, \in\,\{1,\,\cdots,\,d\}} \in \R^d$ satisfy $\lambda_1 = 0$, \eqref{Coercivity} and~\eqref{Sub-ellipticity-Strip}. Then there exists $C>0$ independent of $X_0, X_1$ (and depending only on $c_0$, $m_*$ and $M_*$ in~\eqref{Coercivity} and~\eqref{Sub-ellipticity-Strip}), such that for all $\tau \geq 1$, for all $g \in H^{\frac{1}{2}}(\R^{d-1})$,
\begin{multline}\label{Est-G-tau-g-H-1/2-L-2_*}
\tau^{\frac{3}{4}} \norm{G_{\tau} (g)}_{L^2(\Omega)} +
\norm{\widehat{\partial_1 G_\tau} (g)}_{L^2(\Omega_{1, \tau})} 
+
\tau^{-\frac{1}{4}} \norm{\nabla' G_\tau (g)}_{L^2(\Omega)}\\
+
\norm{G_{\tau} g}_{L^{\frac{2d}{d-2}}(\Omega)}
\leq C \| g \|_{H^{\frac{1}{2}}\left(\R^{d-1}\right)}.
\end{multline}
\end{prop}

\begin{proof}
Note that all the terms in~\eqref{Est-G-tau-g-H-1/2-L-2_*} involving Hilbertian norms have already been estimated in Theorem~\ref{Thm-Improved-Hilbert-Est-G-tau}, so only the estimate on $G_\tau$ in the 
\newline
$\scL(H^{\frac{1}{2}}(\R^{d-1}), L^{\frac{2d}{d-2}}(\Omega))$-norm remains to prove.

For $g \in H^{\frac{1}{2}}(\R^{d-1})$, we introduce $g_0 \in L^2(\R^{d-1})$ so that $\hat g_0 (\xi') = |\xi'|^{\frac{1}{2}} \hat g(\xi')$, and we notice that $G_\tau (g) = G_{\tau,0}(g_0)$, where $G_{\tau, 0}$ is given as follows:
\[
\widehat{G_{\tau,0} g_0} (x_1, \xi') = g_{\tau,0} (x_1, \xi') \widehat g_0(\xi'), \qquad \qquad(x_1, \xi') \in [X_0, X_1] \times \R^{d-1},
\]
with $ g_{\tau,0} (x_1, \xi') = {g_{\tau}(x_1, \xi') }/{|\xi'|^{\frac{1}{2}}}$, and $g_\tau$ as in~\eqref{Def-g-tau-x1-xi'}. It is then clear that the $\scL(H^{\frac{1}{2}}(\R^{d-1}), L^{\frac{2d}{d-2}}(\Omega))$-norm of $G_\tau$ coincides with the $\scL(L^2(\R^{d-1}), L^{\frac{2d}{d-2}}(\Omega))$-norm of $G_{\tau,0}$.

To estimate the $\scL(L^2(\R^{d-1}), L^{\frac{2d}{d-2}}(\Omega))$-norm of $G_{\tau,0}$, we compute $G_{\tau,0}^\star$, where the adjoint is given with respect to the $L^2(\Omega)$ scalar product: for $f \in L^2(\Omega)$,
\[
\widehat{G_{\tau,0}^\star f}(\xi') = \int_{x_1 \,\in\,[X_0, X_1]} g_{\tau,0} (x_1,\xi') \widehat f(x_1, \xi') \, dx_1, \qquad \qquad\xi' \in \R^{d-1}.
\]
Accordingly, for $f \in L^2(\Omega)$,
\[
\widehat{G_{\tau,0} G_{\tau,0}^\star f} (x_1, \xi') = \int_{y_1\,\in\,[X_0, X_1] }g_{\tau,0}\left(x_1, \xi'\right) g_{\tau,0}\left(y_1, \xi'\right) \widehat f\left(y_1, \xi'\right) \, dy_1,
\qquad (x_1, \xi') \in \Omega.
\]
Our next goal is to check that the operator $ G_{\tau,0} G_{\tau,0}^\star$ belongs to $\scL(L^{\frac{2d}{d+2}}(\Omega), L^{\frac{2d}{d-2}}(\Omega))$.

In order to do that, as previously, we first check that there exists a constant $C$ such that for all $x_1 \in [X_0, X_1]$, for all $\lambda \geq 0$,
\begin{equation}\label{Est-g-tau-0}
\norm{ g_{\tau,0}(x_1,\lambda, \cdot)}_{L^\infty\left(\Sigma_{x_1}\right)}
\leq 1_{\lambda > \tau} \frac{C}{\lambda^{\frac{1}{2}}} e^{-(\lambda - \tau)(x_1 - X_0) - \lambda (x_1 - X_0)^2/C}.
\end{equation}

Consequently, for $x_1 \in [X_0, X_1]$ and $y_1 \in [X_0, X_1]$,
\begin{multline*}
\norm{G_{\tau,0} G_{\tau,0}^\star f(x_1, \cdot)
}_{L^{\frac{2d}{d-2}}\left(\R^{d-1}\right)}
\\
\begin{aligned}
&\leq
\int_{y_1\, \in\,[X_0, X_1]} \int_\tau^\infty \norm{ g_{\tau,0}(x_1,\lambda, \cdot)}_{L^\infty\left(\Sigma_{x_1}\right)} \norm{g_{\tau,0}(y_1,\lambda, \cdot)}_{L^\infty(\Sigma_{x_1})} \lambda^{1 - \frac{2}{d}} d \lambda\\
& \mkern250mu  \| f(y_1, \cdot) \|_{L^{\frac{2d}{d+2}}\left(\R^{d-1}\right)} dy_1
\\
& \leq
\int_{y_1\,\in\,[X_0, X_1]} 
\int_\tau^\infty  e^{-(\lambda - \tau)\left((x_1 - X_0)+(y_1-X_0)\right) - \lambda \left((x_1 - X_0)^2+(y_1-X_0)^2\right)/C}\!\lambda^{ - \frac{2}{d}} d \lambda \\
&\mkern250mu \| f(y_1, \cdot) \|_{L^{\frac{2d}{d+2}}\left(\R^{d-1}\right)}  dy_1
\\
& \leq 
\int_{y_1\, \in\,[X_0, X_1]}\left((x_1 -X_0) + (y_1\! -\!X_0) \right)^{-1+\frac{2}{d}} \frac{e^{-\tau \left((x_1 - X_0)^2+(y_1-X_0)^2\right)/C}}{1 + (\tau (x_1 -X_0 + y_1 -X_0))^{\frac{2}{d}}}\\
&\mkern250mu \| f(y_1, \cdot) \|_{L^{\frac{2d}{d+2}}\left(\R^{d-1}\right)} dy_1,
\end{aligned}
\end{multline*}
where we have used Lemma~\ref{lem-general-inequalities} item~\eqref{item-3-lemma-6.1}. With straightforward bounds, we thus get
\[
\norm{G_{\tau,0} G_{\tau,0}^\star f(x_1, \cdot)
}_{L^{\frac{2d}{d-2}}\left(\R^{d-1}\right)}
\leq
\int_{y_1\, \in\,[X_0, X_1]}\left(y_1 -x_1 \right)^{-1+\frac{2}{d}} \| f(y_1, \cdot) \|_{L^{2d/(d+2}\left(\R^{d-1}\right)} \, dy_1,
\]
and then Hardy--Littlewood--Sobolev theorem (Theorem~\ref{Thm-Hardy-Littlewood-Sobolev}) implies that
\begin{equation*}
\left\| G_{\tau,0} G_{\tau,0}^\star f
\right\|_{L^{\frac{2d}{d-2}}(\Omega)}
\leq C \| f \|_{L^{\frac{2d}{d+2}}(\Omega)}.
\end{equation*}
It follows that the $\scL(L^2(\R^{d-1}), L^{\frac{2d}{d-2}}(\Omega))$-norm of $G_{\tau,0}$ is bounded by a constant independent of $\tau$ as announced, and thus this is also the case for the $\scL(H^{\frac{1}{2}}(\R^{d-1}),\newline L^{\frac{2d}{d-2}}(\Omega))$-norm of $G_\tau$.
\end{proof}
\begin{rema}
One may wonder why the above proof does not rely on the estimate~\eqref{Est-K-a-k-L2-Lp'} directly. This is due to the fact that it would correspond to a limit case. Indeed, from the estimate\eqref{Est-g-tau-0} and Proposition~\ref{Prop-Est-Fourier-Op}, for $x_1 \in [X_0, X_1]$,
\begin{multline*}
 \norm{{G_{\tau,0} g_0}(x_1, \cdot) }_{L^{\frac{2d}{d-2}}\left(\R^{d-1}\right)}\\
\begin{aligned}
&\leq C \| g_0 \|_{L^2\left(\R^{d-1}\right)}
\sqrt{\int_\tau^\infty \norm{ g_{\tau,0}(x_1,\lambda, \cdot)}^2_{L^\infty(\Sigma_{x_1})} \lambda^{1 - \frac{2}{d}} \, d\lambda}
\\
&
\leq C\| g_0 \|_{L^2\left(\R^{d-1}\right)}
\sqrt{\int_\tau^\infty e^{-2 (\lambda - \tau)(x_1 - X_0) - 2 \lambda (x_1 - X_0)^2/C}
\lambda^{- \frac{2}{d}} \, d\lambda}
\\
& \leq C \| g_0 \|_{L^2\left(\R^{d-1}\right)} |x_1 -X_0|^{-(d-2)/(2d)} \frac{e^{ - \tau (x_1 - X_0)^2/C}}{1+ (\tau (x_1 - X_0))^{\frac{1}{d}}},
\end{aligned}
\end{multline*}
where we used item~\eqref{item-3-lemma-6.1} in Lemma~\ref{lem-general-inequalities}. But the function
\[
x_1 \mapsto |x_1 -X_0|^{-(d-2)/(2d)} \frac{e^{ - \tau (x_1 - X_0)^2/C}}{1+ (\tau (x_1 - X_0))^{\frac{1}{d}}}
\]
does not belong to $L^{\frac{2d}{d-2}}(X_0, X_1)$, and so we cannot conclude directly that $G_{\tau, 0} g_0 $ belongs to $L^{\frac{2d}{d-2}}(\Omega)$.
\end{rema}

\subsection{Estimates on the operator \texorpdfstring{$R_\tau$}{R tau} in~\eqref{Def-R-tau-Fourier}}

For $f \in H^1(\Omega)$, using that $R_\tau (f) = K_{\tau,0} (g)$ with $ \hat g(x_1, \xi') = \partial_1 \psi(x_1, \xi') \hat f(x_1, \xi')$, and thus with $\| g \|_{L^2(\Omega)} \leq C \| \nabla' f \|_{L^2(\Omega)}$, we immediately deduce the following result from~\eqref{Est-K-tau-0-f-L-2}:

\begin{prop}\label{Prop-OpNorm-Est-R-tau}
Let $\Omega$ be as in~\eqref{Omega=Strip} with $X_0 \! <\! 0\! <\! X_1$ and $\max\{|X_0|, |X_1|\}\linebreak \leq 1$, and assume that the coefficients $(\lambda_j)_{j \in \{1,\cdots,d\}} \in \R^d$ satisfy $\lambda_1 = 0$, \eqref{Coercivity} and~\eqref{Sub-ellipticity-Strip}. Then there exist $C>0$ and $\tau_0 \geq 1$ independent of $X_0, X_1$ (and depending only on $c_0$, $m_*$ and $M_*$ in~\eqref{Coercivity} and~\eqref{Sub-ellipticity-Strip}), such that for all $\tau \geq \tau_0$, for all $f \in H^1(\Omega)$,
\begin{multline*}
\tau^{\frac{3}{4}+\frac{1}{2d}} \norm{R_{\tau} f}_{L^{\frac{2d}{d-2}}(\Omega)} +
\tau^{\frac{3}{2}} \norm{R_{\tau} f}_{L^{2}(\Omega)} +
\tau
\norm{\partial_1 \widehat{ R_{\tau} f}}_{L^2(\Omega_{1, \tau})} +
\tau^{\frac{1}{2}} \norm{\nabla' R_{\tau,0} f}_{L^2(\Omega)}\\
\leq C \norm{\nabla' f}_{L^2(\Omega)}.
\end{multline*}
\end{prop}

\subsection{Proofs of Theorem~\ref{Thm-Carleman-Strip-2} and Theorem~\ref{Thm-Carleman-Strip}}

Within the setting of Theorem~\ref{Thm-Carleman-Strip-2}, we use Proposition~\ref{Prop-parametrix-Strip-1} to write the solution $w$ of~\eqref{Elliptic-Strip-w} satisfying $(\partial_1 w -F_1)(X_1,\cdot) = 0$ in $\R^{d-1}$ under the form~\eqref{parametrix}. Since $w \in H^1(\Omega)$, using the various estimates in Propositions~\ref{Prop-OpNorm-Est-K0}, \ref{Prop-OpNorm-Est-K1}, \ref{Prop-OpNorm-Est-Kj}, \ref{Prop-OpNorm-Est-G-tau} and~\ref{Prop-OpNorm-Est-R-tau}, we obtain, on one hand,
\begin{multline}\label{Almost-1}
\tau^{\frac{3}{2}} \| w\|_{L^2(\Omega)} +
\tau^{\frac{3}{4}} \| \partial_1 w\|_{L^2(\Omega)} +
\tau^{\frac{1}{2}} \norm{\nabla' w}_{L^2(\Omega)}
\leq C\left(
\| f_2 \|_{L^2(\Omega)}
\right.
\\
\left. +
\tau^{\frac{3}{4}- \frac{1}{2d}} \norm{f_{2_*'}}_{L^{\frac{2d}{d+2}}(\Omega)} +
\tau \| F \|_{L^2(\Omega)}+
\tau^{\frac{3}{4}} \| g \|_{H^{\frac{1}{2}}
\left(\{X_0\} \times \R^{d-1}\right)} +
\norm{\nabla' w}_{L^2(\Omega)}
\right),
\end{multline}
and, on the other hand,
\begin{multline}\label{Almost-2}
\tau^{\frac{3}{4}+\frac{1}{2d}} \| w\|_{L^{\frac{2d}{d-2}}(\Omega)}
\leq C\left(
\norm{f_2}_{L^2(\Omega)} +
\tau^{\frac{3}{4}+\frac{1}{2d}} \norm{f_{2_*'}}_{L^{\frac{2d}{d+2}}(\Omega)} +
\tau \| F \|_{L^2(\Omega)}
\right.
\\
\left.+
\tau^{\frac{3}{4}+ \frac{1}{2d}} \| g \|_{H^{\frac{1}{2}}\left(\{X_0\} \times \R^{d-1}\right)} +
\norm{\nabla' w}_{L^2(\Omega)}
\right).
\end{multline}
We then simply take $\tau \geq \tau_0$ with $\tau_0 \geq 1$ large enough in order to absorb the last term in the right-hand side of~\eqref{Almost-1} by $\tau^{\frac{1}{2}} \|\nabla' w\|_{L^2(\Omega)}$, and we get
\begin{multline}\label{Almost-1-bis}
\tau^{\frac{3}{2}} \| w\|_{L^2(\Omega)} +
\tau^{\frac{3}{4}} \| \partial_1 w\|_{L^2(\Omega)} +
\tau^{\frac{1}{2}} \norm{\nabla' w}_{L^2(\Omega)}
\\
\leq C\left(
\| f_2 \|_{L^2(\Omega)} +
\tau^{\frac{3}{4}- \frac{1}{2d}} \norm{f_{2_*'}}_{L^{\frac{2d}{d+2}}(\Omega)} +
\tau \| F \|_{L^2(\Omega)} +
\tau^{\frac{3}{4}} \| g \|_{H^{\frac{1}{2}}
\left(\{X_0\} \times \R^{d-1}\right)}
\right),
\end{multline}
that is estimate~\eqref{Carleman-Strip-w-1}.

Therefore, using~\eqref{Almost-1-bis}, the last term in~\eqref{Almost-2} can be removed, thus yielding~\eqref{Carleman-Strip-w-2}.

In order to prove Theorem~\ref{Thm-Carleman-Strip}, we simply use the correspondence $w(x) = e^{\tau x_1} v(x)$ for $x \in \Omega$. This proves~\eqref{Carleman-Strip-1} and~\eqref{Carleman-Strip-2} for $\tau \geq \tau_0$. We then deduce~\eqref{Carleman-Strip-1} and~\eqref{Carleman-Strip-2} for any $\tau \geq 1$ by changing the constant if necessary through straightforward bounds on $x_1 \mapsto \exp(\tau x_1) $ for $\tau \in [1, \tau_0]$.



\section{General geometrical setting: proof of Theorem~\ref{Thm-Carleman-General}}\label{Sec-Proof-Carl-General}

Here, we provide a proof of Theorem~\ref{Thm-Carleman-General} using Fourier techniques as we did earlier, following the approach developed in~\cite{darde2023short}, and adapted to the case of source terms in $H^{-1}(\Omega)$ and in $L^{\frac{2d}{d+2}}(\Omega)$. This approach is based on a localization argument and a gluing argument, as it is usually done for Carleman estimates. The originality here is that we will localize the functions in balls of size $\tau^{-\frac{1}{3}}$, that is, depending on the Carleman parameter $\tau$. Doing that choice allows to somehow approximate the weight function $\varphi$ by its quadratic approximation, and to reduce the problem through a suitable change of variables to the case of a strip with linear coefficients as in Theorem~\ref{Thm-Carleman-Strip-2} (see Lemma~\ref{Lemma-Local-Carl-x-0} and its proof).

For $\tau \geq 1$, we introduce
\begin{equation*}
w = e^{ \tau \varphi}u,
\qquad
\tilde f_2 = e^{\tau \varphi} (f_2 - \tau \nabla \varphi \cdot F),
\qquad
\tilde f_{2_*'} = e^{\tau \varphi} f_{2_*'},
\qquad
\tilde F = e^{\tau \varphi} F,
\qquad
\tilde g = g,
\end{equation*}
so that the function $u$ solves~\eqref{Elliptic-Eq-general} if and only if $w$ solves
\begin{equation}\label{Eq-Elliptic-w}
\left\{
\begin{array}{ll}
\ds \Delta w - 2\tau \nabla \varphi \cdot \nabla w + \tau^2 |\nabla \varphi|^2 w - \tau \Delta \varphi w = \tilde f_2 + \tilde f_{2_*'} + \div\left(\tilde F\right), &\text{ in } \Omega,
\\
\ds w = \tilde g & \text{ on } \partial \Omega.
\end{array}
\right.
\end{equation}


\subsection{Local estimates}

Our first step is to introduce a local version of~\eqref{Eq-Elliptic-w}. Namely, for $x_0 \in \overline\Omega\setminus \omega$, we introduce $\eta_{x_0}(x)$ a cut-off function, which will be made more precise in~\eqref{Choice-Localization}, and set
\begin{equation*}
w_{x_0}(x) = \eta_{x_0}(x) w(x), \qquad \qquad x \in \Omega,
\end{equation*}
which solves
\begin{equation}\label{Eq-Elliptic-w-x0}
\left\{
\begin{array}{ll}
\ds \Delta w_{x_0} - 2\tau \nabla \varphi \cdot \nabla w_{x_0} + \tau^2 |\nabla \varphi|^2 w_{x_0} = f_{2,x_0} + f_{2_*',x_0} + \div(F_{x_0}), & \text{ in } \Omega,
\\
\ds w_{x_0} = \tilde g_{x_0} & \text{ on } \partial \Omega,
\end{array}
\right.
\end{equation}
where
\begin{align*}
& f_{2,x_0} = \eta_{x_0} \tilde f_2 -\nabla \eta_{x_0} \cdot \tilde F + \tau \Delta \varphi w_{x_0} +2 \nabla \eta_{x_0} \cdot \nabla w + \Delta \eta_{x_0} w - 2 \tau \nabla \varphi \cdot \nabla \eta_{x_0} w,
\\
& f_{2_*',x_0} = \eta_{x_0} \tilde f_{2_*'},
\qquad F_{x_0} = \eta_{x_0} \tilde F,
\qquad g_{x_0} = \eta_{x_0} \tilde g.
\end{align*}

We claim that, provided the localization is strong enough, we can get a local Carleman estimate:
\begin{lemm}\label{Lemma-Local-Carl-x-0}
There exist constants $C>0$ and $\tau_0 \geq 1$ (depending only on $\alpha$, $\beta$, $\| \varphi \|_{C^3(\overline\Omega)}$) such that for all $\tau \geq \tau_0$, for all $x_0 \in \overline\Omega \setminus \omega$, and for all
\begin{equation*}
f_{2,x_0} \in L^2(\Omega),
\quad f_{2_*',x_0} \in L^{\frac{2d}{d+2}}(\Omega),
\quad F_{x_0} \in L^2(\Omega; \C^{d}),
\quad\text{ and } \quad g_{x_0} \in H^{\frac{1}{2}}(\partial\Omega),
\end{equation*}
and $w_{x_0}$ satisfying~\eqref{Eq-Elliptic-w-x0} and supported in $B(x_0, \tau^{-\frac{1}{3}}) \cap \overline{\Omega}$, we have
\begin{multline}\label{Carleman-Local-w-x0-L2}
\tau^{\frac{3}{2}} \| w_{x_0} \|_{L^2(\Omega)} +
\tau^{\frac{1}{2}} \| \nabla w_{x_0} \|_{L^2(\Omega)} +
\tau^{\frac{3}{4}- \frac{1}{2d}} \| w_{x_0}\|_{L^{\frac{2d}{d-2}}(\Omega)}
\\
\leq C\left(
\left\| f_{2,x_0} \right\|_{L^2(\Omega)} +
\tau^{\frac{3}{4}- \frac{1}{2d}}
\left\| f_{2_*',x_0} \right\|_{L^{\frac{2d}{d+2}}(\Omega)} +
\tau \| F_{x_0} \|_{L^2(\Omega)} +
\tau^{\frac{3}{4}} \| g_{x_0} \|_{H^{\frac{1}{2}}(\partial\Omega)}
\right),
\end{multline}
and
\begin{multline}\label{Carleman-Local-w-x0-L2_*}
\tau^{\frac{3}{4}+ \frac{1}{2d}} \left\| w_{x_0} \right\|_{L^{\frac{2d}{d-2}}(\Omega)} +
\tau^{\frac{3}{2}} \left\| w_{x_0} \right\|_{L^2(\Omega)} +
\tau^{\frac{1}{2}} \left\| \nabla w_{x_0}\right\|_{L^2(\Omega)}
\\
\leq C\left(
\left\| f_{2,x_0}\right |_{L^2(\Omega)} +
\tau^{\frac{3}{4}+ \frac{1}{2d}} \left\| f_{2_*',x_0} \right\|_{L^{\frac{2d}{d+2}}(\Omega)} +
\tau \| F_{x_0} \|_{L^2(\Omega)} +
\tau^{\frac{3}{4}+ \frac{1}{2d}} \| g_{x_0} \|_{H^{\frac{1}{2}}(\partial\Omega)}
\right).
\end{multline}
\end{lemm}

The proof of Lemma~\ref{Lemma-Local-Carl-x-0} is postponed to the next section, and will be based on a suitable change of variables and the Carleman estimate in Theorem~\ref{Thm-Carleman-Strip-2}.

In the following, we will thus choose the localization as follows:
\begin{equation}\label{Choice-Localization}
\eta_{x_0}(x) = \eta\left(\tau^{\frac{1}{3}} (x-x_0)\right), \qquad \qquad x \in \R^{d},
\end{equation}
where $\eta$ is a non-negative smooth radial function (in $\scC^\infty_c(\R^d)$) such that $\eta(\rho) = 1$ for $|\rho| \leq \frac{1}{2}$ and vanishing outside the unit ball, so that Lemma~\ref{Lemma-Local-Carl-x-0} applies to $w_{x_0}$, and the estimates~\eqref{Carleman-Local-w-x0-L2} and~\eqref{Carleman-Local-w-x0-L2_*} hold uniformly with respect to $x_0 \in \overline\Omega \setminus \omega$.


It thus remains to prove Lemma~\ref{Lemma-Local-Carl-x-0}, which is done in Section~\ref{Subsec-Proof-Local-Carl}, and to show how to glue the estimates in Lemma~\ref{Lemma-Local-Carl-x-0} to conclude Theorem~\ref{Thm-Carleman-General}, which is explained in Section~\ref{Sec-Gluing}.

\subsection{Proof of Lemma~\ref{Lemma-Local-Carl-x-0}: A suitable change of coordinates}\label{Subsec-Proof-Local-Carl}

The proof of Lemma~\ref{Lemma-Local-Carl-x-0} mainly reduces to a suitable change of variables allowing to link the Carleman estimates~\eqref{Carleman-Local-w-x0-L2} and~\eqref{Carleman-Local-w-x0-L2_*} in small balls around $x_0$ with the Carleman estimates~\eqref{Carleman-Strip-w-1}-\eqref{Carleman-Strip-w-2} in the strip proved in Theorem~\ref{Thm-Carleman-Strip-2}.

We let $x_0 \in \overline\Omega\setminus \omega$, and we introduce $L_1 \in \R^d$ and $A_1 \in \R^{d\times d}$ as follows:
\begin{equation*}
L_1 = \nabla \varphi(x_0) \in \R^d, \quad A_1 = \Hess \varphi(x_0) \in \R^{d\times d}.
\end{equation*}

The bilinear form
\[
\xi \in \R^d \mapsto (\Hess \varphi(x_0)) \xi \cdot \xi
\]
is symmetric on $\R^d$ and on $\Span \{L_1\}^\perp$. Accordingly, there exists a family of orthogonal vectors $(L_j)_{j\,\in\, \{2,\,\cdots,\,d\}} $ of $\Span \{L_1\}^\perp$ which diagonalizes this form, that we normalize so that for all $j \in \{2, \cdots, d\}$, $|L_j | = |L_1|$. Since the family $(L_j)_{j\, \in\,\{2,\, \cdots,\, d\} } $ of $\Span \{L_1\}^\perp$ diagonalizes the form $\xi \mapsto (\Hess \varphi(x_0)) \xi \cdot \xi$ in $\Span \{L_1\}^\perp$, for all $j \in \{2, \cdots, d\}$, there exist $\alpha_j $ and $\mu_j$ in $\R$ such that
\[
(\Hess \varphi(x_0)) L_j = \mu_j L_j + \alpha_j L_1, \qquad \qquad j \in \{2, \cdots, d\}.
\]
Note that by symmetry of $\Hess \varphi(x_0)$, we then necessarily have
\[
(\Hess \varphi(x_0)) L_1 = \mu_1 L_1 + \sum_{k\,\geq\,2} \alpha_k L_k,
\]
where
\[
\mu_1 = \frac{1}{|L_1|^2} (\Hess \varphi(x_0)) L_1 \cdot L_1 =
\frac{1}{|\nabla \varphi(x_0)|^2} (\Hess \varphi(x_0)) \nabla \varphi(x_0) \cdot \nabla \varphi(x_0).
\]
For $j \in \{2, \cdots, d\}$, we then introduce the self-adjoint matrix $A_j \in \R^{d \times d}$ defined by
\begin{equation}\label{Choices-A-j}
\left\{
\begin{array}{ll}
& \ds A_j L_1 = - \alpha_j L_1 - \mu_j L_j,
\\
& \ds A_j L_k = \alpha_k L_j - \alpha_j L_k, \quad \text{ if } k \in \{2,\, \cdots,\, d\}\setminus \{ j\},
\\
& \ds A_j L_j = - \mu_j L_1 + \sum_{k\,\geq\, 2} \alpha_k L_k.
\end{array}
\right.
\end{equation}
(It is easy to check that the matrix $A_j$ defined that way is indeed symmetric.)

We shall then introduce the following change of coordinates for $x$ in a neighbourhood of $x_0$:
\begin{align*}
& y_1(x) = \varphi(x)- \varphi(x_0),
\\
& \text{for} \, j \in \{2, \cdots, d\},
\, y_j(x) = L_j \cdot (x-x_0) + \frac{1}{2} A_j (x-x_0) \cdot (x-x_0).
\end{align*}
By construction, there exists a neighbourhood, whose size depends on the $C^2$ norm of $\varphi$ only, such that $x \mapsto y(x)$ is a local diffeomorphism between a neighbourhood $\clV$ of $x_0$ in $\overline\Omega\setminus \omega$ and a neighbourhood of $0$, that we call $\Omega_y$, and which may thus contain the image of a part of $\partial \Omega$. If it exists, we will denote it by $\Gamma_y = y(\partial \Omega \cap \clV)$.

For $\tau$ large enough, we can ensure that the ball of center $x_0$ and radius $\tau^{-\frac{1}{3}}$, when intersected with $\overline\Omega$, is included in a set on which $x \mapsto y(x)$ is a diffeomorphism, and its image is included in a ball $B(0, C \tau^{-\frac{1}{3}})$.

Therefore, for $w_{x_0}$ solving~\eqref{Eq-Elliptic-w-x0}, we set
\begin{equation*}
\check w(y) = w_{x_0}(x) \quad \text{ for } \quad y = y(x),
\end{equation*}
Explicit computations then give that $\check w$ satisfies, for $x \in \Omega \cap B(x_0, \tau^{-\frac{1}{3}}), $
\begin{multline}\label{Eq-Tilde-w}
 \sum_{j,k = 1}^d b_{j,k}(x) \partial_{y_j} \partial_{y_k} \check w(y(x)) + \nabla_y \check w(y(x)) \cdot \Delta_x y(x) - 2 \tau \sum_{j = 1}^d c_j (x) \partial_{y_j} \check w(y(x))
\\
+
\tau^2 |\nabla \varphi(x)|^2 \check w (y(x)) =
\left(\Delta w_{x_0} - 2\tau \nabla \varphi \cdot \nabla w_{x_0} + \tau^2 |\nabla \varphi|^2 w_{x_0} \right)(x),
\end{multline}
where
\[
b_{j,k}(x) = \nabla_x y_j(x) \cdot \nabla_x y_k(x), \quad \text{ and } \quad c_j (x) = \sum_{i = 1}^d \partial_i \varphi(x) \partial_i y_j(x).
\]
We then remark that $c_j(x) = b_{j,1}(x)$ and that $b_{j,k}(x) = b_{k,j}(x)$ for all $x$. We now briefly analyze the coefficients $b_{j,k}$. By construction of the coordinates $(y_j)_{j \,\in\,\{1,\, \cdots,\, d\}}$, we easily check that for $(j,k,\ell) \in \{1,\, \cdots,\, d\}^2$,
\begin{align*}
b_{j,k}(x_0) &= |L_1|^2 \delta_{j,k},
\\
\partial_\ell b_{j,k} (x_0) & = \sum_{i = 1}^d (\partial_\ell \partial_i y_j (x_0) \partial_i y_k(x_0) + \partial_i y_j(x_0) \partial_\ell \partial_i y_k(x_0))
\notag
\\
& = \sum_{i = 1}^d (A_j e_i \cdot e_\ell) (L_k \cdot e_i)+ (L_j \cdot e_i) (A_k e_i \cdot e_\ell) = \left(A_j L_k + A_k L_j\right) \cdot e_\ell,
\end{align*}
so that we have in particular that
\[
L_\ell \cdot \nabla b_{j,k} (x_0)= \left(A_j L_k + A_k L_j\right) \cdot L_\ell,
\]
For convenience, we also write separately
\[
|\nabla \varphi(x_0)|^2 = |L_1|^2,
\quad
\text{ and }
\quad
\partial_\ell \left(|\nabla \varphi|^2\right)(x_0) = 2 \sum_{i} \partial_\ell \partial_i \varphi(x_0) \partial_i \varphi(x_0) = 2 A_1 L_1 \cdot e_\ell.
\]
We can thus analyze $b_{j,k} / |\nabla \varphi|^2$ close to $x = x_0$:
\begin{align*}
\frac{b_{j,k}}{|\nabla \varphi|^2} (x_0) & = \delta_{j,k},
\\
L_\ell \cdot \nabla \left(\frac{b_{j,k}}{|\nabla \varphi|^2} \right) (x_0) & =
\frac{1}{|L_1|^2} \left(\left(A_j L_k + A_k L_j\right) \cdot L_\ell -2 \delta_{j,k} A_1 L_1 \cdot L_\ell \right).
\end{align*}
In particular, since $A_j L_k + A_k L_j = 0$ for all $j, k \in \{1,\, \cdots,\, d\}$ with $j \neq k$,
\[
\forall j, k \in \{1,\, \cdots,\, d\} \text{ with }j \neq k, \,
\forall \ell \in \{1,\, \cdots,\, d\}, \quad L_\ell \cdot \nabla \left(\frac{b_{j,k}}{|\nabla \varphi|^2} \right) (x_0) = 0.
\]
When $j = k = 1$, it is obvious that
\[
\forall \ell \in \{1,\, \cdots,\, d\}, \quad L_\ell \cdot \nabla \left(\frac{b_{1,1}}{|\nabla \varphi|^2} \right) (x_0) = 0.
\]
When $j = k \geq 2$, the choices~\eqref{Choices-A-j} yield
\begin{align*}
L_\ell \cdot \nabla \left(\frac{b_{j,j}}{|\nabla \varphi|^2} \right) (x_0) &= 0 \text{ when } \ell \in \{2,\, \cdots,\, d \},
\\
\intertext{and}
L_1 \cdot \nabla \left(\frac{b_{j,j}}{|\nabla \varphi|^2} \right) (x_0) &= -\frac{2}{|L_1|^2} \left(A_1 L_j \cdot L_j + A_1 L_1 \cdot L_1\right).
\end{align*}

Consequently, as a consequence of Taylor expansion of $b_{j,k}/|\nabla \varphi|^2$ close to $x = x_0$,
\[
\frac{b_{j,k}(x)}{|\nabla \varphi(x)|^2} =
\delta_{j,k}\left(1 - \lambda_j y_1(x)\right)+ \clO\left(|x-x_0|^2\right),
\]
where
\begin{equation}\label{Link-Lambda-j-Pseudo}
\lambda_1 = 0, \quad \lambda_j = \frac{2}{|L_1|^2} (A_1 L_j \cdot L_j + A_1 L_1 \cdot L_1) \quad \text{ for } j \in \{2,\, \cdots,\, d\}.
\end{equation}
Accordingly, setting
\[
\check b_{j,k}(y) = \frac{b_{j,k}(x)}{|\nabla \varphi(x)|^2} \quad \text{ for } y = y(x),
\]
we have
\[
\check b_{j,k}(y) = \delta_{j,k} (1- \lambda_j y_1) + \clO(|y|^2)
\quad \text{ and } \quad
\nabla_y \left(\check b_{j,k}(y)- \delta_{j,k} \left(1-\lambda_j y_1\right)\right) = \clO(|y|).
\]
Thus, using that $\check w$ is supported in $B(0, C \tau^{-\frac{1}{3}})$, writing
\begin{multline*}
\sum_{j,k = 1}^d \check b_{j,k}(y) \partial_{y_j} \partial_{y_k} \check w(y) -
\sum_{j = 1}^d (1 - \lambda_j y_1) \partial_{y_j}^2 \check w
\\
\begin{aligned}
=
&\sum_{j,k = 1}^d \partial_{y_j} \left(\left(\check b_{j,k}(y) - \delta_{j,k} (1 - \lambda_j y_1) \right) \partial_{y_{k}} \check w \right) \\
-
&\sum_{j,k = 1}^d \partial_{y_j} \left(\check b_{j,k}(y)- \delta_{j,k} \left(1-\lambda_j y_1\right)\right) \partial_{y_{k}} \check w,
\end{aligned}
\end{multline*}
we get that
\[
\sum_{j,k = 1}^d \check b_{j,k}(y) \partial_{y_j} \partial_{y_k} \check w(y) -
\sum_{j = 1}^d \left(1 - \lambda_j y_1\right) \partial_{y_j}^2 \check w =
\check f_{2,a} + \div\left(\check F_a\right),
\]
where
\begin{align*}
\check f_{2,a} (y) &= -
\sum_{j,k = 1}^d \partial_{y_j} \left(\check b_{j,k}(y)- \delta_{j,k} \left(1-\lambda_j y_1\right)\right) \partial_{y_{k}} \check w,
\\
\check F_{a,j}(y) &= \sum_{k = 1}^d \left(\check b_{j,k}(y) - \delta_{j,k} \left(1 - \lambda_j y_1\right) \right) \partial_{y_{k}} \check w
\end{align*}
satisfy
\begin{equation*}
\left\| \check f_{2,a} \right\|_{L^2(\Omega_y)} \leq C \tau^{-\frac{1}{3}} \left\| \nabla_y \check w\right\|_{L^2(\Omega_y)},
\quad \text{and} \quad
\left\| \check F_a\right\|_{L^2(\Omega_y)} \leq C \tau^{-\frac{2}{3}} \left\| \nabla_y \check w\right\|_{L^2(\Omega_y)}.
\end{equation*}
Similarly,
\[
-2 \tau \sum_{j = 1}^d \check b_{1,j}(y) \partial_{y_j} \check w(y) + 2 \tau \partial_{y_1} \check w(y) =
\check f_{2,b}(y),
\]
with
\begin{equation*}
\left\| \check f_{2,b} \right\|_{L^2(\Omega_y)} \leq C \tau^{\frac{1}{3}} \left\| \nabla_y \check w\right\|_{L^2(\Omega_y)}.
\end{equation*}
Finally,
\[
\check f_{2,c} (y) = -\frac{1}{|\nabla \varphi(x(y))|^2} \nabla_y \check w(y) \cdot \Delta_x y(x(y)),
\]
where $x(y)$ denotes the inverse of the change of variables $x \mapsto y(x)$, also satisfies
\begin{equation*}
\left\| \check f_{2,c} \right\|_{L^2(\Omega_y)} \leq C \left\| \nabla_y \check w\right\|_{L^2(\Omega_y)}.
\end{equation*}
We then set
\[
\rho_{j,k}(y) = \frac{\partial_{x_k} y_j(x(y))}{|\nabla \varphi(x(y))|^2},
\]
and introduce
\begin{align*}
&
\check f_{2}(y) = \frac{1}{|\nabla \varphi(x(y))|^2} f_{2,x_0}(x(y)) - \sum_{j,k} \partial_{y_j} \rho_{k,j} F_{x_0,k}(x(y)) + \check f_{2,a}(y) + \check f_{2,b}(y) + \check f_{2,c} (y),
\\
&\check f_{2_*'}(y) = \frac{1}{|\nabla \varphi(x(y))|^2} f_{2_*',x_0}(x(y)),
\\
& \check F_j (y) = \sum_{k = 1}^d \rho_{k,j}(y) F_{x_0,k}(x(y))+ \check F_{j,a}(y), \qquad \qquad j \in \{1, \cdots, d\},
\end{align*}
and we get that $\check w$ satisfies in $\Omega_y$,
\[
\sum_{j = 1}^d (1 - \lambda_j y_1) \partial_{y_j}^2 \check w - 2 \tau \partial_{y_1} \check w +
\tau^2 \check w =
\check f_{2} +
\check f_{2_*'} +
\text{div}_y\, \check F.
\]

If $\Gamma_y$ is not empty, then we simply recall that the weight function $\varphi$ has been chosen such that $\varphi = 0$ on the boundary of $\partial \Omega$. In particular, the set $\Gamma_y$ is simply parametrized by $ y_1(x) = Y_0$ for some $Y_0 = - \varphi(x_0) \leq 0$, and $\Omega$ can be locally defined by $y_1 > Y_{0}$. Thus, in this case, the equation of $\check w$ should be completed with
\begin{align*}
\check w(Y_{0}, y') &= \check g(y'),  &&\text{ for } y' \in \R^{d-1},
\\
\intertext{where}
\check g(y') &= g_{x_0}(x(Y_0, y')),  &&\text{ for } y' \in \R^{d-1} \text{ such that } (Y_0, y') \in \Gamma_y.
\end{align*}

Due to the form of $\Omega_y$ and the fact that we are considering functions which are supported in sets included in balls of the form $B(0, C \tau^{-\frac{1}{3}})$, we can then simply extend all the source terms in a strip of the form $[Y_0, Y_1] \times \R^{d-1}$, where the functions are extended by $0$ outside $\Omega_y$, and $\check w$ then satisfies
\begin{equation}\label{Elliptic-Strip-check w}
\left\{
\begin{array}{ll}
\Delta_y \check w-y_{1} \sum\limits_{j=2}^{d} \lambda_{j} \partial_{y_j}^{2} \check w - 2 \tau \partial_{y_1} \check w + \tau^2 \check w \\
\mkern200mu=- \check f_2 + \check f_{2_*'} +\divv_y\, \check F & \text { in } (Y_0, Y_1) \times \R^{d-1},
\\
\check w\left(Y_{0}, y'\right)=\check g (y'), &\text{ for } y' \in \R^{d-1},
\\
\check w\left(Y_{1}, y'\right)= 0, &\text{ for } y' \in \R^{d-1},
\\
\partial_{y_1} \check w\left(Y_{1}, y'\right)= 0, &\text{ for } y' \in \R^{d-1}.
\end{array}
\right.
\end{equation}
In the following, for convenience, we also write $\Omega_y$ for the strip $(Y_0, Y_1) \times \R^{d-1}$. Now, we come back to the definition of $\lambda_j$ in~\eqref{Link-Lambda-j-Pseudo} and remark that the condition~\eqref{CarlemanWeight-Cond2}, when taken at $x = x_0$, is equivalent to the condition~\eqref{Sub-ellipticity-Strip}. Accordingly, the Carleman estimates in Theorem~\ref{Thm-Carleman-Strip-2} apply: for $\tau \geq \tau_0$, we have
\begin{multline*}
\tau^{\frac{3}{2}} \norm{\check w}_{L^2\left(\Omega_y\right)} +
\tau^{\frac{1}{2}} \norm{\nabla \check w}_{L^2\left(\Omega_y\right)} +
\tau^{\frac{3}{4}- \frac{1}{2d}} \norm{\check w}_{L^{\frac{2d}{d-2}}\left(\Omega_y\right)}
\\
\leq C\left(
\norm{\check f_2}_{L^2\left(\Omega_y\right)} +
\tau^{\frac{3}{4}- \frac{1}{2d}} \norm{\check f_{2_*'}}_{L^{\frac{2d}{d+2}}\left(\Omega_y\right)} +
\tau \norm{\check F}_{L^2\left(\Omega_y\right)} +
\tau^{\frac{3}{4}} \norm{\check g}_{H^{\frac{1}{2}}\left(\{Y_0\} \times \R^{d-1}\right)}
\right),
\end{multline*}
and
\begin{multline*}
\tau^{\frac{3}{4}+ \frac{1}{2d}} \norm{\check w}_{L^{\frac{2d}{d-2}}(\Omega_y)} +
\tau^{\frac{3}{2}} \norm{\check w }_{L^2(\Omega_y)} +
\tau^{\frac{1}{2}} \norm{ \nabla \check w}_{L^2(\Omega_y)}
\\
\leq C\left(
\norm{\check f_2}_{L^2(\Omega_y)} +
\tau^{\frac{3}{4}+ \frac{1}{2d}} \norm{ \check f_{2_*'} }_{L^{\frac{2d}{d+2}}(\Omega_y)} +
\tau \norm{ \check F }_{L^2(\Omega_y)} +
\tau^{\frac{3}{4}+\frac{1}{2d}} \norm{ \check g }_{H^{\frac{1}{2}}\left(\{Y_0\} \times \R^{d-1}\right)}
\right).
\end{multline*}

We then simply remark that, from the expression of $\check f_2, \check f_{2_*'}, \check F$ and $\check g$,
\begin{multline*}
\norm{\check f_2 }_{L^2(\Omega_y)} +
\tau^{\frac{3}{4} -\frac{1}{2d}} \norm{\check f_{2_*'}}_{L^{\frac{2d}{d+2}}(\Omega_y)} +
\tau \norm{\check F}_{L^2(\Omega_y)} +
\tau^{\frac{3}{4}} \norm{\check g}_{H^{\frac{1}{2}}\left(\{Y_0\} \times \R^{d-1}\right)}
\\
\leq C
\left(
\norm{ f_{2,x_0}}_{L^2(\Omega)} +
\tau^{\frac{3}{4}-\frac{1}{2d}} \norm{f_{2_*',x_0}}_{L^{\frac{2d}{d+2}}(\Omega)} +
\tau \norm{F_{x_0}}_{L^2(\Omega)} +
\tau^{\frac{3}{4}} \norm{g_{x_0}}_{H^{\frac{1}{2}}(\partial \Omega)}
\right.\\
\left. +
\tau^{\frac{1}{3}} \norm{\nabla_y \tilde w}_{L^2(\Omega_y)}\right),
\end{multline*}
and
\begin{multline*}
\norm{ \check f_2}_{L^2(\Omega_y)} +
\tau^{\frac{3}{4}+\frac{1}{2d}} \norm{ \check f_{2_*'} }_{L^{\frac{2d}{d+2}}(\Omega_y)} +
\tau \norm{\check F }_{L^2(\Omega_y)} +
\tau^{\frac{3}{4}+\frac{1}{2d}} \norm{\check g}_{H^{\frac{1}{2}}\left(\{Y_0\} \times \R^{d-1}\right)}\\
\leq C
\left(\norm{f_{2,x_0}}_{L^2(\Omega)} +
\tau^{\frac{3}{4}+\frac{1}{2d}} \norm{ f_{2_*',x_0}}_{L^{\frac{2d}{d+2}}(\Omega)} +
\tau \norm{F_{x_0}}_{L^2(\Omega)} +
\tau^{\frac{3}{4}+\frac{1}{2d}} \norm{g_{x_0}}_{H^{\frac{1}{2}}(\partial \Omega)} \right.\\ \left.+
\tau^{\frac{1}{3}} \norm{\nabla_y \tilde w}_{L^2(\Omega_y)}\right).
\end{multline*}
Accordingly, taking $\tau_0\geq 1$ larger if necessary, we get for all $\tau \geq \tau_0$,
\begin{multline*}
\tau^{\frac{3}{2}} \norm{\check w}_{L^2(\Omega_y)} +
\tau^{\frac{1}{2}} \norm{\nabla \check w}_{L^2(\Omega_y)} +
\tau^{\frac{3}{4}- \frac{1}{2d}} \norm{\check w }_{L^{\frac{2d}{d-2}}(\Omega_y)}
\\
\leq C\left(
\norm{ \check f_2 }_{L^2(\Omega_y)} +
\tau^{\frac{3}{4}- \frac{1}{2d}} \norm{\check f_{2_*'} }_{L^{\frac{2d}{d+2}}(\Omega_y)} +
\tau \norm{\check F}_{L^2(\Omega_y)} +
\tau^{\frac{3}{4}} \norm{\check g}_{H^{\frac{1}{2}}\left(\{Y_0\} \times \R^{d-1}\right)}
\right),
\end{multline*}
and
\begin{multline*}
\tau^{\frac{3}{4}+ \frac{1}{2d}} \norm{\check w}_{L^{\frac{2d}{d-2}}(\Omega_y)} +
\tau^{\frac{3}{2}} \norm{\check w}_{L^2(\Omega_y)} +
\tau^{\frac{1}{2}} \norm{\nabla \check w }_{L^2(\Omega_y)}
\\
\leq C\left(
\norm{\check f_2}_{L^2(\Omega_y)} +
\tau^{\frac{3}{4}+ \frac{1}{2d}} \norm{\check f_{2_*'}}_{L^{\frac{2d}{d+2}}(\Omega_y)} +
\tau \norm{\check F}_{L^2(\Omega_y)} +
\tau^{\frac{3}{4}+\frac{1}{2d}} \norm{\check g}_{H^{\frac{1}{2}}\left(\{Y_0\} \times \R^{d-1}\right)}
\right).
\end{multline*}
Undoing the change of variables on the left-hand side, we easily deduce the estimates~\eqref{Carleman-Local-w-x0-L2} and~\eqref{Carleman-Local-w-x0-L2_*}.

The fact that the constants above do not depend on $x_0 \in \overline\Omega \setminus \omega$ can be tracked in the above proof: it comes from uniformity properties of the diffeomorphism $x \mapsto y$, and relies heavily on the uniform bounds~\eqref{CarlemanWeight-Cond1}--\eqref{CarlemanWeight-Cond2}, on the fact that $\varphi \in C^3(\overline\Omega)$, and that the constants in Theorem~\ref{Thm-Carleman-Strip-2} depend only on $c_0$, $m_*$ and $M_*$ in~\eqref{Coercivity}, and~\eqref{Sub-ellipticity-Strip} for $X_0 < 0 < X_1$ with $|X_0|, |X_1| \leq 1$. This ends the proof of Lemma~\ref{Lemma-Local-Carl-x-0}.



\subsection{A gluing argument: end of the proof of Theorem~\ref{Thm-Carleman-General}}\label{Sec-Gluing}

We then perform a gluing argument, which essentially consists in integrating the local Carleman estimates~\eqref{Carleman-Local-w-x0-L2} and~\eqref{Carleman-Local-w-x0-L2_*}, or rather the square of these estimates, with respect to $x_0 \in \overline\Omega \setminus \omega$, in order to deduce estimates~\eqref{Carleman-General-1} and~\eqref{Carleman-General-2}, respectively. We will only explain how to deduce estimate~\eqref{Carleman-General-1} from the estimate~\eqref{Carleman-Local-w-x0-L2}, since the other argument is completely similar.

We thus start from~\eqref{Carleman-Local-w-x0-L2}: There exist constants $C>0$ and $\tau_0 \geq 1$ such that for all $x_0 \in \overline\Omega \setminus \omega$ and $\tau \geq \tau_0$,
\begin{multline*}
\tau^{3} \| w_{x_0} \|_{L^2(\Omega)}^2 +
\tau^{\frac{3}{2}-\frac{1}{d}} \| w_{x_0} \|_{L^{\frac{2d}{d-2}}(\Omega)}^2 +
\tau \| \nabla w_{x_0} \|_{L^2(\Omega)}^2
\\
\leq C\left(
\norm{f_{2,x_0}}_{L^2(\Omega)} ^2 +
\tau^{\frac{3}{2}-\frac{1}{d}} \norm{f_{2_*',x_0} }_{L^{\frac{2d}{d+2}}(\Omega)}^2 +
\tau^2 \norm{F_{x_0}}_{L^2(\Omega)}^2 +
\tau^{\frac{3}{2}} \norm{g_{x_0}}_{H^{\frac{1}{2}}(\partial\Omega)}^2
\right).
\end{multline*}

Using the explicit expressions of the source terms, we obtain:
\begin{multline*}
\tau^{3} \| w_{x_0} \|_{L^2(\Omega)}^2 +
\tau^{\frac{3}{2} - \frac{1}{d}} \| w_{x_0} \|_{L^{\frac{2d}{d-2}}(\Omega)}^2 +
\tau \| \nabla w_{x_0} \|_{L^2(\Omega)}^2
\\
\leq C
\left(
\norm{\eta_{x_0} \tilde f_{2}}_{L^2(\Omega)} ^2 +
\tau^{\frac{3}{2}-\frac{1}{d}} \norm{ \eta_{x_0} \tilde f_{2_*'}}_{L^{\frac{2d}{d+2}}(\Omega)}^2 +
\tau^2 \norm{\eta_{x_0} \tilde F}_{L^2(\Omega)}^2 +
\tau^{\frac{3}{2}} \norm{\eta_{x_0} \tilde g}_{H^{\frac{1}{2}}(\partial\Omega)}^2
\right)
\\
+ C
\left(\norm{\nabla \eta_{x_0} \cdot \tilde F }_{L^2(\Omega)}^2 +
\tau^2 \| w_{x_0} \|_{L^2(\Omega)}^2 +
\norm{ \nabla \eta_{x_0} \cdot \nabla w }_{L^2(\Omega)}^2 +
\norm{\Delta \eta_{x_0} w}_{L^2(\Omega)}^2\right.\\ \left. +
\tau^2 \norm{ |\nabla \eta_{x_0}| w }_{L^2(\Omega)}^2
\right).
\end{multline*}
By taking $\tau_0\geq 1$ larger if necessary (which can be done uniformly in $x_0 \in \overline\Omega\setminus\omega$), we can absorb the term $\tau^2 \| w_{x_0} \|_{L^2(\Omega)}^2$, and we get for all $x_0 \in \overline\Omega\setminus\omega$, for all $\tau \geq \tau_0$,
\begin{multline*}
\tau^{3} \| w_{x_0} \|_{L^2(\Omega)}^2 +
\tau^{\frac{3}{2}-\frac{1}{d}} \| w_{x_0} \|_{L^{\frac{2d}{d-2}}(\Omega)}^2 +
\tau \| \nabla w_{x_0} \|_{L^2(\Omega)}^2
\\
\leq C
\left(
\norm{\eta_{x_0} \tilde f_{2}}_{L^2(\Omega)} ^2 +
\tau^{\frac{3}{2}-\frac{1}{d}} \norm{\eta_{x_0} \tilde f_{2_*'}}_{L^{\frac{2d}{d+2}}(\Omega)}^2 +
\tau^2 \norm{\eta_{x_0} \tilde F}_{L^2(\Omega)}^2 +
\tau^{\frac{3}{2}} \norm{\eta_{x_0} \tilde g}_{H^{\frac{1}{2}}(\partial\Omega)}^2
\right)
\\
+ C
\left(\norm{ \nabla \eta_{x_0} \cdot \tilde F }_{L^2(\Omega)}^2 +
\norm{\nabla \eta_{x_0} \cdot \nabla w}_{L^2(\Omega)}^2 +
\norm{\Delta \eta_{x_0} w}_{L^2(\Omega)}^2 +
\tau^2 \norm{ |\nabla \eta_{x_0}| w}_{L^2(\Omega)}^2
\right).
\end{multline*}

Now, integrating in $x_0$ on $\Omega \setminus \omega$ and using Fubini's identity for the Hilbertian norms, we get
\begin{multline*}
\tau^3 \int_\Omega \rho_0(x) |w(x)|^2 \, dx +
\tau \int_\Omega \rho_0(x) |\nabla w(x)|^2 \, dx +
\tau^{\frac{3}{2}-\frac{1}{d}} \int_{\Omega \setminus \omega} \| \eta_{x_0} w \|_{L^{\frac{2d}{d-2}}(\Omega)}^2 \, dx_0
\\
\begin{aligned}
\leq&
C\left(
\int_\Omega \rho_0(x) \left|\tilde f_2(x)\right|^2\, dx +
\int_\Omega (\tau^2 \rho_0(x) + \rho_{r,1}(x)) \left|\tilde F(x)\right|^2\, dx
\right)
\\
&+ C
\left(
\tau^{\frac{3}{2}-\frac{1}{d}} \int_{\Omega \setminus \omega} \norm{\eta_{x_0} \tilde f_{2_*'} }_{L^{\frac{2d}{d+2}}(\Omega)}^2 \, dx_0 +
\tau^{\frac{3}{2}} \int_{\Omega \setminus \omega} \norm{ \eta_{x_0} \tilde g}_{H^{\frac{1}{2}}(\partial \Omega)}^2 \, dx_0
\right)
\\
& + C
\left(
\int_\Omega \left(\rho_{r,2}(x) + \tau^2 \rho_{r,1}(x)\right) |w(x)|^2 \, dx +
\int_\Omega \rho_{r,1}(x) |\nabla w(x)|^2 \, dx
\right),
\end{aligned}
\end{multline*}
where the weights $\rho_0$, $\rho_{r,i}$ are defined as follows:
\begin{equation*}
\rho_0(x) \!= \!\int_{\Omega \setminus \omega} \!\!\! |\eta_{x_0}(x)|^2 \, dx_0, \quad
\rho_{r,1}(x) \! =\! \int_{\Omega \setminus \omega} \!\!\!|\nabla \eta_{x_0}(x)|^2 \, dx_0,
\quad
\rho_{r,2}(x) \! = \!\int_{\Omega \setminus \omega}\!\!\! |\Delta \eta_{x_0}(x)|^2 \, dx_0.
\end{equation*}
Taking an open subset $\omega_a$ such that $\overline\omega \subset \omega_a$ and $\overline{\omega_a} \subset \omega_1$, it is easy to check from the choice~\eqref{Choice-Localization} that
\begin{align}
\begin{aligned}\label{Est-rho-0-above}
\forall x \in \overline{\Omega} \setminus \omega_a,\quad  |\rho_0(x) | &\geq \frac{\tau^{-\frac{d}{3}}}{3} \norm{\eta}_{L^2}^2,
%\notag
\\
\forall x \in \Omega, \quad  |\rho_{0}(x)| &\leq C \tau^{- \frac{d}{3}},
\end{aligned}
\\
\begin{aligned}\label{Est-rho-2-above}
\forall x \in \Omega, \quad  |\rho_{r,1}(x)| &\leq C \tau^{\frac{2}{3} - \frac{d}{3}},
\\
\forall x \in \Omega, \quad  |\rho_{r,2}(x)| &\leq C \tau^{\frac{4}{3} - \frac{d}{3}}.
\end{aligned}
%\notag
\end{align}
Thus, for $\tau$ large enough,
\begin{multline*}
\tau^3 \int_{\Omega\setminus\omega_a} |w(x)|^2 \, dx +
\tau \int_{\Omega \setminus \omega_a} \left|\nabla w(x)\right|^2 \, dx +
\tau^{\frac{3}{2}+\frac{d}{3}-\frac{1}{d}} \int_{\Omega \setminus \omega} \| \eta_{x_0} w \|_{L^{\frac{2d}{d-2}}(\Omega)}^2 \, dx_0
\\
\leq C \int_\Omega \left|\tilde f_2(x)\right|^2\, dx + C \left(
\tau^2 \int_\Omega \left|\tilde F(x)\right|^2\, dx + \tau^{\frac{3}{2}+\frac{d}{3}-\frac{1}{d}} \int_{\Omega \setminus \omega} \norm{\eta_{x_0} \tilde f_{2_*'}}_{L^{\frac{2d}{d+2}}(\Omega)}^2 \, dx_0
\right)
\\
+ C
\left(
\tau^{\frac{3}{2}+\frac{d}{3}} \int_{\Omega \setminus \omega} \norm{\eta_{x_0} \tilde g }_{H^{\frac{1}{2}}(\partial \Omega)}^2 \, dx_0 + \tau^{\frac{8}{3}} \int_\omega |w(x)|^2 \, dx +
\tau^{\frac{2}{3}} \int_\omega |\nabla w(x)|^2 \, dx
\right).
\end{multline*}
We then add
\[
\tau^3 \int_{\omega_a} |w(x)|^2 \, dx +
\tau \int_{ \omega_a} |\nabla w(x)|^2 \, dx +
\tau^{\frac{3}{2}+\frac{d}{3}-\frac{1}{d}} \int_{ \omega_a} \| \eta_{x_0} w \|_{L^{\frac{2d}{d-2}}(\Omega)}^2 \, dx_0
\]
to both sides of the previous estimate and get
\begin{multline}\label{Carleman-A-Nettoyer}
\tau^3 \int_{\Omega} |w(x)|^2 \, dx +
\tau \int_{\Omega} |\nabla w(x)|^2 \, dx +
\tau^{\frac{3}{2}+\frac{d}{3}-\frac{1}{d}} \int_{\Omega} \| \eta_{x_0} w \|_{L^{\frac{2d}{d-2}}(\Omega)}^2 \, dx_0 \\
\leq C
\left(\int_\Omega \left|\tilde f_2(x)\right|^2\, dx +\tau^2 \int_\Omega \left|\tilde F(x)\right|^2\, dx +
\tau^{\frac{3}{2}+\frac{d}{3}-\frac{1}{d}} \int_{\Omega \setminus \omega} \left\| \eta_{x_0} \tilde f_{2_*'} \right\|_{L^{\frac{2d}{d+2}}(\Omega)}^2 \, dx_0
\right)
\\
+ C
\left(
\tau^{\frac{3}{2}+\frac{d}{3}} \int_{\Omega \setminus \omega} \left\| \eta_{x_0} \tilde g \right\|_{H^{\frac{1}{2}}(\partial \Omega)}^2 \, dx_0+
\tau^{3} \int_{\omega_a} |w(x)|^2 \, dx +
\tau \int_{\omega_a} |\nabla w(x)|^2 \, dx
\right)
\\
+C
\tau^{\frac{3}{2}+\frac{d}{3}-\frac{1}{d}} \int_{\omega_a} \| \eta_{x_0} w \|_{L^{\frac{2d}{d-2}}(\Omega)}^2 \, dx_0.
\end{multline}

We claim the following lemma, proven in Appendix~\ref{Appendix-B}:

\begin{lemm}\label{Lem-Boundary-H-1/2}
With $\eta_{x_0}$ as in~\eqref{Choice-Localization}, there exists a constant $C>0$ such that for all $\tau \geq \tau_0$ and $g \in H^{\frac{1}{2}}(\partial \Omega)$,
\begin{equation}\label{Boundary-H-1/2-norms}
\tau^{\frac{3}{2}+\frac{d}{3}} \int_{\Omega \setminus \omega} \| \eta_{x_0} g \|_{H^{\frac{1}{2}}(\partial \Omega)}^2 \, dx_0
\leq C \left(
\tau^{\frac{3}{2}} \| g \|_{H^{\frac{1}{2}}(\partial \Omega)}^2 + \tau^{11/6} \|g\|_{L^2(\partial \Omega)}^2\right).
\end{equation}
\end{lemm}

We thus have the bound
\begin{equation*}
\tau^{\frac{3}{2}+\frac{d}{3}} \int_{\Omega \setminus \omega} \norm{ \eta_{x_0} \tilde g}_{H^{\frac{1}{2}}(\partial \Omega)}^2 \, dx_0
\leq C\left(
\tau^{\frac{3}{2}} \| \tilde g \|_{H^{\frac{1}{2}}(\partial \Omega)}^2 + \tau^{11/6} \| \tilde g\|_{L^2(\partial \Omega)}^2\right).
\end{equation*}
Now, $\tilde g$ is the trace of the function $w \in H^1(\Omega)$. Taking $X \in W^{1,\infty}(\Omega; \R^d)$ such that $X \cdot n = 1$ on $\partial \Omega$,
\begin{multline*}
\| \tilde g\|_{L^2(\partial \Omega)}^2 =
\int_{\Omega} \div \left(X |w|^2\right) \, dx =
\int_{\Omega} \div(X) |w|^2 \, dx +2 \int_\Omega X \cdot \nabla w w \, dx \\
\leq C \| w\|_{L^2(\Omega)}^2 + \|w \|_{L^2(\Omega)} \| \nabla w \|_{L^2(\Omega)}.
\end{multline*}
It follows that the term $\tau^{11/6} \| \tilde g\|_{L^2(\partial \Omega)}^2$ can in fact be absorbed by the left-hand side of~\eqref{Carleman-A-Nettoyer} by taking $\tau_0$ larger if necessary, and we obtain:
\begin{multline*}
\tau^3 \int_{\Omega} |w(x)|^2 \, dx +
\tau \int_{\Omega} |\nabla w(x)|^2 \, dx +
\tau^{\frac{3}{2}+\frac{d}{3}-\frac{1}{d}} \int_{\Omega} \| \eta_{x_0} w \|_{L^{\frac{2d}{d-2}}(\Omega)}^2 \, dx_0
\\
\leq  C 
\left(
\int_\Omega  \left|\tilde f_2(x)\right|^2 dx+
\tau^2  \int_\Omega \left|\tilde F(x)\right|^2 dx  + 
\tau^{\frac{3}{2}} \norm{\tilde g }_{H^{\frac{1}{2}}(\partial \Omega)}^2\!\right.
\\
\left.+
\tau^{\frac{3}{2}+\frac{d}{3}-\frac{1}{d}} \int_{\Omega \setminus \omega} \norm{\eta_{x_0} \tilde f_{2_*'}}_{L^{\frac{2d}{d+2}}(\Omega)}^2 \, dx_0
 +
\tau^{3} \int_{\omega_a} |w(x)|^2 \, dx 
\right.
\\
\left.
+
\tau \int_{\omega_a} |\nabla w(x)|^2 \, dx +
\tau^{\frac{3}{2}+\frac{d}{3}-\frac{1}{d}} \int_{ \omega_a} \| \eta_{x_0} w \|_{L^{\frac{2d}{d-2}}(\Omega)}^2 \, dx_0
\right).
\end{multline*}
We then go back to the variable $u$ and get the following estimate:
\begin{multline}\label{Carleman-General-almost}
\tau^{3} \norm{e^{\tau \varphi} u}_{L^2(\Omega)}^2 +
\tau \norm{ e^{\tau \varphi} \nabla u }_{L^2(\Omega)}^2 +
\tau^{\frac{3}{2}+\frac{d}{3}-\frac{1}{d}} \int_{\Omega} \norm{\eta_{x_0} u e^{\tau \varphi}}_{L^{\frac{2d}{d-2}}(\Omega)}^2 \, dx_0
\\
\!\leq \! C \!\left(\!
\norm{ e^{\tau \varphi} f_2}_{L^2(\Omega)}^2\! +\!
\tau^{\frac{3}{2}+\frac{d}{3}-\frac{1}{d}}\!\!\! \int_{\Omega}\! \norm{ \eta_{x_0} f_{2_*'} e^{\tau \varphi}}_{L^{\frac{2d}{d+2}}(\Omega)}^2 \!\!dx_0\! + \!
\tau^2 \norm{e^{\tau \varphi} F}_{L^2(\Omega)}^2\! 
\right.
\\
\left. 
+ \!
\tau^{\frac{3}{2}} \norm{e^{\tau \varphi} g}_{H^{\frac{1}{2}}(\partial \Omega)}^2
+
\tau^{3} \norm{e^{\tau \varphi} u}_{L^2(\omega_a)}^2
\right.
\\
\left. 
+
\tau \norm{e^{\tau \varphi} \nabla u}_{L^2(\omega_a)}^2 +
\tau^{\frac{3}{2}+\frac{d}{3}-\frac{1}{d}} \int_{\omega_a} \norm{ \eta_{x_0} u e^{\tau \varphi} }_{L^{\frac{2d}{d-2}}(\Omega)}^2 \, dx_0
\right).
\end{multline}
We finally explain how to remove the term $\tau \| e^{\tau \varphi} \nabla u \|_{L^2(\omega_a)}^2$ from the right-hand side of~\eqref{Carleman-General-almost}. In order to do that, we choose an open subset $\omega_b$ of $\Omega$ such that $\overline{\omega_a} \subset \omega_b$ and $\overline\omega_b \subset \omega_1$, and a smooth compactly supported function $\eta_\omega$, taking value $1$ in $\omega_a$ and vanishing in $\overline\Omega\setminus \omega_b$. We then multiply~\eqref{Elliptic-Eq-general} by $\eta_\omega u e^{2 \tau \varphi}$, which yields:
\begin{multline*}
\int_\Omega \eta_\omega |\nabla u|^2 e^{2 \tau \varphi} \,dx =
\int_{\Omega} \eta_\omega f_2 u e^{2 \tau \varphi} \, dx +
\int_{\Omega} \eta_\omega f_{2_*'} u e^{2 \tau \varphi} \, dx -
\int_\Omega F \cdot \nabla \left(\eta_\omega e^{2 \tau \varphi}\right) u \, dx
\\
-
\int_\Omega F \cdot \nabla u \, \eta_\omega e^{2 \tau \varphi} \, dx +
\frac{1}{2} \int_\Omega \Delta \left(\eta_\omega e^{2 \tau \varphi}\right) |u|^2 \, dx.
\end{multline*}
Using the bound,
\[
\left |
\int_\Omega F \cdot \nabla u \eta_\omega e^{2 \tau \varphi}\, dx
\right|
\leq
\frac{1}{2} \| F e^{\tau \varphi} \|_{L^2(\Omega)}^2 +
\frac{1}{2} \int_\Omega \eta_\omega |\nabla u|^2 e^{2 \tau \varphi} \, dx,
\]
we easily get
\begin{multline*}
\norm{ e^{\tau \varphi} \nabla u}_{L^2(\omega_a)}^2
\leq C\norm{f_2 e^{\tau \varphi}}_{L^2(\Omega)}
\norm{u e^{\tau \varphi}}_{L^2(\omega_b)} + C \norm{ F e^{\tau \varphi} }_{L^2(\Omega)}^2
\\
+ C \tau\norm{ F e^{\tau \varphi}}_{L^2(\Omega)} \norm{u e^{\tau \varphi}}_{L^2(\omega_b)} + C
\tau^2\norm{u e^{\tau \varphi}}_{L^2(\omega_b)}^2 + C \left| \int_{\Omega} \eta_\omega f_{2_*'} u e^{2 \tau \varphi} \, dx\right|.
\end{multline*}
Only the last term is unusual. In order to estimate it, we remark that, for all $x \in \omega_b$ (recall that $\omega_{1}$, hence $\omega_b$, is at a positive distance from $\partial \Omega$), taking $\tau_0\geq 1$ larger if necessary, for all $\tau \geq \tau_0$,
\[
\int_{x_0\, \in\, \Omega} \eta_{x_0}(x)^2 \, dx_0 = \| \eta\|_{L^2\left(\R^d\right)}^2 \tau^{-\frac{d}{3}}.
\]
Accordingly,
\begin{align*}
& \int_{\Omega} \eta_\omega f_{2_*'} u e^{2 \tau \varphi} \, dx =
\frac{\tau^{\frac{d}{3}}}{\| \eta\|_{L^2(\R^d)}^2} \int_{ x_0\, \i\,n \Omega} \int_{x\, \in\, \Omega} \eta_\omega f_{2_*'} u e^{2 \tau \varphi} \eta_{x_0}(x)^2 \, dx \, dx_0
\\
&\leq C \tau^{\frac{d}{3}} \int_{x_0\,\in\, \Omega} \norm{\eta_{x_0} f_{2_*'}e^{\tau \varphi}}_{L^{\frac{2d}{d+2}}(\Omega)}
\norm{ \eta_{x_0} u e^{\tau \varphi}}_{L^{\frac{2d}{d-2}}(\omega_b)} \, dx_0
\\
&\leq C \tau^{\frac{d}{3}} \norm{\eta_{x_0} f_{2_*'}e^{\tau \varphi}}_{L^2_{x_0}\left(\Omega; L^{\frac{2d}{d+2}}_x(\Omega)\right)} \norm{\eta_{x_0} u e^{\tau \varphi}}_{L^2_{x_0}\left(\Omega; L^{\frac{2d}{d-2}}_x(\omega_b)\right)}.
\end{align*}
One then easily gets that
\begin{multline*}
\tau \norm{ e^{\tau \varphi} \nabla u}_{L^2(\omega_a)}^2\\
 \leq  C  \left( 
\norm{e^{\tau \varphi} f_2 }_{L^2(\Omega)}^2\! +
\tau^{\frac{3}{2}+\frac{d}{3}-\frac{1}{d}} \norm{\eta_{x_0} f_{2_*'}e^{\tau \varphi}}_{L^2_{x_0}(\Omega; L^{\frac{2d}{d+2}}_x(\Omega))}^2 + 
\tau^2 \norm{e^{\tau \varphi} F }_{L^2(\Omega)}^2
\right.
\\
\left. +
\tau^{\frac{3}{2}} \norm{ e^{\tau \varphi} g }_{H^{\frac{1}{2}}(\partial \Omega)}^2 +
\tau^{3} \norm{ e^{\tau \varphi} u}_{L^2(\omega_b)}^2 +
\tau^{\frac{3}{2}+\frac{d}{3}-\frac{1}{d}} \norm{ \eta_{x_0} u e^{\tau \varphi}}_{L^2_{x_0}\left(\Omega; L^{\frac{2d}{d-2}}_x(\omega_b)\right)}^2
\right),
\end{multline*}
which concludes the proof of~\eqref{Carleman-General-1} since we obviously have
\[
\norm{ \eta_{x_0} u e^{\tau \varphi}}_{L^2_{x_0}\left(\Omega; L^{\frac{2d}{d-2}}_x(\omega_b)\right)}^2
\leq
\norm{\eta_{x_0} u e^{\tau \varphi}}_{L^2_{x_0}\left(\omega_1; L^{\frac{2d}{d-2}}_x(\omega_1)\right)}^2
\]
for $\tau_0$ sufficiently large so that $\tau_0^{\frac{1}{3}} \geq 1/ d(\omega_b, \Omega \setminus \omega_1)$.


\section{Proof of Theorem~\ref{Thm-QUCP}: Quantitative unique continuation} \label{Sec-Proof-Thm-QUCP}

First, by restricting $\omega$ if necessary, we assume that $\omega$ is a non-empty open subset of $\Omega$ with $\overline\omega \subset \Omega$. Then, for $\omega_0$ a non-empty open subset such that $\overline{\omega_0} \subset \omega$, there exists a function $\varphi$ satisfying conditions~\eqref{CarlemanWeight-Cond1}--\eqref{CarlemanWeight-Cond2} in $\omega_0$ (see for instance~\cite[Lemma~1.1]{Fursikov1996-book} or~\cite[Proposition~3.31]{LeRousseaubookvolume1}), so that the Carleman estimates~\eqref{Carleman-General-1}--\eqref{Carleman-General-2} in Theorem~\ref{Thm-Carleman-General} with $\omega_1 = \omega$ hold.

For $V \in L^{q_0}(\Omega)$, $W_1 \in L^{q_1}(\Omega; \C^d)$, and $W_2 \in L^{q_2}(\Omega; \C^d)$, we consider decompositions of the form
\begin{align*}
& V = V_{\frac{d}{2}} + V_{d}+ V_\infty,
\quad &&\text{ with } V_{\frac{d}{2}} \in L^{\frac{d}{2}}(\Omega),
\ V_{d} \in L^{d}(\Omega),
\ V_{\infty} \in L^{\infty}(\Omega),
\\
& W_1 = W_{1,d} + W_{1,\infty},
\quad &&\text{ with } W_{1,d} \in L^{d}\left(\Omega; \C^d\right),
\ W_{1,\infty} \in L^{\infty}\left(\Omega; \C^d\right),
\\
& W_2 = W_{2,d} + W_{2,\infty},
\quad &&\text{ with } W_{2,d} \in L^{d}\left(\Omega; \C^d\right),
\ W_{2,\infty} \in L^{\infty}\left(\Omega; \C^d\right),
\end{align*}
which will be made precise later.

In particular, applying~\eqref{Carleman-General-1} for $u$ solution of~\eqref{Elliptic-Eq-general}, with $f_{2_*'} = V_{\frac{d}{2}} u + V_{d} u+ W_{1,d} \cdot \nabla u$, $f_2 = V_{\infty} u+ W_{1,\infty} \cdot \nabla u$, and using~\eqref{Lp-larger-Lpetatau} and~\eqref{Holder-in-Lp-tau-eta}, we get the existence of $C_1>0$ such that for $\tau \geq \tau_0$:

\begin{multline}\label{Carleman-General-1-appl}
\tau^{\frac{3}{2}} \| e^{\tau \varphi} u \|_{L^2(\Omega)} +
\tau^{\frac{1}{2}} \| e^{\tau \varphi} \nabla u \|_{L^2(\Omega)}\\
\begin{aligned}
&\leq C_1
\norm{e^{\tau \varphi} (V_{\infty} u+ W_{1,\infty} \cdot \nabla u)}_{L^2(\Omega)}
\\
&\quad +C_1
\tau^{\frac{3}{4}-\frac{1}{2d}} \norm{e^{\tau \varphi}\left(V_{\frac{d}{2}} u + V_{d} u+ W_{1,d} \cdot \nabla u\right)}_{L^{\frac{2d}{d+2}}_{\eta, \tau} (\Omega)}
\\
&\quad + C_1 \tau \norm{ e^{\tau \varphi}\left(W_{2,d} u + W_{2, \infty} u\right)}_{L^2(\Omega)}\\
&\quad + C_1 \left(\tau^{\frac{3}{2}} \norm{ e^{\tau \varphi} u}_{L^2(\omega)} + \tau^{\frac{3}{4}}\norm{ e^{\tau \varphi} u}_{L^{\frac{2d}{d-2}}_{\eta, \tau} (\omega)}\right)
\\
&\leq C_1 \left(\| V_\infty \|_{L^\infty(\Omega)} + \tau^{\frac{3}{4}-\frac{1}{2d}} \| V_{d} \|_{L^d(\Omega)} + \tau \| W_{2, \infty}\|_{L^\infty(\Omega)} \right)
\| e^{\tau \varphi} u \|_{L^2(\Omega)}
\\
&\quad + C_1 \left(\| W_{1,\infty} \|_{L^\infty(\Omega)} + \tau^{\frac{3}{4}-\frac{1}{2d}} \| W_{1,d} \|_{L^d(\Omega)} \right)
\| e^{\tau \varphi} \nabla u \|_{L^2(\Omega)}
\\
&\quad + C_1 \left(\tau^{\frac{3}{4}-\frac{1}{2d}} \norm{V_{\frac{d}{2}}}_{L^{\frac{d}{2}}(\Omega)} + \tau \| W_{2, d}\|_{L^d(\Omega)} \right)
\| e^{\tau \varphi} u \|_{L^{\frac{2d}{d-2}}_{\eta, \tau}(\Omega)}\\
&\quad + C_1 \tau^{\frac{3}{2}} \| e^{\tau \varphi} u\|_{L^{\frac{2d}{d-2}} (\omega)}.
\end{aligned}
\end{multline}
Similarly, applying~\eqref{Carleman-General-2} with $f_{2_*'} = V_{\frac{d}{2}} u + W_{1,d} \cdot \nabla u$, $f_2 = V_{d} u + V_{\infty} u+ W_{1,\infty} \cdot \nabla u$, we obtain the existence of a constant $C_2>0 $ such that for $\tau \geq \tau_0$,
\begin{multline}\label{Carleman-General-2-appl}
\tau^{\frac{3}{4}+\frac{1}{2d}} \| e^{\tau \varphi} u \|_{L^{\frac{2d}{d-2}}_{\eta,\tau} (\Omega)}
\leq C_2 \left(\| V_\infty \|_{L^\infty(\Omega)} + \tau \| W_{2, \infty}\|_{L^\infty(\Omega)} \right)
\norm{e^{\tau \varphi} u}_{L^2(\Omega)}
\\
\begin{aligned}
&+ C_2 \left(\| W_{1,\infty} \|_{L^\infty(\Omega)} + \tau^{\frac{3}{4}+\frac{1}{2d}} \| W_{1,d} \|_{L^d(\Omega)} \right)
\norm{e^{\tau \varphi} \nabla u}_{L^2(\Omega)}
\\
&+ C_2 \left(\| V_{d} \|_{L^d(\Omega)} \! + \tau^{\frac{3}{4}+\frac{1}{2d}} \| V_{\frac{d}{2}} \|_{L^{\frac{d}{2}}(\Omega)}\!\! + \tau \| W_{2, d}\|_{L^d(\Omega)} \right)
\norm{e^{\tau \varphi} u }_{L^{\frac{2d}{d-2}}_{\eta, \tau}(\Omega)} \\
&+ C_2\tau^{\frac{3}{2}} \norm{e^{\tau \varphi} u}_{L^{\frac{2d}{d-2}} (\omega)}.
\end{aligned}
\end{multline}

Thus, if we have
\begin{align}
& 2 C_1 \left(\| V_\infty \|_{L^\infty(\Omega)} + \tau^{\frac{3}{4}-\frac{1}{2d}} \| V_{d} \|_{L^d(\Omega)} + \tau \| W_{2, \infty}\|_{L^\infty(\Omega)} \right)
\leq \tau^{\frac{3}{2}},
\label{Condition-1}
\\
& 2 C_1 \left(\| W_{1,\infty} \|_{L^\infty(\Omega)} + \tau^{\frac{3}{4}-\frac{1}{2d}} \| W_{1,d} \|_{L^d(\Omega)} \right)
\leq \tau^{\frac{1}{2}},
\label{Condition-2}
\\
&2 C_2 \left(\| V_{d} \|_{L^d(\Omega)} + \tau^{\frac{3}{4}+\frac{1}{2d}} \norm{ V_{\frac{d}{2}}}_{L^{\frac{d}{2}}(\Omega)} + \tau \| W_{2, d}\|_{L^d(\Omega)} \right)
\leq \tau^{\frac{3}{4}+\frac{1}{2d}},
\label{Condition-3}
\end{align}
estimates~\eqref{Carleman-General-1-appl}--\eqref{Carleman-General-2-appl} yield:
\begin{multline*}
\tau^{\frac{3}{2}} \norm{ e^{\tau \varphi} u }_{L^2(\Omega)} +
\tau^{\frac{1}{2}} \norm{e^{\tau \varphi} \nabla u }_{L^2(\Omega)}\\
\leq 2C_1 \left(\tau^{\frac{3}{4}-\frac{1}{2d}} \norm{ V_{\frac{d}{2}}}_{L^{\frac{d}{2}}(\Omega)} + \tau \| W_{2, d}\|_{L^d(\Omega)} \right) 
\| e^{\tau \varphi} u \|_{L^{\frac{2d}{d-2}}_{\eta, \tau}(\Omega)}
+ 2C_1\tau^{\frac{3}{2}} \norm{e^{\tau \varphi} u}_{L^{\frac{2d}{d-2}} (\omega)}
\end{multline*}
and
\begin{multline*}
\tau^{\frac{3}{4}+\frac{1}{2d}} \norm{ e^{\tau \varphi} u }_{L^{\frac{2d}{d-2}}_{\eta,\tau} (\Omega)}
\leq 2C_2 \left(\| V_\infty \|_{L^\infty(\Omega)} + \tau \| W_{2, \infty}\|_{L^\infty(\Omega)} \right)
\norm{ e^{\tau \varphi} u }_{L^2(\Omega)}
\\
+ 2C_2 \left(\| W_{1,\infty} \|_{L^\infty(\Omega)} + \tau^{\frac{3}{4}+\frac{1}{2d}} \| W_{1,d} \|_{L^d(\Omega)} \right)
\norm{e^{\tau \varphi} \nabla u}_{L^2(\Omega)} + 2 C_2 \tau^{\frac{3}{2}} \norm{ e^{\tau \varphi} u}_{L^{\frac{2d}{d-2}} (\omega)}.
\end{multline*}
Following, under conditions~\eqref{Condition-1}--\eqref{Condition-2}--\eqref{Condition-3}, we also have
\begin{multline*}
\tau^{\frac{3}{2}} \norm{ e^{\tau \varphi} u }_{L^2(\Omega)} +
\tau^{\frac{1}{2}} \norm{e^{\tau \varphi} \nabla u }_{L^2(\Omega)}
\\
\begin{aligned}
& \leq 4C_1 C_2 \tau^{-\frac{3}{4} - \frac{1}{2d}} 
\left(\tau^{\frac{3}{4}-\frac{1}{2d}} \norm{ V_{\frac{d}{2}}}_{L^{\frac{d}{2}}(\Omega)} + \tau \| W_{2, d}\|_{L^d(\Omega)} \right)\\
&\mkern175mu
\left(\| V_\infty \|_{L^\infty(\Omega)} + \tau \| W_{2, \infty}\|_{L^\infty(\Omega)} \right)
\| e^{\tau \varphi} u \|_{L^2(\Omega)}
\\ &
 +   4 C_1 C_2 \tau^{-\frac{3}{4} - \frac{1}{2d}} 
\left( \tau^{\frac{3}{4}-\frac{1}{2d}} \| V_{\frac{d}{2}} \|_{L^{\frac{d}{2}}(\Omega)} +  \tau \| W_{2, d}\|_{L^d(\Omega)} \right)  
\\
&\mkern175mu
\left( \| W_{1,\infty} \|_{L^\infty(\Omega)} + \tau^{\frac{3}{4}+\frac{1}{2d}} \| W_{1,d} \|_{L^d(\Omega)} \right)
\times
\| e^{\tau \varphi} \nabla u \|_{L^2(\Omega)}
\\
& + \! 2 C_1\left(\! C_2 \left(\tau^{\frac{3}{4}-\frac{1}{2d}} \| V_{\frac{d}{2}} \|_{L^{\frac{d}{2}}(\Omega)} + \tau \| W_{2, d}\|_{L^d(\Omega)} \right) \tau^{\frac{3}{4}}\! + \tau^{\frac{3}{2}}\right)
\norm{e^{\tau \varphi} u}_{L^{\frac{2d}{d-2}} (\omega)}.
\end{aligned}
\end{multline*}
Note that this estimate yields an observation estimate if the additional following conditions are also satisfied:
\begin{align}
& 8 C_1 C_2 \tau^{-\frac{3}{4} - \frac{1}{2d}}
\left(\tau^{\frac{3}{4}-\frac{1}{2d}} \norm{ V_{\frac{d}{2}} }_{L^{\frac{d}{2}}(\Omega)} + \tau \| W_{2, d}\|_{L^d(\Omega)} \right)
\label{Condition-4}
\\
& \mkern200mu\times
\left(\| V_\infty \|_{L^\infty(\Omega)}+ \tau \| W_{2, \infty}\|_{L^\infty(\Omega)} \right)
\leq \tau^{\frac{3}{2}},
\notag
\\
& 8 C_1 C_2 \tau^{-\frac{3}{4} - \frac{1}{2d}}
\left(\tau^{\frac{3}{4}-\frac{1}{2d}} \norm{ V_{\frac{d}{2}}}_{L^{\frac{d}{2}}(\Omega)} + \tau \| W_{2, d}\|_{L^d(\Omega)} \right)
\label{Condition-5}
\\
& \mkern200mu\times
\left(\! \| W_{1,\infty} \|_{L^\infty(\Omega)}+\tau^{\frac{3}{4}+\frac{1}{2d}} \| W_{1,d} \|_{L^d(\Omega)} \right)
\leq \tau^{\frac{1}{2}}.\notag
\end{align}
Indeed, in this case, one would obtain
\begin{multline}\label{Carleman-General-appl-final}
\tau^{\frac{3}{2}} \| e^{\tau \varphi} u \|_{L^2(\Omega)} +
\tau^{\frac{1}{2}} \| e^{\tau \varphi} \nabla u \|_{L^2(\Omega)}\\
\leq 2 C_1\left(C_2 \left(\tau^{\frac{3}{4}-\frac{1}{2d}} \norm{ V_{\frac{d}{2}} }_{L^{\frac{d}{2}}(\Omega)} + \tau \| W_{2, d}\|_{L^\infty(\Omega)} \right) \tau^{\frac{3}{4}} + \tau^{\frac{3}{2}}\right) \norm{ e^{\tau \varphi} u}_{L^{\frac{2d}{d-2}} (\omega)}.
\end{multline}


Therefore, our next step is to understand how, given $V \in L^{q_0}(\Omega)$, $W_1 \in L^{q_1}(\Omega; \C^d)$, and $W_2 \in L^{q_2}(\Omega; \C^d)$, one can minimize the value of $\tau$ for which we can find decompositions such that conditions~\eqref{Condition-1}--\eqref{Condition-2}--\eqref{Condition-3}--\eqref{Condition-4}--\eqref{Condition-5} are satisfied.

Before going further, let us remark that~\eqref{Condition-1}--\eqref{Condition-2}--\eqref{Condition-3}--\eqref{Condition-4}--\eqref{Condition-5} are satisfied provided, for $c_0 >0$ small enough,
\begin{equation}\label{Conditions-Summary}
\left\{
\begin{array}{l}
\ds \| V_\infty \|_{L^\infty(\Omega)} \leq c_0 \tau^{\frac{3}{2}},
\quad
\| V_d\|_{L^{d}(\Omega)} \leq c_0 \tau^{\frac{3}{4} + \frac{1}{2d}},
\quad
\norm{ V_{\frac{d}{2}}}_{L^{\frac{d}{2}}(\Omega)} \leq c_0,
\medskip
\\
\ds \| W_{1, \infty} \|_{L^\infty(\Omega)}+ \| W_{2, \infty} \|_{L^\infty(\Omega)} \leq c_0 \tau^{\frac{1}{2}},\\[0.20em]
\tau^{\frac{1}{4} - \frac{1}{2d}} \left(\| W_{1,d} \|_{L^d(\Omega)} + \| W_{2, d} \|_{L^d(\Omega)}\right) \leq c_0,
\medskip
\\
\ds
\tau^{\frac{1}{2}} \| W_{1,d} \|_{L^d(\Omega)} \| W_{2, d} \|_{L^d(\Omega)} \leq c_0.
\end{array}\right.
\end{equation}
\paragraph{The case $V \in L^{q_0}(\Omega)$ with $q_0 \in (\frac{d}{2}, d]$.}
For $V \in L^{q_0}(\Omega)$ with $q_0 \in (\frac{d}{2}, d]$ and $\lambda_0>0$ to be chosen later, we set $V_{\frac{d}{2}} = \ = V 1_{|V|\,>\,\lambda_0}$, $ V_d = V 1_{|V|\,\leq\,\lambda_0}$, $V_\infty = 0$, for which we have the following estimates:
\[
\lambda_0^{q_0-\frac{d}{2}} \norm{V_{\frac{d}{2}}}_{L^{\frac{d}{2}}(\Omega)}^{\frac{d}{2}} \leq \| V \|_{L^{q_0}(\Omega)}^{q_0},
\qquad
\| V_{d} \|_{L^d(\Omega)}^d \leq \lambda_0^{d - q_0} \| V \|_{L^{q_0}(\Omega)}^{q_0}.
\]

Similarly, for $W_1 \in L^{q_1}(\Omega)$ and $W_2 \in L^{q_2}(\Omega)$ and $\lambda_1, \lambda_2$ positive numbers to be chosen later, we set $W_{1,d} = W_1 1_{|W_1|\, >\,\lambda_1}$, $W_{1, \infty} = W_1 1_{|W_1|\,\leq\, \lambda_1}$, and $W_{2,d} = W_2 1_{|W_2|\,>\, \lambda_2}$, $W_{2, \infty} = W_2 1_{|W_2|\,\leq\,\lambda_2}$, for which we have the estimates:
\begin{equation}\label{Est-W-1-d-infty}
\left\{
\begin{array}{l}
\ds \lambda_1^{q_1-d} \| W_{1,d}\|_{L^{d}(\Omega)}^{d} \leq \| W_1 \|_{L^{q_1}(\Omega)}^{q_1},
\qquad
\| W_{1,\infty} \|_{L^\infty(\Omega)} \leq \lambda_1,
\\[0.15em]
\lambda_2^{q_2-d} \| W_{2,d}\|_{L^{d}(\Omega)}^{d} \leq \| W_2 \|_{L^{q_2}(\Omega)}^{q_2},
\qquad
\| W_{2,\infty} \|_{L^\infty(\Omega)} \leq \lambda_2.
\end{array}
\right.
\end{equation}

Conditions~\eqref{Conditions-Summary} are thus satisfied provided
\begin{equation}\label{Conditions-W1=W2=0-q0<d-a}
\left\{
\begin{array}{l}
\ds \lambda_0^{1 - \frac{q_0}{d}} \|V \|_{L^{q_0}(\Omega)}^{\frac{q_0}{d}} \leq c_0 \tau^{\frac{3}{4}+\frac{1}{2d}},
\quad
\lambda_0^{1- 2\frac{q_0}{d}} \|V \|_{L^{q_0}(\Omega)}^{\frac{2q_0}{d}} \leq c_0,
\\[0.15em]
\ds
\lambda_1 + \lambda_2 \leq c_0 \tau^{\frac{1}{2}}, \quad
\tau^{\frac{1}{4}-\frac{1}{2d}} \left(\lambda_1^{1- \frac{q_1}{d}} \| W_{1} \|_{L^{q_1}(\Omega)}^{ \frac{q_1}{d}} + \lambda_2^{1- \frac{q_2}{d}} \| W_{2} \|_{L^{q_2}(\Omega)}^{ \frac{q_2}{d}}\right) \leq c_0,
\medskip
\\
\ds \tau^{\frac{1}{2}} \lambda_1^{1- \frac{q_1}{d}} \| W_{1} \|_{L^{q_1}(\Omega)}^{ \frac{q_1}{d}} \lambda_2^{1- \frac{q_2}{d}} \| W_{2} \|_{L^{q_2}(\Omega)}^{ \frac{q_2}{d}} \leq c_0.
\end{array}\right.
\end{equation}
We then choose $\lambda_0 = \tau^{\alpha_0} \| V\|_{L^{q_0}(\Omega)}$, $\lambda_1 = \tau^{\alpha_1} \| W_1 \|_{L^{q_1}(\Omega)}$ and $\lambda_2 = \tau^{\alpha_2} \| W_2 \|_{L^{q_2}(\Omega)}$ for some real parameters $\alpha_0,\, \alpha_1, \, \alpha_2$, so that conditions~\eqref{Conditions-W1=W2=0-q0<d-a} yield:
\begin{equation}\label{Conditions-W1=W2=0-q0<d-b}
\left\{
\begin{array}{l}
\ds
\tau^{\alpha_0\left(1 - \frac{q_0}{d}\right)-\frac{3}{4} - \frac{1}{2d}} \|V \|_{L^{q_0}(\Omega)} \leq c_0,
\quad
\tau^{\alpha_0\left(1 - 2 \frac{q_0}{d}\right)} \|V \|_{L^{q_0}(\Omega)} \leq c_0,
\medskip
\\
\ds
\tau^{\alpha_1 - \frac{1}{2}} \| W_1\|_{L^{q_1}(\Omega)} + \tau^{\alpha_2 - \frac{1}{2}} \| W_2\|_{L^{q_2}(\Omega)} \leq c_0, \quad
\medskip
\\
\ds \tau^{\frac{1}{4}-\frac{1}{2d} + \alpha_1\left(1- \frac{q_1}{d}\right)} \| W_{1} \|_{L^{q_1}(\Omega)} +
\tau^{\frac{1}{4}-\frac{1}{2d} + \alpha_2\left(1- \frac{q_2}{d}\right)} \| W_{2} \|_{L^{q_2}(\Omega)}
\leq c_0,
\medskip
\\
\ds
\tau^{\frac{1}{2}+ \alpha_1\left(1- \frac{q_1}{d}\right)+ \alpha_2\left(1- \frac{q_2}{d}\right)} \| W_{1} \|_{L^{q_1}(\Omega)} \| W_{2} \|_{L^{q_2}(\Omega)} \leq c_0.
\end{array}\right.
\end{equation}
For $q_0 \in (\frac{d}{2}, d]$, $q_1 > \frac{3d-2}{2} $, $q_2 > \frac{3d-2}{2} $ with $\frac{1}{q_1} +\frac{1}{q_2} < 4(1-\frac{1}{d})/(3d -2)$, we choose $\alpha_0 = (\frac{3}{4} + \frac{1}{2d})\frac{d}{q_0}$, $\alpha_1 = (\frac{3}{4} - \frac{1}{2d}) \frac{d}{q_1}$, $\alpha_2 = (\frac{3}{4} - \frac{1}{2d}) \frac{d}{q_2}$, so that system~\eqref{Conditions-W1=W2=0-q0<d-b} is satisfied provided, for some $C$ large enough,
\begin{equation}\label{Large-tau-q0<d-Choice-1}
\left\{
\begin{array}{l}
\ds \tau^{(2 - \frac{d}{q_0}) \left(\frac{3}{4}+\frac{1}{2d}\right)} \geq C \|V \|_{L^{q_0}(\Omega)},\\
\ds
\tau^{\frac{1}{2} - \left(\frac{3}{4} - \frac{1}{2d}\right)\frac{d}{q_1}} \geq C \| W_{1} \|_{L^{q_1}(\Omega)},
\\
\tau^{\frac{1}{2} - \left(\frac{3}{4} - \frac{1}{2d}\right)\frac{d}{q_2}} \geq C \| W_{2} \|_{L^{q_2}(\Omega)},
\medskip
\\
\ds \tau^{1-\frac{1}{d} - \left(\frac{3}{4} - \frac{1}{2d}\right) \left(\frac{d}{q_1} + \frac{d}{q_2}\right) } \geq C \| W_{1} \|_{L^{q_1}(\Omega)} \| W_{2} \|_{L^{q_2}(\Omega)},
\end{array}\right.
\end{equation}
that is, with the notations of Theorem~\ref{Thm-QUCP},
\begin{multline}\label{Large-tau-q0<d-Choice-1b}
\tau \geq\\
C \left(\norm{V}_{L^{q_0}(\Omega)}^{\gamma(q_0)}+\norm{W_1}_{L^{q_1}(\Omega)}^{\delta(q_1)} + \norm{W_2}_{L^{q_2}(\Omega)}^{\delta(q_2)} + \left(\norm{W_1}_{L^{q_1}(\Omega)}\norm{W_2}_{L^{q_2}(\Omega)}\right)^{\rho(q_1,q_2)}
\right).
\end{multline}
For $q_0 \in (\frac{d}{2}, d]$, $q_1 > \frac{3d}{2} $, and $q_2 > \frac{3d}{2} $, one can alternatively choose $\alpha_0 = (\frac{3}{4} + \frac{1}{2d})\frac{d}{q_0}$, $\alpha_1 = \frac{3d}{4q_1}$, $\alpha_2 = \frac{3d}{4q_2}$, so that system~\eqref{Conditions-W1=W2=0-q0<d-b} is satisfied provided, for some $C$ large enough,
\begin{equation*}
\ds
\tau^{\left(2 - \frac{d}{q_0}\right) \left(\frac{3}{4}+\frac{1}{2d}\right)} \geq C \|V \|_{L^{q_0}(\Omega)},
\ds
\quad\tau^{\frac{1}{2} - \frac{3d}{4q_1}} \geq C \| W_{1} \|_{L^{q_1}(\Omega)},
\quad
\tau^{\frac{1}{2} - \frac{3d}{4q_2}} \geq C \| W_{2} \|_{L^{q_2}(\Omega)},
\end{equation*}
that is, with the notations of Theorem~\ref{Thm-QUCP},
\begin{equation}\label{Large-tau-q0<d-Choice-2b}
\tau \geq C \left(\norm{V}_{L^{q_0}(\Omega)}^{\gamma(q_0)}+\norm{W_1}_{L^{q_1}(\Omega)}^{\tilde \delta(q_1)} + \norm{W_2}_{L^{q_2}(\Omega)}^{\tilde \delta(q_2)}\right).
\end{equation}
Taking $\tau$ large enough that saturates condition~\eqref{Large-tau-q0<d-Choice-1b}, respectively condition~\eqref{Large-tau-q0<d-Choice-2b}, bounding the weight function $e^{\tau \varphi}$ from below and from above in~\eqref{Carleman-General-appl-final}, we easily deduce Theorem~\ref{Thm-QUCP} item 1, respectively item 2, for $q_0 \in (\frac{d}{2}, d]$.

\paragraph{The case $V \in L^{q_0}(\Omega)$ with $q_0 \in [d,\infty]$}
For $V \in L^{q_0}(\Omega)$ with $q_0 \in [d, \infty]$ and $\lambda_0>0$ to be chosen later, we set $V_{\frac{d}{2}} = 0$, $V_d = V 1_{|V|\,>\,\lambda_0}$, $V_\infty = V 1_{|V| \leq \lambda_0}$, for which we have the following estimates:
\[
\lambda_0^{q_0-d} \| V_{d}\|_{L^{d}(\Omega)}^{d} \leq \| V \|_{L^{q_0}(\Omega)}^{q_0},
\qquad
\| V_{\infty} \|_{L^\infty(\Omega)} \leq \lambda_0.
\]
The potentials $W_1 \in L^{q_1}(\Omega)$ and $W_2 \in L^{q_2}(\Omega)$ are decomposed as before $W_1 = W_{1,d} + W_{1, \infty}$, $W_2 = W_{2,d} + W_{2, \infty}$ with the estimates~\eqref{Est-W-1-d-infty} for positive parameters $\lambda_1$ and $\lambda_2$ to be chosen later. Similarly as before, conditions~\eqref{Conditions-Summary} are thus satisfied provided
\begin{equation}\label{Conditions-W1=W2=0-q0>d-a}
\left\{
\begin{array}{l}
\ds
\lambda_0^{1-\frac{q_0}{d}} \| V \|_{L^{q_0}(\Omega)}^{\frac{q_0}{d}} \leq c_0 \tau^{\frac{3}{4}+\frac{1}{2d}},
\quad
\lambda_0 \leq c_0 \tau^{\frac{3}{2}},
\\[0.15em]
\ds
\lambda_1 + \lambda_2 \leq c_0 \tau^{\frac{1}{2}}, \quad
\tau^{\frac{1}{4}-\frac{1}{2d}} \left(\lambda_1^{1- \frac{q_1}{d}} \| W_{1} \|_{L^{q_1}(\Omega)}^{ \frac{q_1}{d}} + \lambda_2^{1- \frac{q_2}{d}} \| W_{2} \|_{L^{q_2}(\Omega)}^{ \frac{q_2}{d}}\right) \leq c_0,
\\[0.15em]
\ds \tau^{\frac{1}{2}} \lambda_1^{1- \frac{q_1}{d}} \| W_{1} \|_{L^{q_1}(\Omega)}^{ \frac{q_1}{d}} \lambda_2^{1- \frac{q_2}{d}} \| W_{2} \|_{L^{q_2}(\Omega)}^{ \frac{q_2}{d}} \leq c_0.
\end{array}\right.
\end{equation}
Similarly as before, for $q_0 \in [d, \infty]$, $q_1 > \frac{3d-2}{2} $, $q_2 > \frac{3d-2}{2} $ with $\frac{1}{q_1} +\frac{1}{q_2} < 4(1-\frac{1}{d})/(3d -2)$, we choose $\alpha_0 = (\frac{3}{4} - \frac{1}{2d})\frac{d}{q_0}$, $\alpha_1 = (\frac{3}{4} - \frac{1}{2d}) \frac{d}{q_1}$, $\alpha_2 = (\frac{3}{4} - \frac{1}{2d}) \frac{d}{q_2}$, and $\lambda_0 = \tau^{\alpha_0} \| V\|_{L^{q_0}(\Omega)}$, $\lambda_1 = \tau^{\alpha_1} \| W_1 \|_{L^{q_1}(\Omega)}$ and $\lambda_2 = \tau^{\alpha_2} \| W_2 \|_{L^{q_2}(\Omega)}$. We then deduce that system~\eqref{Conditions-W1=W2=0-q0>d-a} is satisfied provided, for some $C$ large enough,
\begin{equation}\label{Large-tau-q0>d-Choice-1}
\left\{
\begin{array}{l}
\ds \tau^{\left(2 - \frac{d}{q_0}\right) \left(\frac{3}{4}-\frac{1}{2d}\right)} \geq C \|V \|_{L^{q_0}(\Omega)},
\\[0.15em]
\ds
\tau^{\frac{1}{2} - \left(\frac{3}{4} - \frac{1}{2d}\right)\frac{d}{q_1}} \geq C \| W_{1} \|_{L^{q_1}(\Omega)},
\\
\tau^{\frac{1}{2} - \left(\frac{3}{4} - \frac{1}{2d}\right)\frac{d}{q_2}} \geq C \| W_{2} \|_{L^{q_2}(\Omega)},
\\[0.15em]
\ds \tau^{1-\frac{1}{d} - \left(\frac{3}{4} - \frac{1}{2d}\right)\left(\frac{d}{q_1} + \frac{d}{q_2}\right)} \geq C \| W_{1} \|_{L^{q_1}(\Omega)} \| W_{2} \|_{L^{q_2}(\Omega)},
\end{array}\right.
\end{equation}
that is, \eqref{Large-tau-q0<d-Choice-1b} with the notations of Theorem~\ref{Thm-QUCP}.

Here again, for $q_0 \in [d,\infty]$, $q_1 > \frac{3d}{2} $, and $q_2 > \frac{3d}{2} $, one can alternatively choose $\alpha_0 = (\frac{3}{4} - \frac{1}{2d})\frac{d}{q_0}$, $\alpha_1 = \frac{3d}{4q_1}$, $\alpha_2 = \frac{3d}{4q_2}$, so that conditions~\eqref{Conditions-W1=W2=0-q0>d-a} are implied by
\begin{equation*}
\tau^{\left(\frac{3}{2} - \frac{d}{q_0}\right) \left(\frac{3}{4}+\frac{1}{2d}\right)} \geq C \|V \|_{L^{q_0}(\Omega)},
\quad \tau^{\frac{1}{2} - \frac{3d}{4q_1}} \geq C \| W_{1} \|_{L^{q_1}(\Omega)},
\quad
\tau^{\frac{1}{2} - \frac{3d}{4q_2}} \geq C \| W_{2} \|_{L^{q_2}(\Omega)},\label{Large-tau-q0>d-Choice-2}
\end{equation*}
that is~\eqref{Large-tau-q0<d-Choice-2b} with the notations of Theorem~\ref{Thm-QUCP}.

We then deduce Theorem~\ref{Thm-QUCP} in the case $q\geq d$ immediately from~\eqref{Carleman-General-appl-final} as in the case $q \leq d$.

\begin{rema}
In fact, if we focus on the conditions
\begin{equation}\label{Conditions-W1W2}
\left\{
\begin{array}{l}
\ds
\lambda_1 + \lambda_2 \leq c_0 \tau^{\frac{1}{2}}, \quad
\tau^{\frac{1}{4}-\frac{1}{2d}} \left(\lambda_1^{1- \frac{q_1}{d}} \| W_{1} \|_{L^{q_1}(\Omega)}^{ \frac{q_1}{d}} + \lambda_2^{1- \frac{q_2}{d}} \| W_{2} \|_{L^{q_2}(\Omega)}^{ \frac{q_2}{d}}\right) \leq c_0,
\\[0.15em]
\ds \tau^{\frac{1}{2}} \lambda_1^{1- \frac{q_1}{d}} \| W_{1} \|_{L^{q_1}(\Omega)}^{ \frac{q_1}{d}} \lambda_2^{1- \frac{q_2}{d}} \| W_{2} \|_{L^{q_2}(\Omega)}^{ \frac{q_2}{d}} \leq c_0,
\end{array}\right.
\end{equation}
which appear in the second and third lines of system~\eqref{Conditions-W1=W2=0-q0<d-a} and~\eqref{Conditions-W1=W2=0-q0>d-a}, and choose $\lambda_1 = \tau^{\alpha_1} \| W_1 \|_{L^{q_1}(\Omega)}$ and $\lambda_2 = \tau^{\alpha_2} \| W_2 \|_{L^{q_2}(\Omega)}$, one can find $\tau$ large enough so that system~\eqref{Conditions-W1W2} is satisfied provided
\begin{equation}\label{Conditions-alpha-1-2}
\left\{
\begin{array}{l}
\ds  \alpha_1 < \frac{1}{2},
\quad
\alpha_2 < \frac{1}{2},
\smallskip
\\
\ds  \alpha_1\! \left(\frac{q_1}{d}-1 \right)  > \frac{1}{4} - \frac{1}{2d},
\quad \alpha_2\! \left(\frac{q_2}{d}-1 \right)  >  \frac{1}{4} - \frac{1}{2d},
\\[0.25em]
\mkern150mu
\ds\alpha_1 \!\left(\frac{q_1}{d}-1 \right) + \alpha_2\! \left(\frac{q_2}{d}-1 \right)  >  \frac{1}{2}.
\end{array}\right.
\end{equation}
Indeed, in this case, it suffices to take, for a sufficiently large constant $C$,
\begin{equation}\label{Choice-tau}
\tau \geq C
\left(
\|W_1\|_{L^{q_1}(\Omega)}^{a_{q_1}(\alpha_1)} +
\|W_2\|_{L^{q_2}(\Omega)}^{a_{q_2} (\alpha_2)} +
\left(\|W_1\|_{L^{q_1}(\Omega)} \|W_2\|_{L^{q_2}(\Omega)}\right)^{b(\alpha_1, \alpha_2)}
\right),
\end{equation}
with
\begin{align*}
a_q(\alpha) &= \max\left\{\frac{1}{\ds \frac{1}{2} - \alpha}, \frac{1}{\ds \alpha \left(\frac{q}{d} -1 \right)- \frac{1}{4} + \frac{1}{2d} } \right\},
\\
b_{q_1, q_2} (\alpha_1, \alpha_2) &= \frac{1}{\ds \alpha_1 \left(\frac{q_1}{d} -1 \right)+ \alpha_2 \left(\frac{q_2}{d} -1 \right)- \frac{1}{2}}.
\end{align*}
Although it is rather easy to check that the system~\eqref{Conditions-alpha-1-2} admits solutions $(\alpha_1, \alpha_2)$ if $q_1$ and $q_2$ satisfy $q_1 > \frac{3d}{2} - 1$, $q_2 > \frac{3d}{2} - 1$ and $q_1 + q_2 > 3d$, it is not clear how to choose $\alpha_1$ and $\alpha_2$ satisfying~\eqref{Conditions-alpha-1-2} to minimize $\tau$ in~\eqref{Choice-tau}. We have thus decided in the above proof of Theorem~\ref{Thm-QUCP} to restrict ourselves to the case in which both terms in $a_q(\alpha)$ are equal (this choice corresponds to item 1 in Theorem~\ref{Thm-QUCP}), or to consider, instead of~\eqref{Conditions-W1W2}, the sufficient conditions
\begin{equation*}
\ds
\lambda_1 + \lambda_2 \leq c_0 \tau^{\frac{1}{2}}, \quad
\tau^{\frac{1}{4}} \left(\lambda_1^{1- \frac{q_1}{d}} \| W_{1} \|_{L^{q_1}(\Omega)}^{ \frac{q_1}{d}} + \lambda_2^{1- \frac{q_2}{d}} \| W_{2} \|_{L^{q_2}(\Omega)}^{ \frac{q_2}{d}}\right) \leq c_0,
\end{equation*}
this choice yielding to item~\eqref{theo1.3.2} in Theorem~\ref{Thm-QUCP}.
\end{rema}



\appendix
\section{Reminder of some classical results in harmonic analysis}

We start by recalling the classical Hardy--Littlewood--Sobolev theorem.

\begin{theo}[{\cite[Theorem~4.5.3.]{hormanderanalysis1}} Hardy--Littlewood--Sobolev theorem]\label{Thm-Hardy-Littlewood-Sobolev}
Let $n \in \N$. For $(p,q,r) \in (1, \infty)^3$ such that
\[
\frac{1}{r}=1-\left(\frac{1}{p}-\frac{1}{q}\right),
\]
there exists a constant $C_{p,q,n}$ such that for all $f \in L^p(\R^n)$,
\[
\left\| x\mapsto \int_{\R^n}|x-y|^{- \frac{n}{r} } f(y)\, d y
\right\|_{L^q(\R^n)}
\leq C_{p, q,n}
\|f\|_{L^p(\R^n)}.
\]
\end{theo}

In the article, we have also used the stationary phase lemma. Although it is a very classical lemma of harmonic analysis, we used the following version proved in~\cite{alazard2017stationary}, which presents the advantage of quantifying precisely the constants in the stationary phase lemma:

\begin{theo}[{\cite[Theorem~1]{alazard2017stationary} Stationary phase lemma}]\label{Thm-Stationary-Phase}
Let $\Phi \in C^{\infty}(\R^d; \R)$, $b \in C_0^{\infty}(\R^d; \R)$. For $\lambda \in \R$, we introduce
\[
I_{\Phi,b}(\lambda)=\int_{\bfR^d} e^{i \lambda \Phi(\xi)} b(\xi) d \xi.
\]
Set $K=\supp b$ and let $V$ be an open neighbourhood of $K$, and use the following notations:
\begin{itemize}
\item $\ds\clM_{d+2}:=\sum_{2\, \leq\,|\alpha| \,\leq\,d+2}\; \sup _{\xi\,\in\,V}\left|D_{\xi}^\alpha \Phi(\xi)\right|$,
\item $\ds \clN_{d+1}:=\sum_{|\alpha|\,\leq \,d+1}\; \sup _{\xi\,\in\,K}\left|D_{\xi}^\alpha b(\xi)\right|$,
\end{itemize}
and assume that there exists $a_0 >0$ such that for all $\xi \in V$, $|\det (\Hess \Phi(\xi)) | \geq a_0$, where $ \Hess \Phi(\xi) $ denotes the Hessian matrix of $\Phi$ at $\xi$.

There exists a constant $C$ independent of $(\Phi, b)$ satisfying the above assumptions, such that for all $\lambda \geq 1$,
\[
|I_{\Phi,b} (\lambda)|
\leq
\frac{C}{a_0^{1+d}}\left(1+\clM_{d+2}^{\frac{d}{2}+d^2}\right) \clN_{d+1} \lambda^{-\frac{d}{2}}.
\]
\end{theo}

\section{Proof of Lemma~\ref{Lem-Boundary-H-1/2}}\label{Appendix-B}

\begin{proof}
The proof of Lemma~\ref{Lem-Boundary-H-1/2} relies on a suitable interpolation estimate.

First, for $\tau \geq \tau_0$, we define the operator $ \Lambda_{\tau}: L^2(\partial \Omega) \to L^2(\Omega\setminus\omega; L^2(\partial\Omega)) $ by
\begin{equation*}
\Lambda_{\tau} g(x_0) = \eta_{x_0, \tau}(\cdot) g (\cdot), \quad \text{ for } x_0 \in \Omega \setminus \omega,
\end{equation*}
where we recall that $\eta_{x_0, \tau}$ is the function given by $\eta_{x_0, \tau} (x) = \eta (\tau^{\frac{1}{3}} (x-x_0))$ for $x \in \R^d$, for a smooth compactly supported function $\eta$.

Using~\eqref{Est-rho-0-above}--\eqref{Est-rho-2-above}, it is easy to check that there exists a constant $C>0$ such that
\begin{align*}
&\forall g \in L^2(\partial \Omega),
\quad
\| \Lambda_\tau g \|_{L^2\left(\Omega\setminus\omega; L^2(\partial\Omega)\right)}^2
\leq C \tau^{-\frac{d}{3}} \|g \|_{L^2(\partial \Omega)}^2,
\\
&\forall g \in H^1(\partial \Omega),
\quad
\| \Lambda_\tau g \|_{L^2\left(\Omega\setminus\omega; H^1(\partial\Omega)\right)}^2
\leq C \tau^{-\frac{d}{3}} \left(\|g \|_{H^1(\partial \Omega)}^2 + \tau^{\frac{2}{3}} \| g \|_{L^2(\partial \Omega)}^2\right).
\end{align*}

We can then deduce easily the estimate~\eqref{Boundary-H-1/2-norms}. More precisely, by interpolation, $\Lambda_\tau$ maps $H^{\frac{1}{2}}(\partial \Omega)$ to $ L^2(\Omega\setminus\omega; H^{\frac{1}{2}}(\partial\Omega))$. To estimate the operator in this norm with appropriate powers of $\tau$, we proceed as follows. We let $(\Phi_j)_{j\, \in\,\N} $ be the basis of eigenfunctions of the Laplace Beltrami operator $-\Delta$ on $\partial \Omega$, with corresponding eigenvalues $(\lambda_j^2)_{j\,\in\,\N}$, which are non-negative and going to infinity. Accordingly, for $g = \sum_j a_j \Phi_j$, the $L^2(\partial \Omega)$, $H^1(\partial \Omega)$ and $H^{\frac{1}{2}}(\partial \Omega)$ norms of $g$ can be read as, respectively, $\| (a_j)\|_{\ell^2(\N)}$, $\|(a_j (\lambda_j +1)) \|_{\ell^2(\N)}$, and $\|(a_j (\lambda_j +1)^{\frac{1}{2}})\|_{\ell^2(\N)}$.

Writing $g \in H^{\frac{1}{2}}(\partial \Omega)$ under the form $g = \sum_{j\,\in\,\N} a_j \Phi_j$, we then introduce the function
\[
f(z) = \sum_{j\,\in\,\N} a_j (\lambda_j +1 + \tau^{\frac{1}{3}})^{\frac{1}{2} - z} \Lambda_\tau \Phi_j,
\qquad \qquad z \in \C \text{ with } \Re(z) \in [0,1].
\]
The function $f$ is holomorphic in $\{z \in \C \text{ with } \Re(z) \in (0,1)\}$ with values in $L^2(\Omega \setminus\omega; L^2(\partial \Omega))$, $f(\frac{1}{2}) = \Lambda_\tau g$, and
\begin{align*}
& \forall \beta \in \R,\quad
\ds \| f(i \beta) \|_{L^2\left(\Omega \setminus\omega; L^2(\partial \Omega)\right)}^2 &&\leq C \tau^{-\frac{d}{3}} \sum_{j} |a_j|^2 \left(\lambda_j + 1+ \tau^{\frac{1}{3}}\right),
\\
& \forall \beta \in \R,\quad
\ds\| f(1+ i \beta) \|_{L^2\left(\Omega \setminus\omega; H^1(\partial \Omega)\right)}^2
&&\leq C \tau^{-\frac{d}{3}}
\sum_{j} |a_j|^2 \left(\lambda_j + 1+ \tau^{\frac{1}{3}}\right).
\end{align*}
Since $L^2(\Omega \setminus\omega; H^{\frac{1}{2}}(\partial \Omega)) = [L^2(\Omega \setminus\omega; L^2(\partial \Omega)), L^2(\Omega \setminus\omega; H^1(\partial \Omega))]_{\frac{1}{2}}$, we deduce from the above estimates that there exists $C>0$ such that for all $g \in H^{\frac{1}{2}}(\partial \Omega)$ and $\tau \geq \tau_0$,
\[
\left\| \Lambda_\tau g\right\|_{L^2\left(\Omega \setminus\omega; H^{\frac{1}{2}}(\partial \Omega)\right) }^2
\leq C \tau^{-\frac{d}{3}} \left(\| g \|_{H^{\frac{1}{2}}(\partial \Omega)}^2 + \tau^{\frac{1}{3}} \| g \|_{L^2(\partial \Omega)}^2\right).
\]
This concludes the proof of Lemma~\ref{Lem-Boundary-H-1/2}.
\end{proof}

\vspace*{45pt}


%\newcommand{\etalchar}[1]{$^{#1}$}
%\providecommand{\bysame}{\leavevmode\text to3em{\hrulefill}\thinspace}
%\providecommand{\MR}{\relax\ifhmode\unskip\space\fi MR }
%\providecommand{\MRhref}[2]{\href{http://www.ams.org/mathscinet-getitem?mr=#1}{#2} }
%\providecommand{\href}[2]{#2}


\begin{thebibliography}{99}
\bibitem{Amrein-Berthier-Georgescu-1981} Amrein, Werner O.; Berthier, Anne-Marie; Georgescu, Vladimir ${L}^ p-$inequalities for the Laplacian and unique continuation, Ann. Inst. Fourier, Volume 31 (1981) no. 3, pp. 153-168 | DOI | Zbl | MR
\bibitem{alazard2017stationary} Alazard, Thomas; Burq, Nicolas; Zuily, Claude A stationary phase type estimate, Proc. Am. Math. Soc., Volume 145 (2017) no. 7, pp. 2871-2880 | DOI | Zbl | MR
\bibitem{Boyer-Fabrie-Book} Boyer, Franck; Fabrie, Pierre Mathematical tools for the study of the incompressible Navier–Stokes equations and related models, Applied Mathematical Sciences, 183, Springer, 2012 | DOI | Zbl | MR
\bibitem{barcelo1988weighted} Barcelo, Bartolomé; Kenig, Carlos E.; Ruiz, Alberto; Sogge, Christopher D. Weighted Sobolev inequalities and unique continuation for the Laplacian plus lower order terms, Ill. J. Math., Volume 32 (1988) no. 2, pp. 230-245 | Zbl | MR
\bibitem{davey2020quantitative} Davey, Blair Quantitative unique continuation for Schrödinger operators, J. Funct. Anal., Volume 279 (2020) no. 4, 108566, 22 pages | DOI | Zbl | MR
\bibitem{darde2023short} Dardé, Jérémi; Ervedoza, Sylvain A short introduction to Carleman estimates, Control of Partial Differential Equations (Series in Contemporary Applied Mathematics), Volume 24, World Scientific, 2023, pp. 1-79 | DOI
\bibitem{ferreira2003sharp} Dos Santos Ferreira, David Sharp ${L}^p$ Carleman estimates and unique continuation, Duke Math. J., Volume 129 (2005) no. 3, pp. 503-550 | DOI | Zbl | MR
\bibitem{davey2019quantitative} Davey, Blair; Zhu, Jiuyi Quantitative uniqueness of solutions to second-order elliptic equations with singular lower order terms, Commun. Partial Differ. Equations, Volume 44 (2019) no. 11, pp. 1217-1251 | DOI | Zbl | MR
\bibitem{Duyckaerts2008} Duyckaerts, Thomas; Zhang, Xu; Zuazua, Enrique On the optimality of the observability inequalities for parabolic and hyperbolic systems with potentials, Ann. Inst. Henri Poincaré, Anal. Non Linéaire, Volume 25 (2008) no. 1, pp. 1-41 | DOI | Zbl | Numdam | MR
\bibitem{Fursikov1996-book} Fursikov, Andreĭ V.; Imanuvilov, Oleg Yu. Controllability of evolution equations, Lecture Notes Series, Seoul, 34, Seoul National University, 1996 | Zbl | MR
\bibitem{Hormander-Uniqueness-1983} Hörmander, Lars Uniqueness theorems for second order elliptic differential equations, Commun. Partial Differ. Equations, Volume 8 (1983) no. 1, pp. 21-64 | DOI | Zbl | MR
\bibitem{hormanderanalysis1} Hörmander, Lars The analysis of linear partial differential operators. I, Grundlehren der Mathematischen Wissenschaften, 256, Springer, 1990 | DOI | Zbl | MR
\bibitem{hormanderanalysis4} Hörmander, Lars The analysis of linear partial differential operators. IV, Classics in Mathematics, Springer, 1994 | DOI | Zbl | MR
\bibitem{imanuvilov2003global} Imanuvilov, Oleg Yu.; Puel, Jean-Pierre Global Carleman estimates for weak solutions of elliptic nonhomogeneous dirichlet problems, Int. Math. Res. Not., Volume 2003 (2003) no. 16, pp. 883-913 | MR | DOI | Zbl
\bibitem{jerison1986carleman} Jerison, David Carleman inequalities for the Dirac and Laplace operators and unique continuation, Adv. Math., Volume 62 (1986) no. 2, pp. 118-134 | DOI | Zbl
\bibitem{jerison1985unique} Jerison, David; Kenig, Carlos E. Unique continuation and absence of positive eigenvalues for Schrödinger operators, Ann. Math., Volume 121 (1985) no. 3, pp. 463-488 | DOI | Zbl
\bibitem{kenig1987uniform} Kenig, Carlos E.; Ruiz, Alberto; Sogge, Christopher D. Uniform Sobolev inequalities and unique continuation for second order constant coefficient differential operators, Duke Math. J., Volume 55 (1987) no. 2, pp. 329-347 | DOI | Zbl | MR
\bibitem{koch2001carleman} Koch, Herbert; Tataru, Daniel Carleman estimates and unique continuation for second-order elliptic equations with nonsmooth coefficients, Commun. Pure Appl. Math., Volume 54 (2001) no. 3, pp. 339-360 | Zbl | DOI | MR
\bibitem{koch2002sharp} Koch, Herbert; Tataru, Daniel Sharp counterexamples in unique continuation for second order elliptic equations, J. Reine Angew. Math., Volume 542 (2002), pp. 133-146 | DOI | Zbl | MR
\bibitem{koch2005dispersive} Koch, Herbert; Tataru, Daniel Dispersive estimates for principally normal pseudodifferential operators, Pure Appl. Math., Volume 58 (2005) no. 2, pp. 217-284 | DOI | Zbl | MR
\bibitem{LeRousseau2013interface} Le Rousseau, Jérôme; Lerner, Nicolas Carleman estimates for anisotropic elliptic operators with jumps at an interface, Anal. PDE, Volume 6 (2013) no. 7, pp. 1601-1648 | DOI | Zbl | MR
\bibitem{LeRousseaubookvolume1} Le Rousseau, Jérôme; Lebeau, Gilles; Robbiano, Luc Elliptic Carleman estimates and applications to stabilization and controllability. Vol. I. Dirichlet boundary conditions on Euclidean space, Progress in Nonlinear Differential Equations and their Applications, 97, Birkhäuser, 2022 | DOI | Zbl | MR
\bibitem{malinnikova2012quantitative} Malinnikova, Eugenia; Vessella, Sergio Quantitative uniqueness for elliptic equations with singular lower order terms, Math. Ann., Volume 353 (2012) no. 4, pp. 1157-1181 | DOI | Zbl | MR
\bibitem{sogge1988concerning} Sogge, Christopher D. Concerning the ${L}^p$ norm of spectral clusters for second-order elliptic operators on compact manifolds, J. Funct. Anal., Volume 77 (1988) no. 1, pp. 123-138 | DOI | Zbl | MR
\bibitem{Sogge1989uniquecontinuation} Sogge, Christopher D. Oscillatory integrals and unique continuation for second order elliptic differential equations, J. Am. Math. Soc., Volume 2 (1989) no. 3, pp. 491-515 | DOI | Zbl | MR
\bibitem{Sogge.Strong.Elliptic1990} Sogge, Christopher D. Strong uniqueness theorems for second order elliptic differential equations, Am. J. Math., Volume 112 (1990) no. 6, pp. 943-984 | DOI | Zbl | MR
\bibitem{sogge2017fourier} Sogge, Christopher D. Fourier integrals in classical analysis, Cambridge Tracts in Mathematics, 210, Cambridge University Press, 2017 | DOI | Zbl | MR
\bibitem{stein2016singular} Stein, Elias M. Singular integrals and differentiability properties of functions, Princeton Mathematical Series, 30, Princeton University Press, 1970 | Zbl | MR
\bibitem{stein1993harmonic} Stein, Elias M. Harmonic analysis: real-variable methods, orthogonality, and oscillatory integrals, Princeton Mathematical Series, 3, Princeton University Press, 1993 | Zbl | MR
\bibitem{tomas1975restriction} Tomas, Peter A. A restriction theorem for the Fourier transform, Bull. Am. Math. Soc., Volume 81 (1975) no. 1, pp. 477-478 | DOI | Zbl | MR
\bibitem{wolff1992property} Wolff, Thomas H. A property of measures in $\mathbb{R}^{N}$ and an application to unique continuation, Geom. Funct. Anal., Volume 2 (1992) no. 2, pp. 225-284 | DOI | Zbl | MR
\bibitem{Wolff-1993} Wolff, Thomas H. Recent work on sharp estimates in second-order elliptic unique continuation problems, J. Geom. Anal., Volume 3 (1993) no. 6, pp. 621-650 | DOI | Zbl | MR
\end{thebibliography}
\end{document}
